\chapter{Intrinsic Correlative Operations of the Same Continuum}
\label{chap:pdfv2-operaciones-intrinsecas-correlativas}

\begin{thesisbox}
The symbols \(\mathrm{sp},\mathrm{sol},\mathrm{gau},\mathrm{coh}\) and
\(\mathrm{det}\) that appear in this chapter mark five correlative operations
over a single APP--TRIT--TPK state\@. Each operation receives
the same surviving extension, the
same orientation, the same carry, the same memory, the same boundary, and the
same route. The following sections are successive mathematical stages of a
single construction and not autonomous expositions.
\end{thesisbox}

\section{Generating edge and common genealogical state}

We do not repeat here the state, tree, or measure constructed in the preceding
chapter. We take
\(\widehat x=\widehat x_k(\gamma)\in
\widehat{\mathcal X}^{\rm tot}_k\) from
\eqref{eq:continuo-comun-estado} and all its admissible emissions
\(\alpha\in\operatorname{Out}(\widehat x)\). If
\(\widehat y=\alpha_*\widehat x\), we write
\begin{equation}
\label{eq:pdfv2-cor-hijos-supervivientes}
 \operatorname{Ch}_k(\widehat x):=
 \{(\alpha,\alpha_*\widehat x):
      \alpha\in\operatorname{Out}(\widehat x)\},
 \qquad
 d_k(\widehat x):=\#\operatorname{Ch}_k(\widehat x).
\end{equation}
The finiteness and non-emptiness of this set are already
\eqref{eq:continuo-comun-out-finito}; survival at all horizons,
the complete multisection, and its measure are
\cref{thm:continuo-comun-total-superviviente,thm:continuo-comun-multiseccion-no-vacia,thm:continuo-comun-medida-canonica}. A second non-emptiness is not again assumed or proved.

Each pair \((\alpha,\widehat y)\) of
\eqref{eq:pdfv2-cor-hijos-supervivientes} preserves the unique record
\begin{equation}
\label{eq:pdfv2-cor-registro-extension}
 \begin{aligned}
 \mathfrak e_\alpha:=\bigl(&
 \alpha,\operatorname{Ext}_{\gamma_\alpha};
 \operatorname{Sh}_{s(\alpha)},\operatorname{Sh}_{t(\alpha)};\\
 &
 \operatorname{res}^{\Sigma,\Pi}(\alpha),
 \operatorname{quo}^{\Sigma,\Pi}(\alpha);\\
 &\operatorname{or}(\alpha),\kappa(\alpha),\mathsf C_\alpha;
 \mathsf Q(s(\alpha)),\mathsf Q(t(\alpha)),g(s(\alpha)),g(t(\alpha)),
 q^{(9)}(s(\alpha)),q^{(9)}(t(\alpha));\\
 &
 m(\gamma\alpha;m_0),\partial\alpha,\gamma\alpha,
 \operatorname{Led}(\gamma\alpha);\\
 &\{\operatorname{Term}_{k+1,N}(\widehat y)\}_{N\geq k+1},
 \operatorname{Msect}(\widehat y),\mu_{\widehat y}\bigr).
 \end{aligned}
\end{equation}
Here \(\operatorname{Ext}_{\gamma_\alpha}\) denotes the complete relation that
the admissible emission already satisfies; it introduces no additional filter. The subsequent
operations are coordinated applications of \(\mathfrak e_\alpha\) and do not
reconstruct the record from their outputs.

\section{Intrinsic Correlative Actions on the Common Generator}

The internal marks are fixed without changing their source: \(\mathrm{sp}\) denotes transport and grading of the two sheets; \(\mathrm{sol}\), signed increment, balance, and memory; \(\mathrm{gau}\), incidence and carry curvature; \(\mathrm{coh}\), ports, rectangular difference, and cellular complex; and \(\mathrm{det}\), based complex, correlative filling, and determinant line. The following formulas display their couplings over the same \(\mathfrak e_\alpha\); their actions literally share source, target, route, and ledger.

The double APP evaluation provides, on the same support, the labels
\(\Sigma\) and \(\Pi\). For an extension \(\alpha\), their multiplicities
are calculated from the complete ledger, without an external counter:
\begin{equation}
\label{eq:pdfv2-cor-multiplicidades-hojas}
 N_\diamond(\alpha):=
 \sum_{\epsilon\ {\rm en}\ \gamma\alpha}
 \mathbf1_{\{\operatorname{res}^{\diamond}(\epsilon)\neq0\}},
 \qquad \diamond\in\{\Sigma,\Pi\}.
\end{equation}
Let there also be the total oriented coefficient
\[
 \Delta_{\gamma_\alpha}:=
 \sum_{\epsilon\ {\rm en}\ \gamma\alpha}\operatorname{or}(\epsilon)
 \in\mathbb Z.
\]
The signed action is
\begin{equation}
\label{eq:pdfv2-cor-b-alpha}
 b_\alpha:=\Delta_{\gamma_\alpha}
 \bigl(N_\Sigma(\alpha)-N_\Pi(\alpha)\bigr),
 \qquad
 b_\alpha^\pm:=\frac{|b_\alpha|\pm b_\alpha}{2}.
\end{equation}
By construction,
\begin{equation}
\label{eq:pdfv2-cor-b-identidades}
 b_\alpha=b_\alpha^+-b_\alpha^-,\qquad
 |b_\alpha|=b_\alpha^++b_\alpha^-,\qquad
 b_\alpha^+b_\alpha^-=0.
\end{equation}

The same set of extensions simultaneously produces the ports
\begin{equation}
\label{eq:pdfv2-cor-puertos-comunes}
 I_k(\widehat x_k):=
 \operatorname{Ch}_k(\widehat x_k)\times\{\Sigma\},
 \qquad
 J_k(\widehat x_k):=
 \operatorname{Ch}_k(\widehat x_k)\times\{\Pi\}.
\end{equation}
A visible leaf is not assumed to contain more prolongations than the other:
each surviving extension is evaluated on both leaves of the same support.
By \eqref{eq:continuo-comun-out-finito}, both sets of ports are
nonempty.

For \(A_N=\mathbb Z/3^N\mathbb Z\), the cellular action on generators is
\begin{equation}
\label{eq:pdfv2-cor-kappa-celda}
 \begin{aligned}
 \kappa_{k,N}:A_N^{I_k}\oplus A_N^{J_k}
 &\longrightarrow A_N^{I_k\times J_k},\\
 \kappa_{k,N}(u,v)_{(\alpha,\Sigma),(\beta,\Pi)}
 &:=u_{(\alpha,\Sigma)}-v_{(\beta,\Pi)}.
 \end{aligned}
\end{equation}
Selected first ports according to the TPK order, the rectangular difference
\begin{equation}
\label{eq:pdfv2-cor-delta-celda}
 \delta(w)_{ij}:=w_{ij}-w_{ij_0}-w_{i_0j}+w_{i_0j_0}
\end{equation}
download the split sequence
\begin{equation}
\label{eq:pdfv2-cor-sucesion-celda}
 0\longrightarrow A_N\xrightarrow{\Delta}
 A_N^{I_k}\oplus A_N^{J_k}
 \xrightarrow{\kappa_{k,N}}A_N^{I_k\times J_k}
 \xrightarrow{\delta}A_N^{(|I_k|-1)(|J_k|-1)}
 \longrightarrow0.
\end{equation}
This is an identity of the same APP fibre: it adds no cell foreign to the
generator \(\mathfrak e_\alpha\).

The local cellular action uses the two integer coordinates accumulated by that
same ledger. For
\(\widehat x=\widehat x_k(\gamma)\) and
\(\widehat y=\alpha_*\widehat x=\widehat x_{k+1}(\gamma\alpha)\), we set
\begin{equation}
\label{eq:pdfv2-cor-c-v}
 \begin{aligned}
 c_{\widehat x}&:=\kappa(\gamma), &
 c_\alpha:=c_{\widehat y}&:=\kappa(\gamma\alpha),\\
 v_{\widehat x}&:=\sum_{\epsilon\ {\rm en}\ \gamma}
                  \operatorname{or}(\epsilon), &
 v_{\widehat y}&:=\sum_{\epsilon\ {\rm en}\ \gamma\alpha}
                  \operatorname{or}(\epsilon),\\
 v_\alpha&:=v_{\widehat y}-v_{\widehat x}
           =\operatorname{or}(\alpha).&&
 \end{aligned}
\end{equation}
On the generators \(f,r,e,b,g\) is defined
\begin{equation}
\label{eq:pdfv2-cor-complejo-local}
 \partial f=3c_\alpha e+b,\qquad
 \partial r=0,\qquad
 \partial e=g,\qquad
 \partial b=-3c_\alpha g,\qquad
 \partial g=0.
\end{equation}
The cancellation
\begin{equation}
\label{eq:pdfv2-cor-dos-cero}
 \partial^2f=3c_\alpha g-3c_\alpha g=0
\end{equation}
is literal and uses the same carry with the two signs imposed by the
orientation.

The incidence seed is obtained from the free module of the two leaves,
\begin{equation}
\label{eq:pdfv2-cor-v-tpk}
 V_{\rm TPK}:=\mathbf F_3e_\Sigma\oplus\mathbf F_3e_\Pi,
 \qquad
 \mathscr I:=\mathbf P^1(V_{\rm TPK}),
 \qquad |\mathscr I|=4,
\end{equation}
and of its unordered pairs
\begin{equation}
\label{eq:pdfv2-cor-incidencias}
 \mathscr E:=\binom{\mathscr I}{2},
 \qquad |\mathscr E|=6,
 \qquad
 \mathscr B:=\mathscr I\times\{+,-\},
 \qquad |\mathscr B|=8.
\end{equation}
The incidence application is determined on the base by
\begin{equation}
\label{eq:pdfv2-cor-b-incidencia-generadores}
 B[e_L]=\sum_{L'\neq L}e_{\{L,L'\}},
 \qquad L\in\mathscr I,
\end{equation}
and, by linearity,
\begin{equation}
\label{eq:pdfv2-cor-b-incidencia-coordenadas}
 (Ba)_{\{L,L'\}}=a_L+a_{L'}.
\end{equation}
If \(s(q)=\sum_{E\in\mathscr E}q_E\), each \(L\) appears three times and
\begin{equation}
\label{eq:pdfv2-cor-sb-cero}
 s(Ba)=0.
\end{equation}
Therefore the function
\begin{equation}
\label{eq:pdfv2-cor-wc}
 W_c(q):=\mathbf1_{\{s(q)=0\}}-\frac13
\end{equation}
satisfies \(W_c(q+Ba)=W_c(q)\). The identities
\eqref{eq:pdfv2-cor-v-tpk}--\eqref{eq:pdfv2-cor-wc} are internal and exact.
The material lift of \(q\) from the cells and common carries is treated in
the assembly layer; this incidence identity is not conflated with the
construction of that lift.

\section{Unitary transport, polarization, and balance on the common fiber}

The non-emptiness of \eqref{eq:pdfv2-cor-hijos-supervivientes} allows one to fix,
without selecting a prolongation,
\begin{equation}
\label{eq:pdfv2-cor-peso-conteo}
 p_k(\alpha,\widehat y\mid\widehat x)
 :=\frac1{d_k(\widehat x)}.
\end{equation}
Let the space of the two leaves of the same state be
\begin{equation}
\label{eq:pdfv2-cor-hilbert-hojas}
 \mathscr H_k:=
 \bigoplus_{\widehat x\in\widehat{\mathcal X}^{\rm tot}_k}
 \bigl(\mathsf H_{\widehat x}^{\Sigma}
       \oplus\mathsf H_{\widehat x}^{\Pi}\bigr),
\end{equation}
and let \(\iota_{\widehat x}\) be the inclusion of the indicated fiber. The TPK
transport of \eqref{eq:continuo-comun-transporte-hojas} gives the exact action
\begin{equation}
\label{eq:pdfv2-cor-u-comun}
 U_k\iota_{\widehat x}v:=
 \sum_{(\alpha,\widehat y)\in\operatorname{Ch}_k(\widehat x)}
 p_k(\alpha,\widehat y\mid\widehat x)^{1/2}
 \iota_{\widehat y}U_\alpha v.
\end{equation}
The children of two distinct parents are disjoint because truncation is a
function. Moreover, two distinct witnesses \(\alpha\) do not collapse into the same
direct summand: the complete route \(\gamma\), including its last step, is a
coordinate of \(\widehat y\). Consequently, the Gram identity is
\begin{equation}
\label{eq:pdfv2-cor-u-gram}
 \left.U_k^*U_k\right|_{\mathsf H_{\widehat x}^{\Sigma}
       \oplus\mathsf H_{\widehat x}^{\Pi}}
 =
 \sum_{(\alpha,\widehat y)\in\operatorname{Ch}_k(\widehat x)}
 p_k(\alpha,\widehat y\mid\widehat x)\,U_\alpha^*U_\alpha.
\end{equation}
Therefore \(U_k^*U_k=I\) exactly when, in each initial fibre,
\begin{equation}
\label{eq:pdfv2-cor-gram-promedio-normalizado}
 \sum_{(\alpha,\widehat y)\in\operatorname{Ch}_k(\widehat x)}
 p_k(\alpha,\widehat y\mid\widehat x)\,U_\alpha^*U_\alpha=I.
\end{equation}
The isometry of every sheet transport \(U_\alpha\) is sufficient for this equality, together with normalization of the weights; the necessary and sufficient condition is normalization of the average Gram matrix displayed in \eqref{eq:pdfv2-cor-gram-promedio-normalizado}. For \(\ell>k\), the complete history is preserved through
\begin{equation}
\label{eq:pdfv2-cor-u-iterado}
 U_{\ell,k}:=U_{\ell-1}\cdots U_k,
 \qquad U_{k,k}:=I.
\end{equation}

The grading acts on the two coordinates already present in each fiber of \eqref{eq:pdfv2-cor-hilbert-hojas}:
\begin{equation}
\label{eq:pdfv2-cor-sigma}
 \sigma_k\iota_{\widehat x}(v_\Sigma,v_\Pi)
 :=\iota_{\widehat x}(v_\Sigma,-v_\Pi),
\qquad
 P_k:=\frac{I+\sigma_k}{2},\quad
 Q_k:=\frac{I-\sigma_k}{2}.
\end{equation}
The TRIT threshold belongs to the orientation and to the carry of
\(\widehat x\); it erases neither of the two APP sheets and does not require a third
case in \eqref{eq:pdfv2-cor-sigma}.

The leaf-preserving part and the leaf-changing part are two components of
the same \(U_k\):
\begin{align}
 A_k&:=P_{k+1}U_kP_k+Q_{k+1}U_kQ_k,
 \label{eq:pdfv2-cor-a-comun}\\
 \eta_k&:=Q_{k+1}U_kP_k+P_{k+1}U_kQ_k.
 \label{eq:pdfv2-cor-eta-comun}
\end{align}
By orthogonality of \(P_{k+1}\) and \(Q_{k+1}\),
\begin{equation}
\label{eq:pdfv2-cor-balance-hojas}
 \begin{aligned}
 U_k&=A_k+\eta_k,\\
 A_k^*A_k+\eta_k^*\eta_k
 &=P_kU_k^*U_kP_k+Q_kU_k^*U_kQ_k.
 \end{aligned}
\end{equation}
\begin{proof}[Proof of the sheetwise Gram balance]
Expanding the products, orthogonality of the target projectors gives
\[
\begin{aligned}
 A_k^*A_k
 &=P_kU_k^*P_{k+1}U_kP_k+Q_kU_k^*Q_{k+1}U_kQ_k,\\
 \eta_k^*\eta_k
 &=P_kU_k^*Q_{k+1}U_kP_k+Q_kU_k^*P_{k+1}U_kQ_k.
\end{aligned}
\]
Their sum and \(P_{k+1}+Q_{k+1}=I\) prove \eqref{eq:pdfv2-cor-balance-hojas}. The difference between \(U_k^*U_k\) and that balance is \(P_kU_k^*U_kQ_k+Q_kU_k^*U_kP_k\); it vanishes exactly when both cross blocks vanish. In particular, condition \eqref{eq:pdfv2-cor-gram-promedio-normalizado} yields \(U_k^*U_k=I\) and reduces the balance to \(I\).
\end{proof}
Proposition~\ref{x_reciprocidad:iv:cont:gram} develops this same general identity in the subsequent realization on complete histories, whose isometry is proved by retaining the label of each prolongation. This identity and the increment \eqref{eq:pdfv2-cor-b-alpha} are not two transports: \(A_k,\eta_k\) act on the history basis and \(b_\alpha\) reads the edge of that same basis.

\subsection{Signed polarization, non-autonomous telescoping, and return without reinitialization}
No independent spectral branch is introduced here. The notations are
restricted to the same space \(\mathscr H_k\) and the same already constructed
transport \(U_k\)
\[
 \mathcal H_k^{\mathrm{sp}}:=\mathscr H_k,
 \qquad
 \Lambda(\gamma_{\widehat x})
 :=\bigl(\operatorname{or},\operatorname{Carr},\operatorname{Mem},
          \partial,\gamma\bigr)(\widehat x),
\]
and \(\Sigma_{\rm reg}\Lambda(\gamma_{\widehat x})\) denotes its family of
compatible truncations. Thus, the following signed reader, defect, and
terminal are actions of the common correlative constructor, not a
retrospective selection of histories.
\begingroup
  \let\section\subsection
  % Extracción quirúrgica del propietario V2 05a: líneas 293--394.
% El resto permanece en este archivo como procedencia, pero no se publica.
\iffalse
\chapter{Spectral macrostructure--internal polar continuum}
\label{chap:continuo-macro-sp-interna}

\section{Strictly internal scope}

The construction in this chapter terminates before any external recognition. Its sole data are the two-sheet APP support, the TRIT orientation, the TPK extensions and carry, the projective system of histories, the cylindrical measure, and the signed ledger. The result is a graded macrostructure with transport, a sheet-change expression, a signed reader, and nonautonomous telescoping. No external function is used to define any of its maps.

\section{Surviving domain of the branch}

For \(x_k\in\mathcal A_k\), let
\begin{equation}
\label{eq:continuo-sp-msect-n}
 \operatorname{Msect}^{\mathrm{sp},(N)}_k(x_k)
 =\left\{(x_j)_{j=k}^{N}:
 \tau_{j+1,j}x_{j+1}=x_j,
 \ (x_j,x_{j+9})\in\Gamma_{9,j},
 \ \text{and the sheaf is preserved }\{\Sigma,\Pi\}\right\}.
\end{equation}
We define
\begin{equation}
\label{eq:continuo-sp-dominio-superviviente}
 \mathcal C_k^{\mathrm{sp}}
 =\left\{x_k:
 \operatorname{Msect}^{\mathrm{sp},(N)}_k(x_k)\neq\varnothing
 \text{ for every }N\geq k\right\},
 \qquad
 \mathcal C_{\rm cont}^{\mathrm{sp}}
 =\varprojlim_k\mathcal C_k^{\mathrm{sp}}.
\end{equation}
The TRIT threshold is not excluded from this domain. The bundle preserves the two
arithmetic sheets even if a visible chart publishes neither of them in a
specific event.

\section{Thirteen-coordinate package}

For \(x_k\in\mathcal C_k^{\mathrm{sp}}\), the dependent package is
\begin{equation}
\label{eq:continuo-sp-d13}
\begin{aligned}
 D_{\mathrm{sp},k}(x_k)=\bigl(&
 \mathcal F_{\mathrm{sp},k}(x_k),
 \operatorname{Ext}_{\Gamma,\mathrm{sp},k}(x_k),
 \operatorname{Rel}_{\mathrm{sp},k}(x_k),
 \mathcal N_{\mathrm{sp},k}(x_k),
 \operatorname{or}_{\mathrm{sp},k}(x_k),\\
 &\operatorname{Carr}_{\mathrm{sp},k}(x_k),
 \operatorname{Mem}_{\mathrm{sp},k}(x_k),
 \partial_{\mathrm{sp},k}(x_k),
 \gamma_{\mathrm{sp},k}(x_k),
 \operatorname{Msect}_{\mathrm{sp},k}(x_k),\\
 &\operatorname{Surv}_{\mathrm{sp},k}(x_k),
 \operatorname{Term}_{\mathrm{sp},k}(x_k),
 \mathcal L^{\rm sgn}_{\mathrm{sp},k}(x_k)\bigr).
\end{aligned}
\end{equation}
Their components are determined correlatively as follows.

The fiber \(\mathcal F_{\mathrm{sp},k}\) contains the TPK state and the ordered
bundle \(\{\Sigma,\Pi\}\). The extension is the complete restriction of
\(\operatorname{Ext}_\Gamma\) to histories satisfying
\eqref{eq:continuo-sp-dominio-superviviente}. The relations contain the
identity and composition laws, the residue--quotient lift, the carry
cocycle, boundary compatibility and truncation naturality. The
numerical coordinate contains depth, phase, residues, quotients and internal
signature. Orientation stores sheet and TRIT orientation separately.
Memory contains the quotient, the ledger and the turn increment. The
boundary and route remain uncontracted. The multisection is the complete
system of prolongations of \eqref{eq:continuo-sp-msect-n}.
Survival is the stable membership of
\eqref{eq:continuo-sp-dominio-superviviente}. The terminal and reader are
constructed in the following sections.

The truncations of all coordinates form a map
\(u^{\mathrm{sp}}_{\ell,k}\) satisfying
\begin{equation}
\label{eq:continuo-sp-d-naturalidad}
 u^{\mathrm{sp}}_{\ell,k}D_{\mathrm{sp},\ell}
 =D_{\mathrm{sp},k}\tau_{\ell,k}.
\end{equation}

\section{Graph of the constraint and dependent section}

We define
\begin{equation}
\label{eq:continuo-sp-g-graph}
 \mathcal G_{\mathrm{sp},k}
 =\operatorname{Graph}(D_{\mathrm{sp},k})
 =\{(x_k,D_{\mathrm{sp},k}x_k):
     x_k\in\mathcal C_k^{\mathrm{sp}}\}.
\end{equation}
The total restriction is
\begin{equation}
\label{eq:continuo-sp-res-graph}
 \operatorname{Res}_{\mathrm{sp}}=s_{\mathrm{sp}}:
 \mathcal C_{\rm cont}^{\mathrm{sp}}
 \longrightarrow\mathcal G_{\mathrm{sp}},
 \qquad
 s_{\mathrm{sp}}(x)=(x,D_{\mathrm{sp}}x).
\end{equation}
The projection \(\pi_1(x,d)=x\) verifies
\begin{equation}
\label{eq:continuo-sp-res-inversa}
 \pi_1s_{\mathrm{sp}}=I,
 \qquad
 s_{\mathrm{sp}}\pi_1=I_{\mathcal G_{\mathrm{sp}}}.
\end{equation}

Let \(r_{\mathrm{sp},k}(x_k)\) be the multisection of sheets and paths that already
belongs to \(D_{\mathrm{sp},k}(x_k)\). Define
\begin{equation}
\label{eq:continuo-sp-x-dependiente}
 X_{\chi,\mathrm{sp},k}
 =\{(r_{\mathrm{sp},k}x_k,x_k,D_{\mathrm{sp},k}x_k):
     x_k\in\mathcal C_k^{\mathrm{sp}}\},
\end{equation}
\begin{equation}
\label{eq:continuo-sp-pr-x}
 \operatorname{pr}_{X,\mathrm{sp}}(x,D_{\mathrm{sp}}x)
 =(r_{\mathrm{sp}}x,x,D_{\mathrm{sp}}x).
\end{equation}
Forgetting the first coordinate of
\eqref{eq:continuo-sp-x-dependiente} is the inverse of
\(\operatorname{pr}_{X,\mathrm{sp}}\).

\section{Bundle of leaves and projective measure}

Let
\begin{equation}
\label{eq:continuo-sp-haz-hojas}
 \widetilde{\mathcal C}_k^{\mathrm{sp}}
 =\{(x_k,h):x_k\in\mathcal C_k^{\mathrm{sp}},
              h\in\{\Sigma,\Pi\}\}.
\end{equation}
The projective measure of the continuum induces probabilities
\(\widetilde\mu_k\) on this bundle. We restrict to the positive support; for
\(z\) in that support,
\begin{equation}
\label{eq:continuo-sp-consistencia-medida}
 \widetilde\mu_k(z)
 =\sum_{\widetilde\tau_{k+1,k}(z')=z}
  \widetilde\mu_{k+1}(z').
\end{equation}
Level linearization is
\begin{equation}
\label{eq:continuo-sp-hilbert-nivel}
 \mathcal H_k^{\mathrm{sp}}
 =\ell^2(\operatorname{supp}\widetilde\mu_k),
\end{equation}
with orthonormal basis \(e_z\).

\section{Transport of all extensions and carry}

For \(z'\) extending \(z\), we define the positive conditional weight
\begin{equation}
\label{eq:continuo-sp-peso-condicional}
 w_k(z'\mid z)
 =\left(\frac{\widetilde\mu_{k+1}(z')}
               {\widetilde\mu_k(z)}\right)^{1/2}.
\end{equation}
Transport is
\begin{equation}
\label{eq:continuo-sp-u-accion}
 U_ke_z
 =\sum_{\substack{z':\widetilde\tau_{k+1,k}(z')=z\\
            (z,z')\in\operatorname{Ext}_{\Gamma,\mathrm{sp},k}}}
 w_k(z'\mid z)e_{z'}.
\end{equation}
The sum contains all admissible extensions. It does not depend on a selector.

By \eqref{eq:continuo-sp-consistencia-medida},
\[
 \sum_{z':\widetilde\tau(z')=z}|w_k(z'\mid z)|^2=1.
\]
The sets of descendants of two distinct parents are disjoint. Consequently,
\begin{equation}
\label{eq:continuo-sp-u-isometria}
 U_k^*U_k=I_{\mathcal H_k^{\mathrm{sp}}}.
\end{equation}

Each basis \(e_{z'}\) includes the carry of the corresponding extension. For
composable paths, we preserve
\begin{align}
 \kappa(\gamma_2\gamma_1)
 &=\kappa(\gamma_2)+
   \mathsf C(\gamma_2)\kappa(\gamma_1),
 \label{eq:continuo-sp-carry-nativo}\\
 \operatorname{Carr}(\gamma_2\gamma_1)
 &=\operatorname{Carr}(\gamma_2)H(\gamma_1)
  +Y(\gamma_2)\operatorname{Carr}(\gamma_1).
 \label{eq:continuo-sp-carry-operador}
\end{align}
The conditional weights multiply along the unique chain of
truncations of a descendant. Equations
\eqref{eq:continuo-sp-carry-nativo}--
\eqref{eq:continuo-sp-carry-operador} ensure that the carry labels
coincide. Therefore,
\begin{equation}
\label{eq:continuo-sp-u-composicion}
 U_{\ell,k}=U_{\ell-1}\cdots U_k.
\end{equation}

\section{Graduation, projectors and decomposition of transport}

The sheet graduation is the involution.
\begin{equation}
\label{eq:continuo-sp-sigma}
 \sigma_ke_{(x,\Sigma)}=+e_{(x,\Sigma)},
 \qquad
 \sigma_ke_{(x,\Pi)}=-e_{(x,\Pi)}.
\end{equation}
Thus,
\[
 \sigma_k^*=\sigma_k,
 \qquad
 \sigma_k^2=I.
\]
The internal projectors are
\begin{equation}
\label{eq:continuo-sp-p-q}
 P_k=\frac{I+\sigma_k}{2},
 \qquad
 Q_k=\frac{I-\sigma_k}{2}.
\end{equation}
Satisfy
\begin{equation}
\label{eq:continuo-sp-p-q-relaciones}
 P_k^2=P_k,
 \quad Q_k^2=Q_k,
 \quad P_kQ_k=0,
 \quad P_k+Q_k=I.
\end{equation}

We separate the transport part that preserves sheet from the one that changes it:
\begin{align}
 A_k&=P_{k+1}U_kP_k+Q_{k+1}U_kQ_k,
 \label{eq:continuo-sp-a}\\
 \eta_k&=Q_{k+1}U_kP_k+P_{k+1}U_kQ_k.
 \label{eq:continuo-sp-eta}
\end{align}
Then
\begin{equation}
\label{eq:continuo-sp-u-a-eta}
 U_k=A_k+\eta_k,
\end{equation}
and the graduation verifies
\begin{equation}
\label{eq:continuo-sp-paridad-a-eta}
 \sigma_{k+1}A_k=A_k\sigma_k,
 \qquad
 \sigma_{k+1}\eta_k=-\eta_k\sigma_k.
\end{equation}

\begin{proposition}[Leaf local balance]
\label{prop:continuo-sp-balance-local}
For each level,
\begin{equation}
\label{eq:continuo-sp-balance-local}
 A_k^*A_k+\eta_k^*\eta_k=I.
\end{equation}
\end{proposition}

\begin{proof}
The orthogonality of \(P_{k+1}\) and \(Q_{k+1}\) gives
\begin{align*}
 A_k^*A_k
 &=P_kU_k^*P_{k+1}U_kP_k
   +Q_kU_k^*Q_{k+1}U_kQ_k,\\
 \eta_k^*\eta_k
 &=P_kU_k^*Q_{k+1}U_kP_k
   +Q_kU_k^*P_{k+1}U_kQ_k.
\end{align*}
Adding and using \(P_{k+1}+Q_{k+1}=I\),
\[
 A_k^*A_k+\eta_k^*\eta_k
 =P_kU_k^*U_kP_k+Q_kU_k^*U_kQ_k
 =P_k+Q_k=I.
\]
\end{proof}

The operator \(\eta_k\) records departure from the sheet-preserving channel. It
is not the integer carry: the coordinates
\eqref{eq:continuo-sp-carry-nativo}--
\eqref{eq:continuo-sp-carry-operador} remain in the package
\eqref{eq:continuo-sp-d13}.

\fi
\section{Signed reader with genealogical register}

For \(v\in\mathcal H_k^{\mathrm{sp}}\), the signed polarization is
\begin{equation}
\label{eq:continuo-sp-polarizacion-firmada}
 \operatorname{pol}_k(v)
 =\langle v,\sigma_kv\rangle
 =\|P_kv\|^2-\|Q_kv\|^2.
\end{equation}
The specialized reader furthermore preserves the source:
\begin{equation}
\label{eq:continuo-sp-lector-firmado}
 \mathcal L^{\rm sgn}_{\mathrm{sp},k}(x_k;v)
 =\bigl(
   \Lambda(\gamma_{x_k}),
   \Sigma_{\rm reg}(\Lambda(\gamma_{x_k})),
   \|P_kv\|^2,
   \|Q_kv\|^2,
   \operatorname{pol}_k(v)\bigr).
\end{equation}
The TRIT orientation is preserved in \(\Lambda(\gamma_{x_k})\); it is not
identified with the sheet sign. At horizon \(N\), the record keeps the
family of readers for all levels \(k\leq N\). Truncating means
removing only subsequent entries.

\section{Non-autonomous and terminal product}

We define
\begin{equation}
\label{eq:continuo-sp-f-producto}
 F_{0,0}=I_{\mathcal H_0^{\mathrm{sp}}},
 \qquad
 F_{0,k}=A_{k-1}\cdots A_1A_0
 \quad(k\geq1).
\end{equation}
Therefore,
\begin{equation}
\label{eq:continuo-sp-f-recursion}
 F_{0,k+1}=A_kF_{0,k}.
\end{equation}
The term published upon leaving the channel at step \(k\) and the terminal at
horizon \(N\) are
\begin{equation}
\label{eq:continuo-sp-defecto-terminal}
 \mathcal D_k=F_{0,k}^*\eta_k^*\eta_kF_{0,k},
 \qquad
 \mathcal T_N=F_{0,N}^*F_{0,N}.
\end{equation}

\begin{theorem}[Spectral telescope--non-autonomous polar]
\label{thm:continuo-sp-telescopia-no-autonoma}
For every finite horizon \(N\),
\begin{equation}
\label{eq:continuo-sp-telescopia-no-autonoma}
 \boxed{
 I=\sum_{k=0}^{N-1}
 F_{0,k}^*\eta_k^*\eta_kF_{0,k}
 +F_{0,N}^*F_{0,N}.}
\end{equation}
Moreover,
\[
 0\preceq\mathcal T_{N+1}
 \preceq\mathcal T_N\preceq I,
\]
so that the strong positive limit exists
\[
 \mathcal T_\infty
 =\operatorname{s-lim}_{N\to\infty}\mathcal T_N
\]
and
\begin{equation}
\label{eq:continuo-sp-telescopia-infinita}
 I=\sum_{k=0}^{\infty}
 F_{0,k}^*\eta_k^*\eta_kF_{0,k}
 +\mathcal T_\infty
\end{equation}
in the strong topology. The theorem does not assert that
\(\mathcal T_\infty\) is zero.
\end{theorem}

\begin{proof}
Multiplying \eqref{eq:pdfv2-cor-balance-hojas} by
\(F_{0,k}^*\) and \(F_{0,k}\), and using
\eqref{eq:continuo-sp-f-recursion}, one obtains
\[
 F_{0,k}^*F_{0,k}
 =F_{0,k+1}^*F_{0,k+1}
  +F_{0,k}^*\eta_k^*\eta_kF_{0,k}.
\]
The sum from \(k=0\) to \(N-1\) cancels all intermediate
terms and leaves precisely
\eqref{eq:continuo-sp-telescopia-no-autonoma}. Since
\(A_k^*A_k\preceq I\),
\[
 \mathcal T_{N+1}
 =F_{0,N}^*A_N^*A_NF_{0,N}
 \preceq F_{0,N}^*F_{0,N}=\mathcal T_N.
\]
Monotonicity and the positive bound provide the strong limit. Passing to the
limit in the finite identity gives
\eqref{eq:continuo-sp-telescopia-infinita}.
\end{proof}

\iffalse
\section{Assembler as graph}

For
\[
 z_N=(r_{\mathrm{sp},N}x_N,x_N,D_{\mathrm{sp},N}x_N)
 \in X_{\chi,\mathrm{sp},N},
\]
the material ensemble is defined
\begin{equation}
\label{eq:continuo-sp-a-accion}
\begin{aligned}
 a_{\mathrm{sp},N}(z_N)=\bigl(&
 (\mathcal H_k^{\mathrm{sp}},\sigma_k,P_k,Q_k)_{k\leq N},
 (U_k,A_k,\eta_k)_{k<N},\\
 &(F_{0,k},\mathcal D_k,\mathcal T_k,
   \mathcal L^{\rm sgn}_{\mathrm{sp},k})_{k\leq N}\bigr).
\end{aligned}
\end{equation}
The macrostructure is
\begin{equation}
\label{eq:continuo-sp-m-graph}
 \mathfrak M_{\mathrm{sp},N}
 =\operatorname{Graph}(a_{\mathrm{sp},N}),
 \qquad
 \mathcal A_{\mathrm{sp},N}(z_N)
 =(z_N,a_{\mathrm{sp},N}z_N).
\end{equation}
Projection onto \(z_N\) is the inverse of
\(\mathcal A_{\mathrm{sp},N}\). Truncation removes the entries whose
index exceeds the new horizon and satisfies
\begin{equation}
\label{eq:continuo-sp-a-naturalidad}
 m^{\mathrm{sp}}_{\ell,k}a_{\mathrm{sp},\ell}
 =a_{\mathrm{sp},k}\xi^{\mathrm{sp}}_{\ell,k}.
\end{equation}

\section{Publisher as graph}

The publisher does not contract the macrostructure to a scalar. Return the
complete certificate.
\begin{equation}
\label{eq:continuo-sp-p-accion}
 p_{\mathrm{sp},N}(m_N)
 =\left(
 \mathcal L^{\rm sgn}_{\mathrm{sp},k},
 \mathcal D_k,
 \mathcal T_k,
 I=\sum_{j=0}^{k-1}\mathcal D_j+\mathcal T_k
 \right)_{0\leq k\leq N}.
\end{equation}
We define
\begin{equation}
\label{eq:continuo-sp-c-graph}
 \mathcal C_{\mathrm{sp},N}
 =\operatorname{Graph}(p_{\mathrm{sp},N}),
 \qquad
 \mathcal P_{\mathrm{sp},N}(m_N)
 =(m_N,p_{\mathrm{sp},N}m_N).
\end{equation}
Projection onto \(m_N\) is its inverse. Truncating the certificate removes only
subsequent identities, so
\begin{equation}
\label{eq:continuo-sp-p-naturalidad}
 c^{\mathrm{sp}}_{\ell,k}p_{\mathrm{sp},\ell}
 =p_{\mathrm{sp},k}m^{\mathrm{sp}}_{\ell,k}.
\end{equation}

\section{Theorem of the complete internal chain}

\begin{theorem}[Spectral-polar macrostructure generated by the continuum]
\label{thm:continuo-sp-cadena-interna}
The actions
\eqref{eq:continuo-sp-res-graph},
\eqref{eq:continuo-sp-pr-x},
\eqref{eq:continuo-sp-m-graph}, and
\eqref{eq:continuo-sp-c-graph} form the natural chain
\begin{equation}
\label{eq:continuo-sp-cadena-interna}
 \mathcal C_{\rm cont}^{\rm disc}
 \supset\mathcal C_{\rm cont}^{\mathrm{sp}}
 \xrightarrow[\cong]{\operatorname{Res}_{\mathrm{sp}}}
 \mathcal G_{\mathrm{sp}}
 \xrightarrow[\cong]{\operatorname{pr}_{X,\mathrm{sp}}}
 X_{\chi,\mathrm{sp}}
 \xrightarrow[\cong]{\mathcal A_{\mathrm{sp}}}
 \mathfrak M_{\mathrm{sp}}
 \xrightarrow[\cong]{\mathcal P_{\mathrm{sp}}}
 \mathcal C_{\mathrm{sp}}.
\end{equation}
Its published content consists of the two-sheet bundle, the grading
\(\sigma_k\), the projectors \(P_k,Q_k\), the total transport \(U_k\),
the separation \(U_k=A_k+\eta_k\), the carry and the memory, the signed reader,
and the telescoping structure
\eqref{eq:continuo-sp-telescopia-no-autonoma} with an explicit terminal object. Each
object preserves a canonical projection to all its antecedents.
\end{theorem}

\begin{proof}
The identities
\eqref{eq:continuo-sp-res-inversa} prove the first equivalence. The
dependent definition \eqref{eq:continuo-sp-x-dependiente} proves the
second. The graphs
\eqref{eq:continuo-sp-m-graph} and
\eqref{eq:continuo-sp-c-graph} prove the remaining two. Naturality of
the four steps is
\eqref{eq:continuo-sp-d-naturalidad},
\eqref{eq:continuo-sp-a-naturalidad}, and
\eqref{eq:continuo-sp-p-naturalidad}, together with naturality of
\(r_{\mathrm{sp},k}\). The operatorial content is proved by
\eqref{eq:continuo-sp-u-isometria}, the
\cref{prop:continuo-sp-balance-local}, and the
\cref{thm:continuo-sp-telescopia-no-autonoma}. All extensions appear
in \eqref{eq:continuo-sp-u-accion}; consequently, no choice of
route intervenes.
\end{proof}

\section{Phase return and state evolution}

The branch explicitly preserves
\begin{equation}
\label{eq:continuo-sp-no-reset}
 g(k+9)=g(k),
 \qquad
 x_{k+9}\neq x_k\quad\text{in general}.
\end{equation}
The second member records the increase in memory and possible changes
of sheet, carry, cylinder and boundary. Neither \(\sigma_k\), nor the projectors, nor
the signed reader contract this difference.

\section{Origin}

The two APP sheets, the enriched TRIT state, the TPK fibers and extensions,
the carry, return without resetting, survival and the ledger are
\texttt{PREEXISTING\_AUTHORIAL\_ARCHITECTURE}. The obligation to preserve the
finite terminal is a \texttt{RECOVERED\_RESULT}. The explicit bundle of sheets,
the involution \(\sigma_k\), the projectors, the linearization that sums all
extensions, the decomposition \(U_k=A_k+\eta_k\), the nonautonomous
product \(F_{0,k}\) and the dependent graphs are a
\texttt{NEW\_FORMALIZATION}. The
\cref{thm:continuo-sp-telescopia-no-autonoma,thm:continuo-sp-cadena-interna}
constitutes the \texttt{NEW\_CERTIFICATE} of that formalization.

\section{Falsadores internos}

\begin{proposition}[Falsifiers of the macrostructure]
\label{prop:continuo-sp-falsadores}
The conclusion of
\cref{thm:continuo-sp-cadena-interna} is not transferable if any
of the following substitutions occurs:
\begin{enumerate}
 \item choosing a prolongation in
       \eqref{eq:continuo-sp-u-accion} and deleting the remaining ones;
 \item replacing \(\mathcal G_{\mathrm{sp}}\) with the external image of
       \(D_{\mathrm{sp}}\);
 \item deleting \((x,D_{\mathrm{sp}}x)\) when forming
       \(X_{\chi,\mathrm{sp}}\);
 \item identifying the sheet with the TRIT orientation or deleting the threshold;
 \item breaking the consistency
       \eqref{eq:continuo-sp-consistencia-medida};
 \item deleting the carry or violating
       \eqref{eq:continuo-sp-carry-nativo}--
       \eqref{eq:continuo-sp-carry-operador};
 \item identifying \(\eta_k\) with the entire carry;
 \item replacing the family \((A_k)_k\) with a single operator without proving
       stationarity;
 \item removing \(F_{0,N}^*F_{0,N}\) at a finite horizon;
 \item asserting \(\mathcal T_\infty=0\) without additional proof;
 \item reconstructing a history from the signed polarization;
 \item using \(g(k+9)=g(k)\) to assert \(x_{k+9}=x_k\);
 \item introducing an external domain, function, or criterion to
define any of the internal arrows;
 \item introducing a subsequent output as an antecedent of the branch.
\end{enumerate}
\end{proposition}
\fi

  % Extracción quirúrgica del propietario V2 05a: líneas 512--523.
% El resto permanece en este archivo como procedencia, pero no se publica.
\iffalse
\chapter{Spectral macrostructure--internal polar continuum}
\label{chap:continuo-macro-sp-interna}

\section{Strictly internal scope}

The construction in this chapter terminates before any external recognition. Its sole data are the two-sheet APP support, the TRIT orientation, the TPK extensions and carry, the projective system of histories, the cylindrical measure, and the signed ledger. The result is a graded macrostructure with transport, a sheet-change expression, a signed reader, and nonautonomous telescoping. No external function is used to define any of its maps.

\section{Surviving domain of the branch}

For \(x_k\in\mathcal A_k\), let
\begin{equation}
\label{eq:continuo-sp-msect-n}
 \operatorname{Msect}^{\mathrm{sp},(N)}_k(x_k)
 =\left\{(x_j)_{j=k}^{N}:
 \tau_{j+1,j}x_{j+1}=x_j,
 \ (x_j,x_{j+9})\in\Gamma_{9,j},
 \ \text{and the sheaf is preserved }\{\Sigma,\Pi\}\right\}.
\end{equation}
We define
\begin{equation}
\label{eq:continuo-sp-dominio-superviviente}
 \mathcal C_k^{\mathrm{sp}}
 =\left\{x_k:
 \operatorname{Msect}^{\mathrm{sp},(N)}_k(x_k)\neq\varnothing
 \text{ for every }N\geq k\right\},
 \qquad
 \mathcal C_{\rm cont}^{\mathrm{sp}}
 =\varprojlim_k\mathcal C_k^{\mathrm{sp}}.
\end{equation}
The TRIT threshold is not excluded from this domain. The bundle preserves the two
arithmetic sheets even if a visible chart publishes neither of them in a
specific event.

\section{Thirteen-coordinate package}

For \(x_k\in\mathcal C_k^{\mathrm{sp}}\), the dependent package is
\begin{equation}
\label{eq:continuo-sp-d13}
\begin{aligned}
 D_{\mathrm{sp},k}(x_k)=\bigl(&
 \mathcal F_{\mathrm{sp},k}(x_k),
 \operatorname{Ext}_{\Gamma,\mathrm{sp},k}(x_k),
 \operatorname{Rel}_{\mathrm{sp},k}(x_k),
 \mathcal N_{\mathrm{sp},k}(x_k),
 \operatorname{or}_{\mathrm{sp},k}(x_k),\\
 &\operatorname{Carr}_{\mathrm{sp},k}(x_k),
 \operatorname{Mem}_{\mathrm{sp},k}(x_k),
 \partial_{\mathrm{sp},k}(x_k),
 \gamma_{\mathrm{sp},k}(x_k),
 \operatorname{Msect}_{\mathrm{sp},k}(x_k),\\
 &\operatorname{Surv}_{\mathrm{sp},k}(x_k),
 \operatorname{Term}_{\mathrm{sp},k}(x_k),
 \mathcal L^{\rm sgn}_{\mathrm{sp},k}(x_k)\bigr).
\end{aligned}
\end{equation}
Their components are determined correlatively as follows.

The fiber \(\mathcal F_{\mathrm{sp},k}\) contains the TPK state and the ordered
bundle \(\{\Sigma,\Pi\}\). The extension is the complete restriction of
\(\operatorname{Ext}_\Gamma\) to histories satisfying
\eqref{eq:continuo-sp-dominio-superviviente}. The relations contain the
identity and composition laws, the residue--quotient lift, the carry
cocycle, boundary compatibility and truncation naturality. The
numerical coordinate contains depth, phase, residues, quotients and internal
signature. Orientation stores sheet and TRIT orientation separately.
Memory contains the quotient, the ledger and the turn increment. The
boundary and route remain uncontracted. The multisection is the complete
system of prolongations of \eqref{eq:continuo-sp-msect-n}.
Survival is the stable membership of
\eqref{eq:continuo-sp-dominio-superviviente}. The terminal and reader are
constructed in the following sections.

The truncations of all coordinates form a map
\(u^{\mathrm{sp}}_{\ell,k}\) satisfying
\begin{equation}
\label{eq:continuo-sp-d-naturalidad}
 u^{\mathrm{sp}}_{\ell,k}D_{\mathrm{sp},\ell}
 =D_{\mathrm{sp},k}\tau_{\ell,k}.
\end{equation}

\section{Graph of the constraint and dependent section}

We define
\begin{equation}
\label{eq:continuo-sp-g-graph}
 \mathcal G_{\mathrm{sp},k}
 =\operatorname{Graph}(D_{\mathrm{sp},k})
 =\{(x_k,D_{\mathrm{sp},k}x_k):
     x_k\in\mathcal C_k^{\mathrm{sp}}\}.
\end{equation}
The total restriction is
\begin{equation}
\label{eq:continuo-sp-res-graph}
 \operatorname{Res}_{\mathrm{sp}}=s_{\mathrm{sp}}:
 \mathcal C_{\rm cont}^{\mathrm{sp}}
 \longrightarrow\mathcal G_{\mathrm{sp}},
 \qquad
 s_{\mathrm{sp}}(x)=(x,D_{\mathrm{sp}}x).
\end{equation}
The projection \(\pi_1(x,d)=x\) verifies
\begin{equation}
\label{eq:continuo-sp-res-inversa}
 \pi_1s_{\mathrm{sp}}=I,
 \qquad
 s_{\mathrm{sp}}\pi_1=I_{\mathcal G_{\mathrm{sp}}}.
\end{equation}

Let \(r_{\mathrm{sp},k}(x_k)\) be the multisection of sheets and paths that already
belongs to \(D_{\mathrm{sp},k}(x_k)\). Define
\begin{equation}
\label{eq:continuo-sp-x-dependiente}
 X_{\chi,\mathrm{sp},k}
 =\{(r_{\mathrm{sp},k}x_k,x_k,D_{\mathrm{sp},k}x_k):
     x_k\in\mathcal C_k^{\mathrm{sp}}\},
\end{equation}
\begin{equation}
\label{eq:continuo-sp-pr-x}
 \operatorname{pr}_{X,\mathrm{sp}}(x,D_{\mathrm{sp}}x)
 =(r_{\mathrm{sp}}x,x,D_{\mathrm{sp}}x).
\end{equation}
Forgetting the first coordinate of
\eqref{eq:continuo-sp-x-dependiente} is the inverse of
\(\operatorname{pr}_{X,\mathrm{sp}}\).

\section{Bundle of leaves and projective measure}

Let
\begin{equation}
\label{eq:continuo-sp-haz-hojas}
 \widetilde{\mathcal C}_k^{\mathrm{sp}}
 =\{(x_k,h):x_k\in\mathcal C_k^{\mathrm{sp}},
              h\in\{\Sigma,\Pi\}\}.
\end{equation}
The projective measure of the continuum induces probabilities
\(\widetilde\mu_k\) on this bundle. We restrict to the positive support; for
\(z\) in that support,
\begin{equation}
\label{eq:continuo-sp-consistencia-medida}
 \widetilde\mu_k(z)
 =\sum_{\widetilde\tau_{k+1,k}(z')=z}
  \widetilde\mu_{k+1}(z').
\end{equation}
Level linearization is
\begin{equation}
\label{eq:continuo-sp-hilbert-nivel}
 \mathcal H_k^{\mathrm{sp}}
 =\ell^2(\operatorname{supp}\widetilde\mu_k),
\end{equation}
with orthonormal basis \(e_z\).

\section{Transport of all extensions and carry}

For \(z'\) extending \(z\), we define the positive conditional weight
\begin{equation}
\label{eq:continuo-sp-peso-condicional}
 w_k(z'\mid z)
 =\left(\frac{\widetilde\mu_{k+1}(z')}
               {\widetilde\mu_k(z)}\right)^{1/2}.
\end{equation}
Transport is
\begin{equation}
\label{eq:continuo-sp-u-accion}
 U_ke_z
 =\sum_{\substack{z':\widetilde\tau_{k+1,k}(z')=z\\
            (z,z')\in\operatorname{Ext}_{\Gamma,\mathrm{sp},k}}}
 w_k(z'\mid z)e_{z'}.
\end{equation}
The sum contains all admissible extensions. It does not depend on a selector.

By \eqref{eq:continuo-sp-consistencia-medida},
\[
 \sum_{z':\widetilde\tau(z')=z}|w_k(z'\mid z)|^2=1.
\]
The sets of descendants of two distinct parents are disjoint. Consequently,
\begin{equation}
\label{eq:continuo-sp-u-isometria}
 U_k^*U_k=I_{\mathcal H_k^{\mathrm{sp}}}.
\end{equation}

Each basis \(e_{z'}\) includes the carry of the corresponding extension. For
composable paths, we preserve
\begin{align}
 \kappa(\gamma_2\gamma_1)
 &=\kappa(\gamma_2)+
   \mathsf C(\gamma_2)\kappa(\gamma_1),
 \label{eq:continuo-sp-carry-nativo}\\
 \operatorname{Carr}(\gamma_2\gamma_1)
 &=\operatorname{Carr}(\gamma_2)H(\gamma_1)
  +Y(\gamma_2)\operatorname{Carr}(\gamma_1).
 \label{eq:continuo-sp-carry-operador}
\end{align}
The conditional weights multiply along the unique chain of
truncations of a descendant. Equations
\eqref{eq:continuo-sp-carry-nativo}--
\eqref{eq:continuo-sp-carry-operador} ensure that the carry labels
coincide. Therefore,
\begin{equation}
\label{eq:continuo-sp-u-composicion}
 U_{\ell,k}=U_{\ell-1}\cdots U_k.
\end{equation}

\section{Graduation, projectors and decomposition of transport}

The sheet graduation is the involution.
\begin{equation}
\label{eq:continuo-sp-sigma}
 \sigma_ke_{(x,\Sigma)}=+e_{(x,\Sigma)},
 \qquad
 \sigma_ke_{(x,\Pi)}=-e_{(x,\Pi)}.
\end{equation}
Thus,
\[
 \sigma_k^*=\sigma_k,
 \qquad
 \sigma_k^2=I.
\]
The internal projectors are
\begin{equation}
\label{eq:continuo-sp-p-q}
 P_k=\frac{I+\sigma_k}{2},
 \qquad
 Q_k=\frac{I-\sigma_k}{2}.
\end{equation}
Satisfy
\begin{equation}
\label{eq:continuo-sp-p-q-relaciones}
 P_k^2=P_k,
 \quad Q_k^2=Q_k,
 \quad P_kQ_k=0,
 \quad P_k+Q_k=I.
\end{equation}

We separate the transport part that preserves sheet from the one that changes it:
\begin{align}
 A_k&=P_{k+1}U_kP_k+Q_{k+1}U_kQ_k,
 \label{eq:continuo-sp-a}\\
 \eta_k&=Q_{k+1}U_kP_k+P_{k+1}U_kQ_k.
 \label{eq:continuo-sp-eta}
\end{align}
Then
\begin{equation}
\label{eq:continuo-sp-u-a-eta}
 U_k=A_k+\eta_k,
\end{equation}
and the graduation verifies
\begin{equation}
\label{eq:continuo-sp-paridad-a-eta}
 \sigma_{k+1}A_k=A_k\sigma_k,
 \qquad
 \sigma_{k+1}\eta_k=-\eta_k\sigma_k.
\end{equation}

\begin{proposition}[Leaf local balance]
\label{prop:continuo-sp-balance-local}
For each level,
\begin{equation}
\label{eq:continuo-sp-balance-local}
 A_k^*A_k+\eta_k^*\eta_k=I.
\end{equation}
\end{proposition}

\begin{proof}
The orthogonality of \(P_{k+1}\) and \(Q_{k+1}\) gives
\begin{align*}
 A_k^*A_k
 &=P_kU_k^*P_{k+1}U_kP_k
   +Q_kU_k^*Q_{k+1}U_kQ_k,\\
 \eta_k^*\eta_k
 &=P_kU_k^*Q_{k+1}U_kP_k
   +Q_kU_k^*P_{k+1}U_kQ_k.
\end{align*}
Adding and using \(P_{k+1}+Q_{k+1}=I\),
\[
 A_k^*A_k+\eta_k^*\eta_k
 =P_kU_k^*U_kP_k+Q_kU_k^*U_kQ_k
 =P_k+Q_k=I.
\]
\end{proof}

The operator \(\eta_k\) records departure from the sheet-preserving channel. It
is not the integer carry: the coordinates
\eqref{eq:continuo-sp-carry-nativo}--
\eqref{eq:continuo-sp-carry-operador} remain in the package
\eqref{eq:continuo-sp-d13}.

\section{Signed reader with ledger}

For \(v\in\mathcal H_k^{\mathrm{sp}}\), the signed polarization is
\begin{equation}
\label{eq:continuo-sp-polarizacion-firmada}
 \operatorname{pol}_k(v)
 =\langle v,\sigma_kv\rangle
 =\|P_kv\|^2-\|Q_kv\|^2.
\end{equation}
The specialized reader furthermore preserves the source:
\begin{equation}
\label{eq:continuo-sp-lector-firmado}
 \mathcal L^{\rm sgn}_{\mathrm{sp},k}(x_k;v)
 =\bigl(
   \Lambda(\gamma_{x_k}),
   \Sigma_{\rm reg}(\Lambda(\gamma_{x_k})),
   \|P_kv\|^2,
   \|Q_kv\|^2,
   \operatorname{pol}_k(v)\bigr).
\end{equation}
The TRIT orientation is preserved in \(\Lambda(\gamma_{x_k})\); it is not
identified with the sheet sign. At horizon \(N\), the record keeps the
family of readers for all levels \(k\leq N\). Truncating means
removing only subsequent entries.

\section{Non-autonomous and terminal product}

We define
\begin{equation}
\label{eq:continuo-sp-f-producto}
 F_{0,0}=I_{\mathcal H_0^{\mathrm{sp}}},
 \qquad
 F_{0,k}=A_{k-1}\cdots A_1A_0
 \quad(k\geq1).
\end{equation}
Therefore,
\begin{equation}
\label{eq:continuo-sp-f-recursion}
 F_{0,k+1}=A_kF_{0,k}.
\end{equation}
The term published upon leaving the channel at step \(k\) and the terminal at
horizon \(N\) are
\begin{equation}
\label{eq:continuo-sp-defecto-terminal}
 \mathcal D_k=F_{0,k}^*\eta_k^*\eta_kF_{0,k},
 \qquad
 \mathcal T_N=F_{0,N}^*F_{0,N}.
\end{equation}

\begin{theorem}[Spectral telescope--non-autonomous polar]
\label{thm:continuo-sp-telescopia-no-autonoma}
For every finite horizon \(N\),
\begin{equation}
\label{eq:continuo-sp-telescopia-no-autonoma}
 \boxed{
 I=\sum_{k=0}^{N-1}
 F_{0,k}^*\eta_k^*\eta_kF_{0,k}
 +F_{0,N}^*F_{0,N}.}
\end{equation}
Moreover,
\[
 0\preceq\mathcal T_{N+1}
 \preceq\mathcal T_N\preceq I,
\]
so that the strong positive limit exists
\[
 \mathcal T_\infty
 =\operatorname{s-lim}_{N\to\infty}\mathcal T_N
\]
and
\begin{equation}
\label{eq:continuo-sp-telescopia-infinita}
 I=\sum_{k=0}^{\infty}
 F_{0,k}^*\eta_k^*\eta_kF_{0,k}
 +\mathcal T_\infty
\end{equation}
in the strong topology. The theorem does not assert that
\(\mathcal T_\infty\) is zero.
\end{theorem}

\begin{proof}
Multiplying \eqref{eq:continuo-sp-balance-local} by
\(F_{0,k}^*\) and \(F_{0,k}\), and using
\eqref{eq:continuo-sp-f-recursion}, one obtains
\[
 F_{0,k}^*F_{0,k}
 =F_{0,k+1}^*F_{0,k+1}
  +F_{0,k}^*\eta_k^*\eta_kF_{0,k}.
\]
The sum from \(k=0\) to \(N-1\) cancels all intermediate
terms and leaves precisely
\eqref{eq:continuo-sp-telescopia-no-autonoma}. Since
\(A_k^*A_k\preceq I\),
\[
 \mathcal T_{N+1}
 =F_{0,N}^*A_N^*A_NF_{0,N}
 \preceq F_{0,N}^*F_{0,N}=\mathcal T_N.
\]
Monotonicity and the positive bound provide the strong limit. Passing to the
limit in the finite identity gives
\eqref{eq:continuo-sp-telescopia-infinita}.
\end{proof}

\section{Assembler as graph}

For
\[
 z_N=(r_{\mathrm{sp},N}x_N,x_N,D_{\mathrm{sp},N}x_N)
 \in X_{\chi,\mathrm{sp},N},
\]
the material ensemble is defined
\begin{equation}
\label{eq:continuo-sp-a-accion}
\begin{aligned}
 a_{\mathrm{sp},N}(z_N)=\bigl(&
 (\mathcal H_k^{\mathrm{sp}},\sigma_k,P_k,Q_k)_{k\leq N},
 (U_k,A_k,\eta_k)_{k<N},\\
 &(F_{0,k},\mathcal D_k,\mathcal T_k,
   \mathcal L^{\rm sgn}_{\mathrm{sp},k})_{k\leq N}\bigr).
\end{aligned}
\end{equation}
The macrostructure is
\begin{equation}
\label{eq:continuo-sp-m-graph}
 \mathfrak M_{\mathrm{sp},N}
 =\operatorname{Graph}(a_{\mathrm{sp},N}),
 \qquad
 \mathcal A_{\mathrm{sp},N}(z_N)
 =(z_N,a_{\mathrm{sp},N}z_N).
\end{equation}
Projection onto \(z_N\) is the inverse of
\(\mathcal A_{\mathrm{sp},N}\). Truncation removes the entries whose
index exceeds the new horizon and satisfies
\begin{equation}
\label{eq:continuo-sp-a-naturalidad}
 m^{\mathrm{sp}}_{\ell,k}a_{\mathrm{sp},\ell}
 =a_{\mathrm{sp},k}\xi^{\mathrm{sp}}_{\ell,k}.
\end{equation}

\section{Publisher as graph}

The publisher does not contract the macrostructure to a scalar. Return the
complete certificate.
\begin{equation}
\label{eq:continuo-sp-p-accion}
 p_{\mathrm{sp},N}(m_N)
 =\left(
 \mathcal L^{\rm sgn}_{\mathrm{sp},k},
 \mathcal D_k,
 \mathcal T_k,
 I=\sum_{j=0}^{k-1}\mathcal D_j+\mathcal T_k
 \right)_{0\leq k\leq N}.
\end{equation}
We define
\begin{equation}
\label{eq:continuo-sp-c-graph}
 \mathcal C_{\mathrm{sp},N}
 =\operatorname{Graph}(p_{\mathrm{sp},N}),
 \qquad
 \mathcal P_{\mathrm{sp},N}(m_N)
 =(m_N,p_{\mathrm{sp},N}m_N).
\end{equation}
Projection onto \(m_N\) is its inverse. Truncating the certificate removes only
subsequent identities, so
\begin{equation}
\label{eq:continuo-sp-p-naturalidad}
 c^{\mathrm{sp}}_{\ell,k}p_{\mathrm{sp},\ell}
 =p_{\mathrm{sp},k}m^{\mathrm{sp}}_{\ell,k}.
\end{equation}

\section{Theorem of the complete internal chain}

\begin{theorem}[Spectral-polar macrostructure generated by the continuum]
\label{thm:continuo-sp-cadena-interna}
The actions
\eqref{eq:continuo-sp-res-graph},
\eqref{eq:continuo-sp-pr-x},
\eqref{eq:continuo-sp-m-graph}, and
\eqref{eq:continuo-sp-c-graph} form the natural chain
\begin{equation}
\label{eq:continuo-sp-cadena-interna}
 \mathcal C_{\rm cont}^{\rm disc}
 \supset\mathcal C_{\rm cont}^{\mathrm{sp}}
 \xrightarrow[\cong]{\operatorname{Res}_{\mathrm{sp}}}
 \mathcal G_{\mathrm{sp}}
 \xrightarrow[\cong]{\operatorname{pr}_{X,\mathrm{sp}}}
 X_{\chi,\mathrm{sp}}
 \xrightarrow[\cong]{\mathcal A_{\mathrm{sp}}}
 \mathfrak M_{\mathrm{sp}}
 \xrightarrow[\cong]{\mathcal P_{\mathrm{sp}}}
 \mathcal C_{\mathrm{sp}}.
\end{equation}
Its published content consists of the two-sheet bundle, the grading
\(\sigma_k\), the projectors \(P_k,Q_k\), the total transport \(U_k\),
the separation \(U_k=A_k+\eta_k\), the carry and the memory, the signed reader,
and the telescoping structure
\eqref{eq:continuo-sp-telescopia-no-autonoma} with an explicit terminal object. Each
object preserves a canonical projection to all its antecedents.
\end{theorem}

\begin{proof}
The identities
\eqref{eq:continuo-sp-res-inversa} prove the first equivalence. The
dependent definition \eqref{eq:continuo-sp-x-dependiente} proves the
second. The graphs
\eqref{eq:continuo-sp-m-graph} and
\eqref{eq:continuo-sp-c-graph} prove the remaining two. Naturality of
the four steps is
\eqref{eq:continuo-sp-d-naturalidad},
\eqref{eq:continuo-sp-a-naturalidad}, and
\eqref{eq:continuo-sp-p-naturalidad}, together with naturality of
\(r_{\mathrm{sp},k}\). The operatorial content is proved by
\eqref{eq:continuo-sp-u-isometria}, the
\cref{prop:continuo-sp-balance-local}, and the
\cref{thm:continuo-sp-telescopia-no-autonoma}. All extensions appear
in \eqref{eq:continuo-sp-u-accion}; consequently, no choice of
route intervenes.
\end{proof}

\fi
\section{Phase return and non-reinitialized evolution of the common state}

The branch explicitly preserves
\begin{equation}
\label{eq:continuo-sp-no-reset}
 g(k+9)=g(k),
 \qquad
 x_{k+9}\neq x_k\quad\text{in general}.
\end{equation}
The second member records the increase in memory and possible changes
of sheet, carry, cylinder and boundary. Neither \(\sigma_k\), nor the projectors, nor
the signed reader contract this difference.

\iffalse
\section{Origin}

The two APP sheets, the enriched TRIT state, the TPK fibers and extensions,
the carry, return without resetting, survival and the ledger are
\texttt{PREEXISTING\_AUTHORIAL\_ARCHITECTURE}. The obligation to preserve the
finite terminal is a \texttt{RECOVERED\_RESULT}. The explicit bundle of sheets,
the involution \(\sigma_k\), the projectors, the linearization that sums all
extensions, the decomposition \(U_k=A_k+\eta_k\), the nonautonomous
product \(F_{0,k}\) and the dependent graphs are a
\texttt{NEW\_FORMALIZATION}. The
\cref{thm:continuo-sp-telescopia-no-autonoma,thm:continuo-sp-cadena-interna}
constitutes the \texttt{NEW\_CERTIFICATE} of that formalization.

\section{Falsadores internos}

\begin{proposition}[Falsifiers of the macrostructure]
\label{prop:continuo-sp-falsadores}
The conclusion of
\cref{thm:continuo-sp-cadena-interna} is not transferable if any
of the following substitutions occurs:
\begin{enumerate}
 \item choosing a prolongation in
       \eqref{eq:continuo-sp-u-accion} and deleting the remaining ones;
 \item replacing \(\mathcal G_{\mathrm{sp}}\) with the external image of
       \(D_{\mathrm{sp}}\);
 \item deleting \((x,D_{\mathrm{sp}}x)\) when forming
       \(X_{\chi,\mathrm{sp}}\);
 \item identifying the sheet with the TRIT orientation or deleting the threshold;
 \item breaking the consistency
       \eqref{eq:continuo-sp-consistencia-medida};
 \item deleting the carry or violating
       \eqref{eq:continuo-sp-carry-nativo}--
       \eqref{eq:continuo-sp-carry-operador};
 \item identifying \(\eta_k\) with the entire carry;
 \item replacing the family \((A_k)_k\) with a single operator without proving
       stationarity;
 \item removing \(F_{0,N}^*F_{0,N}\) at a finite horizon;
 \item asserting \(\mathcal T_\infty=0\) without additional proof;
 \item reconstructing a history from the signed polarization;
 \item using \(g(k+9)=g(k)\) to assert \(x_{k+9}=x_k\);
 \item introducing an external domain, function, or criterion to
define any of the internal arrows;
 \item introducing a subsequent output as an antecedent of the branch.
\end{enumerate}
\end{proposition}
\fi

\endgroup

The conditional expectation associated with the same weight is
\begin{equation}
\label{eq:pdfv2-cor-esperanza-comun}
 (E_kF)(\widehat x):=
 \sum_{(\alpha,\widehat y)\in\operatorname{Ch}_k(\widehat x)}
 p_k(\alpha,\widehat y\mid\widehat x)F(\alpha,\widehat y).
\end{equation}
For the fine increment \(b_f\), set
\begin{equation}
\label{eq:pdfv2-cor-kappa-jensen}
 b_c:=E_kb_f,\qquad
 \kappa_k^{\rm can}:=\frac{E_k|b_f|-|b_c|}{2}\geq0.
\end{equation}
Then
\begin{equation}
\label{eq:pdfv2-cor-jensen-hojas}
 E_kb_f^+=b_c^++\kappa_k^{\rm can},\qquad
 E_kb_f^-=b_c^-+\kappa_k^{\rm can}.
\end{equation}
The superscript \({\rm can}\) distinguishes this cancellation carry from the integer
carry \(\kappa_\alpha\); both arise from the same transport, but have
distinct types.

Signed memories require no external initial condition: for
\(\gamma=(\epsilon_1,\ldots,\epsilon_k)\) they are defined by summing the ledger itself,
\begin{equation}
\label{eq:pdfv2-cor-memorias}
 \begin{aligned}
 M_0^\pm&:=0,\\
 M_k^\pm(\gamma)&:=\sum_{j=1}^k b_{\epsilon_j}^\pm,\\
 M_{k+1}^\pm(\gamma\alpha)&:=M_k^\pm(\gamma)+b_\alpha^\pm,\\
 D_k^0&:=M_k^+-M_k^-, &
 H_k^0&:=M_k^++M_k^-.
 \end{aligned}
\end{equation}
Thus, the signed balance, the total memory, and the polarization
\(\langle v,\sigma_kv\rangle\) are three coordinated readings of the same
prefix, not three states reconstructed separately.

\subsection{Nonadic Refinement, Orthogonal Losses and Terminal Memory}
The index \(j\) in the following development is the level of the same tree of
prolongations, and its thirteen-coordinate record is the restriction of the
common record of \eqref{eq:pdfv2-cor-registro-extension}. Accordingly,
\(D_{{\rm sol},j}:=D_j^{(13)}\) is fixed and no additional domain is postulated.
\begingroup
  \let\section\subsection
  % Extracción quirúrgica del propietario V2 05b: líneas 100--359.
% El resto permanece en este archivo como procedencia, pero no se publica.
\iffalse
% Macroestructura solenoidal estrictamente interna del continuo discreto.
% Este módulo no introduce ninguna ecuación clásica ni un problema exterior.

\chapter{Internal macrostructure of balance, memory, and telescopia}
\label{chap:pdfv2-sol-interna}

\begin{statusbox}
The only input of this chapter is the discrete structure of the continuum.
The branch \(\mathrm{sol}\) is constructed before any exterior scalar evaluation
and before any classical differential formulation.  Its complete chain is
\[
 \mathcal C_{\rm cont}^{\rm disc}
 \supset\mathcal C_{\rm cont}^{\rm sol}
 \xrightarrow{\operatorname{Res}_{\rm sol}}
 \mathcal G_{\rm sol}
 \xrightarrow{\operatorname{pr}_{X,{\rm sol}}}
 X_{\chi,{\rm sol}}
 \xrightarrow{\mathcal A_{\rm sol}}
 \mathfrak M_{\rm sol}
 \xrightarrow{\mathcal P_{\rm sol}}
 \mathcal C_{\rm sol}^{\rm HMT}.
\]
Each arrow preserves as coordinates the history and the apparatus that
produces it.  No image in this chain is a quotient that forgets its source.
\end{statusbox}

\section{Proper domain prior to any subsequent evaluation}
\label{sec:pdfv2-sol-dominio}

Let
\[
 x=(x_0\xrightarrow{e_0}x_1\xrightarrow{e_1}\cdots)
 \in\mathcal C_{\rm cont}^{\rm disc}
\]
a surviving history.  Each edge \(e_j\) preserves its two sheets,
orientation, residue, quotient, and integer carry.  Let
\(N_j^+(e_j)\) and \(N_j^-(e_j)\) be the multiplicities of the two sheets, and let
\(\Delta_{\gamma_j}\in\mathbb Z\) be the oriented coefficient of the route, with the
carry of the composition already incorporated.  The internal signed increment is
\begin{equation}\label{eq:pdfv2-sol-b}
 b_j=b(e_j):=
 \Delta_{\gamma_j}\bigl(N_j^+(e_j)-N_j^-(e_j)\bigr)\in\mathbb Z.
\end{equation}
It is not obtained by evaluating an external path: it is a coordinate of the ledger of the edge itself.

Its positive and negative parts are constructed by
\begin{equation}\label{eq:pdfv2-sol-bpm}
 b_j^+:=\frac{|b_j|+b_j}{2},
 \qquad
 b_j^-:=\frac{|b_j|-b_j}{2}.
\end{equation}
Therefore,
\begin{equation}\label{eq:pdfv2-sol-bpm-identidades}
 b_j=b_j^+-b_j^-,
 \qquad
 |b_j|=b_j^++b_j^-,
 \qquad
 b_j^+b_j^-=0.
\end{equation}

The two original memoirs and their two readers are
\begin{align}
 M_0^+=M_0^-&=0,
 &M_j^+&=\sum_{0\leq i<j}b_i^+,
 &M_j^-&=\sum_{0\leq i<j}b_i^-;
 \label{eq:pdfv2-sol-memorias}\\
 D_j^0&:=M_j^+-M_j^-,
 &H_j^0&:=M_j^++M_j^-.
 \label{eq:pdfv2-sol-dh}
\end{align}
Consequently,
\begin{equation}\label{eq:pdfv2-sol-dh-diferencias}
 D_{j+1}^0-D_j^0=b_j,
 \qquad
 H_{j+1}^0-H_j^0=|b_j|.
\end{equation}
Signed memory \(D^0\) and total memory \(H^0\) remain distinct
objects; retaining only their final difference would destroy the genealogy.

The substructure \(\mathcal C_{\rm cont}^{\rm sol}\) consists of the histories
for which, in every prefix and every nonadic refinement:
\begin{enumerate}
 \item the increments \eqref{eq:pdfv2-sol-b} and the two memories
       \eqref{eq:pdfv2-sol-memorias} exist;
 \item the fine cylinders are grouped into fibres of nine extensions with the
       same projective probability law;
 \item the depth and resolution projections defined below are
       compatible with the history extensions;
 \item each loss row and each terminal column is finite at a finite
       horizon;
 \item every prefix admits at least one compatible extension to any
       subsequent horizon.
\end{enumerate}
These are conditions of membership in the branch, not silent premises
attributed to an arbitrary history.

\fi
\section{Common nonadic refinement: inclusion, balance expectation and memory}
\label{sec:pdfv2-sol-e9j9}

Let \(Y_j(x)\) be the finite set of cylinders compatible with the prefix
\(x_{\leq j}\), endowed with the projective measure \(\mu_j\).  The truncation
\[
 \tau_{j+1,j}:Y_{j+1}(x)\longrightarrow Y_j(x)
\]
has nine admissible extensions over each cylinder. Finite spaces are considered.
\[
 \mathscr H_j(x):=\ell^2\bigl(Y_j(x),\mu_j\bigr).
\]
The inclusion and the nonadic expectation are
\begin{align}
 J_{9,j}:\mathscr H_j&\longrightarrow\mathscr H_{j+1},
 &(J_{9,j}f)(y)&=f(\tau_{j+1,j}y),
 \label{eq:pdfv2-sol-j9}\\
 E_{9,j}:\mathscr H_{j+1}&\longrightarrow\mathscr H_j,
 &(E_{9,j}g)(z)&=\frac1{9}
   \sum_{\tau_{j+1,j}y=z}g(y).
 \label{eq:pdfv2-sol-e9}
\end{align}
The projective compatibility of \(\mu_{j+1}\) and \(\mu_j\) gives
\begin{equation}\label{eq:pdfv2-sol-ej-projector}
 E_{9,j}J_{9,j}=I_{\mathscr H_j},
 \qquad
 J_{9,j}^*=E_{9,j},
 \qquad
 P_{9,j}:=J_{9,j}E_{9,j}=P_{9,j}^*=P_{9,j}^2.
\end{equation}
The complementary microscopic projector is
\begin{equation}\label{eq:pdfv2-sol-qmicro-preliminar}
 Q_{9,j}:=I-P_{9,j}.
\end{equation}

If \(b_f\) is the increment on the fine fiber and
\begin{equation}\label{eq:pdfv2-sol-bcoarse}
 b_c:=E_{9,j}b_f,
\end{equation}
the elementary convexity of \(|\cdot|\) constructs, rather than postulates, the carry of
cancellation
\begin{equation}\label{eq:pdfv2-sol-kappa}
 \kappa_j
 :=\frac{E_{9,j}|b_f|-|b_c|}{2}\geq0.
\end{equation}
Indeed, using \(t^\pm=(|t|\pm t)/2\), one obtains
\begin{align}
 E_{9,j}b_f^+
 &=\frac{E_{9,j}|b_f|+E_{9,j}b_f}{2}
   =b_c^++\kappa_j,
 \label{eq:pdfv2-sol-jensen-plus}\\
 E_{9,j}b_f^-
 &=\frac{E_{9,j}|b_f|-E_{9,j}b_f}{2}
   =b_c^-+\kappa_j.
 \label{eq:pdfv2-sol-jensen-minus}
\end{align}
Both sheets lose the same amount when compressed; therefore, the signed
difference is preserved while total memory records the cancellation:
\begin{equation}\label{eq:pdfv2-sol-jensen-dh}
 E_{9,j}(b_f^+-b_f^-)=b_c,
 \qquad
 E_{9,j}(b_f^++b_f^-)=|b_c|+2\kappa_j.
\end{equation}

\section{Orthogonal factorization of the degrees of depth and of the resolution layers in \(\mathscr H_j^{\mathrm{depth}}\widehat\otimes\mathscr H_j^{\mathrm{res}}\)}
\label{sec:pdfv2-sol-micro-shell}

To prevent the same loss from being counted both as depth and as
resolution, the space at each level is typed as
\begin{equation}\label{eq:pdfv2-sol-factor-hilbert}
 \mathscr H_j
 =\mathscr H_j^{\rm depth}\widehat\otimes
  \mathscr H_j^{\rm res}.
\end{equation}
For \(j\geq1\), the first factor carries
\[
 P_j^{\rm depth}:=J_{9,j-1}E_{9,j-1}
 \quad\text{on }\mathscr H_j^{\rm depth},
\]
and \(P_0^{\rm depth}=I\) is fixed. The second factor carries the finite
shell partition inherited from the cylinder boundary. If
\(\sigma_j:Y_j^{\rm res}\to S_j\) is that partition, its operators are
\begin{align}
 J_j^{\rm res}f&=f\circ\sigma_j,
 &E_j^{\rm res}g(s)&=
 \frac1{\#\sigma_j^{-1}(s)}
 \sum_{\sigma_j(y)=s}g(y),
 \label{eq:pdfv2-sol-shell-je}\\
 P_j^{\rm res}&:=J_j^{\rm res}E_j^{\rm res}.
 \label{eq:pdfv2-sol-shell-p}
\end{align}

The three mutually orthogonal projectors are defined.
\begin{align}
 P_j^{\rm term}
 &:=P_j^{\rm depth}\otimes P_j^{\rm res},
 \label{eq:pdfv2-sol-pterm}\\
 Q_j^{\rm micro}
 &:=(I-P_j^{\rm depth})\otimes I,
 \label{eq:pdfv2-sol-qmicro}\\
 Q_j^{\rm shell}
 &:=P_j^{\rm depth}\otimes(I-P_j^{\rm res}).
 \label{eq:pdfv2-sol-qshell}
\end{align}
Then
\begin{equation}\label{eq:pdfv2-sol-ortogonalidad}
 Q_j^{\rm micro}Q_j^{\rm shell}=0,
 \qquad
 I=P_j^{\rm term}+Q_j^{\rm micro}+Q_j^{\rm shell},
\end{equation}
and, for all \(\xi\in\mathscr H_j\),
\begin{equation}\label{eq:pdfv2-sol-pitagoras}
 \|\xi\|^2
 =\|P_j^{\rm term}\xi\|^2
  +\|Q_j^{\rm micro}\xi\|^2
  +\|Q_j^{\rm shell}\xi\|^2.
\end{equation}
The term
\((I-P_j^{\rm depth})\otimes(I-P_j^{\rm res})\) does not appear a second time: it is already
contained exactly once in \(Q_j^{\rm micro}\).

\section{Cylinder extensions and full row of losses}
\label{sec:pdfv2-sol-gamma-loss}

An extension of cylinders from level \(j\) to level \(j+1\) induces the map
\begin{equation}\label{eq:pdfv2-sol-gamma}
 \Gamma_j:\mathscr H_j\longrightarrow\mathscr H_{j+1},
 \qquad
 \Gamma_j f=f\circ\tau_{j+1,j},
\end{equation}
with all leaf, orientation, carry, memory, and path labels
transported in the cylinder basis.  For \(n>j\), we write
\begin{equation}\label{eq:pdfv2-sol-gamma-composition}
 \Gamma_{j:n}:=\Gamma_{n-1}\cdots\Gamma_j,
 \qquad
 \Gamma_{j:j}:=I.
\end{equation}

On \(\mathscr H_j\), the nonnegative increments define multiplication
operators
\begin{equation}\label{eq:pdfv2-sol-loss-multipliers}
 B_j^+:=M_{\sqrt{b_j^+}},
 \qquad
 B_j^-:=M_{\sqrt{b_j^-}},
 \qquad
 K_j:=M_{\sqrt{2\kappa_j}}.
\end{equation}
Each root is taken pointwise over the finite set of cylinders; if a
coordinate is zero, its operator is zero.  The space of losses is the external
orthogonal sum
\[
 \mathscr E_j
 :=\mathscr H_j^{(+)}\oplus\mathscr H_j^{(-)}
   \oplus\mathscr H_j^{(\kappa)}
   \oplus\mathscr H_j^{({\rm micro})}
   \oplus\mathscr H_j^{({\rm shell})},
\]
and the complete row is
\begin{equation}\label{eq:pdfv2-sol-loss-row}
 L_j:\mathscr H_j\longrightarrow\mathscr E_j,
 \qquad
 L_j\xi
 =\bigl(
 B_j^+\xi,B_j^-\xi,K_j\xi,
 Q_j^{\rm micro}\xi,Q_j^{\rm shell}\xi
 \bigr).
\end{equation}
Therefore,
\begin{equation}\label{eq:pdfv2-sol-loss-norm}
 \|L_j\xi\|^2
 =\|B_j^+\xi\|^2+\|B_j^-\xi\|^2+\|K_j\xi\|^2
  +\|Q_j^{\rm micro}\xi\|^2
  +\|Q_j^{\rm shell}\xi\|^2.
\end{equation}
The positive part, negative part, cancellation carry, microstructure, and shell
each appear once. They are neither merged into an anonymous constant nor
repeated in two summands.

\section{Complete history terminal}
\label{sec:pdfv2-sol-terminal}

Let \(\operatorname{Hist}_{J}^{\rm sol}\) be the set of histories
surviving to the horizon \(J\), each with its package of thirteen
coordinates.  The terminal space is defined as the free space on these
complete records:
\begin{equation}\label{eq:pdfv2-sol-terminal-space}
 \mathscr T_J^{\rm sol}
 :=\ell^2\!\left(
 \left\{(h,D_{{\rm sol},J}(h)):
 h\in\operatorname{Hist}_{J}^{\rm sol}\right\}
 \right).
\end{equation}
The terminal evaluation
\begin{equation}\label{eq:pdfv2-sol-terminal-evaluation}
 E_J^{\rm term}:\mathscr H_J\longrightarrow\mathscr T_J^{\rm sol}
\end{equation}
sends each basic cylinder vector to the complete record of the history that
generates it and extends linearly.  In particular, the terminal retains
fiber, relations, numbers, orientation, carry, memory, boundary, route,
multisection, survival, and reader.  It is not replaced by the last visible
value of the route.

\section{Observation column, Gram operator, and telescoping}
\label{sec:pdfv2-sol-gram}

For \(0\leq j\leq J\), the complete future column is
\begin{equation}\label{eq:pdfv2-sol-g-column}
 \begin{aligned}
 G_{j;J}:\mathscr H_j&\longrightarrow
 \mathscr T_J^{\rm sol}\oplus
 \bigoplus_{n=j}^{J-1}\mathscr E_n,\\
 G_{j;J}\xi
 &:={}
 E_J^{\rm term}\Gamma_{j:J}\xi
 \ \oplus\ 
 \bigoplus_{n=j}^{J-1}L_n\Gamma_{j:n}\xi.
 \end{aligned}
\end{equation}
For \(j=J\), the sum of perdidas is empty and
\(G_{J;J}=E_J^{\rm term}\).  The Gram terminal is
\begin{equation}\label{eq:pdfv2-sol-u-gram}
 U_{j;J}:=G_{j;J}^*G_{j;J}\succeq0.
\end{equation}

Rearranging the orthogonal sum of the codomain yields the literal identity.
\begin{equation}\label{eq:pdfv2-sol-g-recursion}
 G_{j;J}\xi
 \simeq
 \bigl(L_j\xi,\,G_{j+1;J}\Gamma_j\xi\bigr).
\end{equation}
Consequently,
\begin{equation}\label{eq:pdfv2-sol-stein}
 \boxed{U_{j;J}
 =L_j^*L_j+\Gamma_j^*U_{j+1;J}\Gamma_j}
\end{equation}
and
\begin{equation}\label{eq:pdfv2-sol-stein-difference}
 U_{j;J}-\Gamma_j^*U_{j+1;J}\Gamma_j
 =L_j^*L_j\succeq0.
\end{equation}
There is no sign mutation: the loss term is on the positive side
of the identity because it arises from a sum of squares.

Iterating \eqref{eq:pdfv2-sol-stein} yields the exact telescoping identity
\begin{equation}\label{eq:pdfv2-sol-telescopia-operatorial}
 U_{j;J}
 =\Gamma_{j:J}^*(E_J^{\rm term})^*E_J^{\rm term}\Gamma_{j:J}
  +\sum_{n=j}^{J-1}
   \Gamma_{j:n}^*L_n^*L_n\Gamma_{j:n},
\end{equation}
o, in quadratic form,
\begin{equation}\label{eq:pdfv2-sol-telescopia-normas}
 \langle U_{j;J}\xi,\xi\rangle
 =\|E_J^{\rm term}\Gamma_{j:J}\xi\|^2
  +\sum_{n=j}^{J-1}\|L_n\Gamma_{j:n}\xi\|^2.
\end{equation}
The first term preserves the complete terminal future; the sum preserves,
level by level, all the losses and all the carries that led to it.

\iffalse
\section{The thirteen coordinates of the constraint}
\label{sec:pdfv2-sol-trece}

For a prefix \(x_{\leq j}\), it is defined
\begin{equation}\label{eq:pdfv2-sol-d-package}
 \begin{aligned}
 D_{{\rm sol},j}(x):=\bigl(&
 \mathcal F_{{\rm sol},j},
 \operatorname{Ext}_{\Gamma,{\rm sol},j},
 \operatorname{Rel}_{{\rm sol},j},
 \mathcal N_{{\rm sol},j},
 \operatorname{or}_{{\rm sol},j},
 \operatorname{Carr}_{{\rm sol},j},
 \operatorname{Mem}_{{\rm sol},j},\\
 &\partial_{{\rm sol},j},
 \gamma_{{\rm sol},j},
 \operatorname{Msect}_{{\rm sol},j},
 \operatorname{Surv}_{{\rm sol},j},
 \operatorname{Term}_{{\rm sol},j},
 \mathcal L_{{\rm sol},j}\bigr).
 \end{aligned}
\end{equation}
Its components are as follows.
\begin{enumerate}
 \item \(\mathcal F_{{\rm sol},j}\) contains the cylinders \(Y_j\), their projective
       measure, \(\mathscr H_j^{\rm depth}\),
       \(\mathscr H_j^{\rm res}\), their tensor products, and the five
       spaces of losses.
 \item \(\operatorname{Ext}_{\Gamma,{\rm sol},j}\) contains
       \(J_{9,j},E_{9,j},\Gamma_j,\Gamma_{j:n}\), the shell inclusions, and
       all the truncation and refinement maps.
 \item \(\operatorname{Rel}_{{\rm sol},j}\) contains
       \eqref{eq:pdfv2-sol-bpm-identidades},
       \eqref{eq:pdfv2-sol-ej-projector},
       \eqref{eq:pdfv2-sol-jensen-plus}--
       \eqref{eq:pdfv2-sol-jensen-minus},
       \eqref{eq:pdfv2-sol-ortogonalidad} and
       \eqref{eq:pdfv2-sol-stein}.
 \item \(\mathcal N_{{\rm sol},j}\) conserves
       \(j,J,9,b_j,b_j^+,b_j^-,\kappa_j\), sheet multiplicities, fiber sizes,
       and shell dimensions.
 \item \(\operatorname{or}_{{\rm sol},j}\) preserves
       \(\Delta_{\gamma_j}\), the path order, and the inversion action
       \(b_j\mapsto-b_j\), which interchanges \(b_j^+\) and \(b_j^-\) without changing
       \(|b_j|\) or \(\kappa_j\).
 \item \(\operatorname{Carr}_{{\rm sol},j}\) preserves the integer TPK carry, the composition carry and the cancellation carry \(\kappa_j\) as distinct coordinates.
 \item \(\operatorname{Mem}_{{\rm sol},j}\) is
       \((M_j^+,M_j^-,D_j^0,H_j^0)\) together with the ordered ledger of all
       preceding increments.
 \item \(\partial_{{\rm sol},j}\) contains the oriented cylinder boundary,
       the increments \(b_j\), the row \(L_j\), and the Stein differences.
 \item \(\gamma_{{\rm sol},j}\) is the complete route \(e_0e_1\cdots e_{j-1}\), not merely
its endpoints.
 \item \(\operatorname{Msect}_{{\rm sol},j}\) contains all compatible
       nonadic and shell extensions; it does not choose a prolongation.
 \item \(\operatorname{Surv}_{{\rm sol},j}\) records the horizons
       \(J\geq j\) for which \(G_{j;J}\), \(U_{j;J}\), and
       compatible terminal objects exist.
 \item \(\operatorname{Term}_{{\rm sol},j}\) contains
       \(E_J^{\rm term}\Gamma_{j:J}\), the complete record of history and
       the trece coordinates in each horizon.
 \item \(\mathcal L_{{\rm sol},j}\) internally publishes
       \((b^\pm,M^\pm,D^0,H^0,\kappa,Q^{\rm micro},Q^{\rm shell},
       L,G,U)\) and its identities.
\end{enumerate}

For an extension \(u:x_j\to x_{j'}\), the coordinated transport
\(u_{{\rm sol},j',j}\) satisfies
\begin{equation}\label{eq:pdfv2-sol-naturality-d}
 u_{{\rm sol},j',j}D_{{\rm sol},j'}(x_{j'})
 =D_{{\rm sol},j}(\tau_{j',j}x_{j'}).
\end{equation}
The equality includes all five components of \(L\), the complete terminal, and
the composition of \(\Gamma\); it is not an equality of cardinalities alone.

\section{Constraint graph and dependent object}
\label{sec:pdfv2-sol-graphs}

We define
\begin{equation}\label{eq:pdfv2-sol-g-graph}
 \mathcal G_{\rm sol}:=\operatorname{Graph}(D_{\rm sol}),
 \qquad
 \operatorname{Res}_{\rm sol}(x):=(x,D_{\rm sol}x).
\end{equation}
The first projection is inverse over the graph:
\begin{equation}\label{eq:pdfv2-sol-res-inverse}
 \pi_1\operatorname{Res}_{\rm sol}=I,
 \qquad
 \operatorname{Res}_{\rm sol}\pi_1=I_{\mathcal G_{\rm sol}}.
\end{equation}

Let \(\chi_{\rm sol}(x)\) be the complete multisection formed by all the
cylinders, refinements, shells, rows of losses, and terminals compatible
with \(x\).  Then
\begin{equation}\label{eq:pdfv2-sol-xchi}
 X_{\chi,{\rm sol}}
 :=\{(\chi_{\rm sol}(x),x,D_{\rm sol}x):
 x\in\mathcal C_{\rm cont}^{\rm sol}\},
\end{equation}
and
\begin{equation}\label{eq:pdfv2-sol-prx}
 \operatorname{pr}_{X,{\rm sol}}(x,D_{\rm sol}x)
 :=(\chi_{\rm sol}(x),x,D_{\rm sol}x).
\end{equation}
Again, the projection that preserves \((x,D_{\rm sol}x)\) is inverse on
this dependent image.

\section{Internal assembler and publisher}
\label{sec:pdfv2-sol-assembler-publisher}

The ambient space \(\mathscr M_{\rm sol}^{\rm amb}\) consists of compatible
systems
\[
 \bigl(
 \{b,b^\pm,M^\pm,D^0,H^0,\kappa\},
 \{E_9,J_9,P^{\rm term},Q^{\rm micro},Q^{\rm shell}\},
 \{\Gamma,L,E^{\rm term},G,U\}
 \bigr)
\]
with all the previous identities. The raw assembler is
\begin{equation}\label{eq:pdfv2-sol-assembler}
 \begin{aligned}
 a_{\rm sol}:X_{\chi,{\rm sol}}&\longrightarrow
 \mathscr M_{\rm sol}^{\rm amb},\\
 a_{\rm sol}(\chi(x),x,Dx)&:=
 \bigl(
 \{b,b^\pm,M^\pm,D^0,H^0,\kappa\}_x,
 \{E_9,J_9,P,Q\}_x,
 \{\Gamma,L,E^{\rm term},G,U\}_x
 \bigr).
 \end{aligned}
\end{equation}
The macrostructure preserves its antecedent through
\begin{equation}\label{eq:pdfv2-sol-macro-graph}
 \mathfrak M_{\rm sol}:=\operatorname{Graph}(a_{\rm sol}),
 \qquad
 \mathcal A_{\rm sol}(z):=(z,a_{\rm sol}z).
\end{equation}

Let \(\mathscr Y_{\rm sol}^{\rm int}\) be the type of certificates formed by
the identities
\begin{equation}\label{eq:pdfv2-sol-certificate-list}
 \begin{gathered}
 b=b^+-b^-,\qquad |b|=b^++b^-,\\
 E_9J_9=I,\qquad J_9E_9=P_9,\\
 E_9b_f^\pm=b_c^\pm+\kappa,\qquad\kappa\geq0,\\
 I=P^{\rm term}+Q^{\rm micro}+Q^{\rm shell},
 \qquad Q^{\rm micro}Q^{\rm shell}=0,\\
 U_{j;J}=L_j^*L_j+\Gamma_j^*U_{j+1;J}\Gamma_j,\\
 \langle U_{j;J}\xi,\xi\rangle
 =\|E_J^{\rm term}\Gamma_{j:J}\xi\|^2
  +\sum_{n=j}^{J-1}\|L_n\Gamma_{j:n}\xi\|^2,
 \end{gathered}
\end{equation}
together with the naturality of all arrows. The raw publisher
\begin{equation}\label{eq:pdfv2-sol-publisher}
 p_{\rm sol}^{\rm raw}:\mathfrak M_{\rm sol}
 \longrightarrow\mathscr Y_{\rm sol}^{\rm int}
\end{equation}
publish that certificate without removing the macroobject. Finally,
\begin{equation}\label{eq:pdfv2-sol-output-graph}
 \mathcal C_{\rm sol}^{\rm HMT}
 :=\operatorname{Graph}(p_{\rm sol}^{\rm raw}),
 \qquad
 \mathcal P_{\rm sol}(m):=(m,p_{\rm sol}^{\rm raw}m).
\end{equation}

\section{Solenoidal Generation Internal Theorem}
\label{sec:pdfv2-sol-theorem}

\begin{theorem}[Internal generation of the balance macrostructure]
\label{thm:pdfv2-sol-internal-generation}
For every \(x\in\mathcal C_{\rm cont}^{\rm sol}\), the chain
\[
 \mathcal C_{\rm cont}^{\rm sol}
 \xrightarrow{\operatorname{Res}_{\rm sol}}
 \mathcal G_{\rm sol}
 \xrightarrow{\operatorname{pr}_{X,{\rm sol}}}
 X_{\chi,{\rm sol}}
 \xrightarrow{\mathcal A_{\rm sol}}
 \mathfrak M_{\rm sol}
 \xrightarrow{\mathcal P_{\rm sol}}
 \mathcal C_{\rm sol}^{\rm HMT}
\]
naturally constructs a macrostructure containing both memories, the cancellation carry, the orthogonal micro/shell separation, the complete row of losses, the future column, the full-history terminal object, the Gram matrix, and the telescoping structure. The thirteen coordinates of the continuum can be recovered by successive projections from any point in the chain.
\end{theorem}

\begin{proof}
The identities \eqref{eq:pdfv2-sol-bpm-identidades} and
\eqref{eq:pdfv2-sol-dh-diferencias} are obtained directly from the
positive and negative decomposition of the integer increment; therefore the
memories are not added data.  The formulas
\eqref{eq:pdfv2-sol-j9}--\eqref{eq:pdfv2-sol-e9}, together with the projective
measure, prove \(E_9J_9=I\), \(J_9^*=E_9\), and that \(P_9\) is an
orthogonal projector.

The triangle inequality gives
\(E_9|b_f|\geq|E_9b_f|=|b_c|\), so that
\(\kappa\geq0\).  Substituting
\(b^\pm=(|b|\pm b)/2\) proves
\eqref{eq:pdfv2-sol-jensen-plus} and
\eqref{eq:pdfv2-sol-jensen-minus} without an additional hypothesis.

The tensorial definitions
\eqref{eq:pdfv2-sol-pterm}--\eqref{eq:pdfv2-sol-qshell} prove that
\(P^{\rm term}\), \(Q^{\rm micro}\), and \(Q^{\rm shell}\) are orthogonal and
sum to the identity. Consequently, row \(L_j\) counts each sector only
once.

The definition \eqref{eq:pdfv2-sol-g-column} orthogonally separates the terminal
and the losses of each level.  Separating the first summand of losses produces
\eqref{eq:pdfv2-sol-g-recursion}.  Taking the adjoint and composing proves the
Stein identity \eqref{eq:pdfv2-sol-stein}; iterating it proves
\eqref{eq:pdfv2-sol-telescopia-operatorial} and
\eqref{eq:pdfv2-sol-telescopia-normas}.  Every term is a squared norm,
so the opposite sign cannot appear.

Compatibility of cylinder truncation, expectation, projections, and
extensions \(\Gamma\) yields the naturality
\eqref{eq:pdfv2-sol-naturality-d}. Finally,
\(\mathcal G_{\rm sol}\), \(\mathfrak M_{\rm sol}\), and
\(\mathcal C_{\rm sol}^{\rm HMT}\) are graphs, and
\(X_{\chi,{\rm sol}}\) is a dependent object that explicitly preserves
its antecedent. The first projection at each stage recovers the preceding
stage. Thus no composition erases history, memory, multisection,
terminal, or reader.
\end{proof}

\section{Probative origin and forgers}
\label{sec:pdfv2-sol-provenance-falsifiers}

\paragraph{Result recovered.}
The signed increment with carry, the separation \(b=b^+-b^-\), the memories
\(M^\pm,D^0,H^0\), the pair \(E_9,J_9\), the Jensen carry \(\kappa\), cylinder
extension, the loss row, the terminal, the Gram matrix, Stein, and telescoping
belong to the already constructed architecture.

\paragraph{New Formalization.}
Explicit formalizations of that architecture are the typing prior to every
external evaluation, the factorization
\(\mathscr H^{\rm depth}\widehat\otimes\mathscr H^{\rm res}\), the orthogonal
separation preventing double counting, the free terminal on histories with
thirteen coordinates, and the four nonablative graph packagings.

\begin{criterion}[Internal falsifiers of the branch \(\mathrm{sol}\)]
\label{crit:pdfv2-sol-falsifiers}
The construction fails if any of the following occurs:
\begin{enumerate}
 \item \(b\neq b^+-b^-\) o \(|b|\neq b^++b^-\);
 \item \(E_9J_9\neq I\), the measure is not projective or \(P_9\) is not a
       projector;
 \item \(\kappa<0\) or fail
       \(E_9b_f^\pm=b_c^\pm+\kappa\);
 \item \(Q^{\rm micro}Q^{\rm shell}\neq0\), or the same sector appears in
       both summands of row \(L\);
 \item one of \(b^+,b^-,\kappa,Q^{\rm micro},Q^{\rm shell}\) is omitted or
       counted twice;
 \item the terminal preserves only the visible value and not the
history with its thirteen coordinates;
 \item the Stein identity fails or the telescoping contains a sign
       contrary to a sum of squares;
 \item an extension does not commute with \(E_9,J_9,L,G,U\);
 \item the multisection of continuations is replaced with a retrospectively chosen prolongation;
 \item any of the thirteen coordinates is omitted.
\end{enumerate}
In the last case, the binding diagnosis is: \emph{object replaced;
conclusion not transferable}.
\end{criterion}

\begin{warningbox}
This chapter concludes only the internal construction of
\(\mathcal C_{\rm sol}^{\rm HMT}\). It contains no classical differential
equation, classical initial data, map from external data, or
conclusion concerning an external problem. Those intertwiners belong to another
document and do not govern the existence of this macrostructure.
\end{warningbox}
\fi

\endgroup

\section{Orientation, carry, and simultaneous composition of paths}

Path inversion acts only once on
\(\mathfrak e_\alpha\). Let
\(\mathsf J(v_\Sigma,v_\Pi):=(v_\Pi,v_\Sigma)\) be the leaf exchange,
and let \(P_{\rm tr}\) be the transposition of port matrices. Its five
local expressions are
\begin{equation}
\label{eq:pdfv2-cor-orientacion-cruzada}
 \begin{gathered}
 \mathsf J\sigma_k=-\sigma_k\mathsf J,
 \qquad b_{\bar\alpha}=-b_\alpha,
 \qquad (b_{\bar\alpha}^+,b_{\bar\alpha}^-)
       =(b_\alpha^-,b_\alpha^+),\\
 \vartheta_L(K,\epsilon):=
 \begin{cases}
  (L,-\epsilon),&K=L,\\
  (K,\epsilon),&K\neq L,
 \end{cases}
 \qquad
 \kappa_{J,I}(v,u)=-P_{\rm tr}\kappa_{I,J}(u,v),\\
 \kappa(\bar\alpha)
 =-\mathsf C_\alpha^{-1}\kappa(\alpha),
 \qquad v_{\bar\alpha}=-v_\alpha.
 \end{gathered}
\end{equation}
Here \(\mathsf J\) exchanges the two sheets, \(P_{\rm tr}\) transposes the incidence matrix, and \(\vartheta_L\) exchanges exactly the two faces in the direction \(L\), fixing the other six. In the free basis of \(\mathscr B=\mathscr I\times\{+,-\}\), the eight vectors are distinct, \(\vartheta_L^2=I\) and \(\vartheta_L\neq\vartheta_{L'}\) for \(L\neq L'\). The same orientation explains all signs in \eqref{eq:pdfv2-cor-orientacion-cruzada}.

For composable extensions \(\alpha:x\to y\) and \(\beta:y\to z\), the
native carry satisfies
\begin{equation}
\label{eq:pdfv2-cor-carry-nativo}
 \kappa(\beta\alpha)=
 \kappa(\beta)+\mathsf C(\beta)\kappa(\alpha).
\end{equation}
The matrix lift is not left as a name. In the ordered basis
\((f,r,e,b,g)\), let
\(\widetilde L_N(\alpha)\in\operatorname{Mat}_5(\mathbb Z)\) be the
centered integral representative of the matrix of
\(\Psi_\alpha\) of \eqref{eq:pdfv2-cor-psi-generadores} modulo \(3^N\), and
define
\begin{equation}
\label{eq:pdfv2-cor-carry-matricial-def}
 c_N(\beta,\alpha):=
 \frac{\widetilde L_N(\beta)\widetilde L_N(\alpha)
       -\widetilde L_N(\beta\alpha)}{3^N}
 \in\operatorname{Mat}_5(\mathbb Z).
\end{equation}
Divisibility follows from equality of the two compositions modulo
\(3^N\). Writing \(L_N=[\widetilde L_N]_{3^N}\), the integer lifts
satisfy
\begin{equation}
\label{eq:pdfv2-cor-carry-matricial}
 \widetilde L_N(\beta)\widetilde L_N(\alpha)
 =\widetilde L_N(\beta\alpha)+3^Nc_N(\beta,\alpha)
\end{equation}
and associativity produces the cocycle
\begin{equation}
\label{eq:pdfv2-cor-cociclo-matricial}
 c_N(\delta,\beta\alpha)+\widetilde L_N(\delta)c_N(\beta,\alpha)
 =c_N(\delta\beta,\alpha)
  +c_N(\delta,\beta)\widetilde L_N(\alpha).
\end{equation}

The same carry increment acts upon the local complex through
\begin{equation}
\label{eq:pdfv2-cor-psi-generadores}
 \begin{gathered}
 \Psi_\alpha(r_x)=r_y,\quad
 \Psi_\alpha(e_x)=e_y,\quad
 \Psi_\alpha(g_x)=g_y,\quad
 \Psi_\alpha(f_x)=f_y,\\
 \Psi_\alpha(b_x)=b_y+3(c_y-c_x)e_y.
 \end{gathered}
\end{equation}
The oriented boundary change remains in
\begin{equation}
\label{eq:pdfv2-cor-k-alpha}
 K_\alpha:=(v_y-v_x)f_y,
 \qquad
 K_{\beta\alpha}=K_\beta+\Psi_\beta K_\alpha.
\end{equation}
The equations \eqref{eq:pdfv2-cor-carry-nativo},
\eqref{eq:pdfv2-cor-cociclo-matricial}, and
\eqref{eq:pdfv2-cor-k-alpha} are three degrees of the same ledger of composition.
None authorizes erasing the other two.

\subsection{Nerve of extensions and Alexander--Whitney diagonal}
For a history of the common ledger, write
\(A_n:=\mathbf Z/3^n\mathbf Z\) and \(\mathsf T_m(x)\) for the finite category
of \emph{all} its TPK prefixes of depth \(m\), with the
extension morphisms already constructed. The following diagonal and counit
are applied to that same nerve; they do not define a fifth autonomous branch.
\begingroup
  \let\section\subsection
  % Extracción quirúrgica del propietario V2 05e: líneas 68--110.
% El resto permanece en este archivo como procedencia, pero no se publica.
\iffalse
\chapter[Internal Determinantal Macrostructure]{Internal Determinantal Macrostructure Generated by the Discrete Continuum}
\label{chap:detint-macroestructura}

\begin{thesisbox}
The \({\rm det}\) branch starts exclusively from surviving histories of
\(\mathcal C_{\rm cont}^{\rm disc}\). The nerve of extensions, the Alexander--Whitney diagonal,
the local complex, the unit, the two publications, the filler, the terminal
reduction, and the \(3\)-adic tower are constructed before any subsequent
evaluation. The balance and acyclicity conditions are incorporated into the
typed domain and are never tacitly assumed.
\end{thesisbox}

\section{Surviving determinant domain}

Let
\begin{equation}\label{eq:detint-an}
 A_n:=\mathbf Z/3^n\mathbf Z,
 \qquad
 \rho_{n+1,n}:A_{n+1}\longrightarrow A_n.
\end{equation}
For \(x\in\mathcal C_{\rm cont}^{\rm disc}\), let
\(\mathsf T_m(x)\) be the finite category of TPK prefixes of depth \(m\)
compatible with \(x\). Each object preserves sheet, orientation, residue,
quotient, integer carry, memory, boundary, route, multisection, survival,
and terminal; its morphisms are the admissible extensions that preserve those
coordinates.

From each object \(y\in\mathsf T_m(x)\), the internal integers are extracted
\begin{equation}\label{eq:detint-c-v-lifts}
 \widetilde c_y,\widetilde v_y\in\mathbf Z,
 \qquad
 c_{y,n}:=\widetilde c_y\bmod3^n,
 \qquad
 v_{y,n}:=\widetilde v_y\bmod3^n.
\end{equation}
The first coordinate is the unoriented carry and the second, the oriented
boundary coefficient. Reversal satisfies
\begin{equation}\label{eq:detint-or-cv}
 c_{\bar y,n}=c_{y,n},
 \qquad
 v_{\bar y,n}=-v_{y,n}.
\end{equation}

We write
\begin{equation}\label{eq:detint-z3-k3}
 \mathbf Z_3:=\varprojlim_n A_n,
 \qquad
 K_3:=\operatorname{Frac}(\mathbf Z_3).
\end{equation}
The determinantial restriction is defined by
\begin{equation}\label{eq:detint-dominio}
\begin{split}
 \mathcal C_{\rm cont}^{\rm det}:=
 \bigl\{x\in\mathcal C_{\rm cont}^{\rm disc}:{}&
 x\text{ survives at every depth};\\
 &\widehat P_x^{\rm red}\text{ is finite and based};\\
 &\chi(\widehat P_x^{\rm red})=0;\\
 &H_*(\widehat P_x^{\rm red}\otimes_{\mathbf Z_3}K_3)=0;\\
 &\text{the canonical reduction produces a square block}\\
 &\text{stable under refinement}\bigr\}.
\end{split}
\end{equation}
The objects \(\widehat P_x^{\rm red}\) and the square block will be constructed
in \cref{sec:detint-terminal}. Therefore
\eqref{eq:detint-dominio} is a verifiable membership condition and not
a conclusion inserted into the definition of the maps.

\fi
\section{Nerve of routes, Alexander--Whitney diagonal and counit}

The normalized complex is defined
\begin{equation}\label{eq:detint-gmn}
 G_{m,n}(x):=
 N_*^{\rm norm}\!\left(N\mathsf T_m(x);A_n\right).
\end{equation}
For a nondegenerate simplex
\(\sigma=[y_0\to y_1\to\cdots\to y_q]\), the diagonal is
\begin{equation}\label{eq:detint-aw}
 \Delta_{\rm AW}(\sigma)
 =\sum_{j=0}^q
 [y_0\to\cdots\to y_j]\otimes
 [y_j\to\cdots\to y_q].
\end{equation}
The counit is given by
\begin{equation}\label{eq:detint-counidad-g}
 \epsilon_G([y])=1,
 \qquad
 \epsilon_G(\sigma)=0\quad(q>0).
\end{equation}
In particular,
\begin{equation}\label{eq:detint-vertice-grouplike}
 \Delta_{\rm AW}[y]=[y]\otimes[y],
 \qquad
 \epsilon_G([y])=1.
\end{equation}

If \(m'\ge m\) and \(n'\ge n\), truncation and coefficient reduction
induce
\begin{equation}\label{eq:detint-rmn}
 R_{m',m}^{n',n}:G_{m',n'}(x)\longrightarrow G_{m,n}(x),
\end{equation}
and they are verified
\begin{equation}\label{eq:detint-aw-natural}
 R\partial=\partial R,
 \qquad
 (R\otimes R)\Delta_{\rm AW}=\Delta_{\rm AW}R,
 \qquad
 \epsilon_GR=\rho_{n',n}\epsilon_G.
\end{equation}
The path is not reduced to its endpoints: the entire simplex remains in
\(G_{m,n}(x)\).

\iffalse
\section{Internal local complex}

For \(c\in A_n\), define
\begin{equation}\label{eq:detint-p0-grados}
 P^0_{n,1}(c)=A_nf,
 \qquad
 P^0_{n,0}(c)=A_nr\oplus A_ne\oplus A_nb,
 \qquad
 P^0_{n,-1}(c)=A_ng,
\end{equation}
with differential
\begin{equation}\label{eq:detint-p0-diferencial}
 \partial f=3ce+b,
 \qquad
 \partial r=0,
 \qquad
 \partial e=g,
 \qquad
 \partial b=-3cg.
\end{equation}
The chain identity is literal:
\begin{equation}\label{eq:detint-d2}
 \partial^2f=3c\,\partial e+\partial b=3cg-3cg=0,
 \qquad
 \partial^2e=\partial^2b=\partial^2r=0.
\end{equation}

The increase
\begin{equation}\label{eq:detint-epsilon-p}
 \epsilon_P:P_n^0(c)\longrightarrow A_n[0]
\end{equation}
is fixed by
\[
 \epsilon_P(r)=1,
 \qquad
 \epsilon_P(e)=\epsilon_P(b)=\epsilon_P(f)=\epsilon_P(g)=0.
\]
The complex possesses the differential coalgebra.
\begin{equation}\label{eq:detint-coalgebra-p}
 \Delta_P(r)=r\otimes r,
 \qquad
 \Delta_P(q)=q\otimes r+r\otimes q
 \quad(q\in\{e,b,f,g\}).
\end{equation}
Thus, \(r\) is group-like and the other four generators are primitive
relative to \(r\). Equations
\eqref{eq:detint-p0-diferencial} prove that \(\Delta_P\) commutes with the
tensor differential.

\section{Local extension system and memory ledger}

For an extension \(\alpha:y\to y'\), we define
\begin{equation}\label{eq:detint-psi}
 \Psi_\alpha:P_n^0(c_y)\longrightarrow P_n^0(c_{y'})
\end{equation}
through means of
\begin{equation}\label{eq:detint-psi-generadores}
\begin{gathered}
 \Psi_\alpha(r_y)=r_{y'},\quad
 \Psi_\alpha(e_y)=e_{y'},\quad
 \Psi_\alpha(g_y)=g_{y'},\quad
 \Psi_\alpha(f_y)=f_{y'},\\
 \Psi_\alpha(b_y)=b_{y'}+3(c_{y'}-c_y)e_{y'}.
\end{gathered}
\end{equation}
It is a map of complexes because
\begin{equation}\label{eq:detint-psi-b-chain}
 \partial\Psi_\alpha(b_y)
 =-3c_{y'}g+3(c_{y'}-c_y)g
 =-3c_yg
 =\Psi_\alpha(\partial b_y),
\end{equation}
and
\begin{equation}\label{eq:detint-psi-f-chain}
 \Psi_\alpha(\partial f_y)
 =3c_ye+b+3(c_{y'}-c_y)e
 =3c_{y'}e+b
 =\partial\Psi_\alpha(f_y).
\end{equation}
Moreover,
\begin{equation}\label{eq:detint-psi-functorial}
 \Psi_\beta\Psi_\alpha=\Psi_{\beta\alpha}.
\end{equation}

The oriented coordinate change does not erase. It is recorded by
\begin{equation}\label{eq:detint-k-alpha}
 K_\alpha:=(v_{y'}-v_y)f_{y'}.
\end{equation}
For composable extensions,
\begin{equation}\label{eq:detint-k-coherencia}
 K_{\beta\alpha}=K_\beta+\Psi_\beta K_\alpha.
\end{equation}
This equality is the exact ledger of the boundary increment and makes visible
the memory that a purely nominal composition would lose.

The complete local diagram is the Grothendieck construction.
\begin{equation}\label{eq:detint-grothendieck-local}
 \int_{y\in\mathsf T_m(x)}P_n^0(c_y).
\end{equation}
In it, an arrow over \(\alpha:y\to y'\) carries \(\Psi_\alpha\)
as its linear component and \(K_\alpha\) as its coherence. When applied to an
element supported at a vertex, the notation is used
\begin{equation}\label{eq:detint-psi-soporte-vertice}
 \Psi_\alpha(p,[y]):=(\Psi_\alpha p,[y']).
\end{equation}
The edge \([y\to y']\) remains in the path coordinate of the nerve. This
convention will be used in the naturality formulas for
\(z_\pm\) and \(h_*\); no global map erasing the remaining
simplices of \(G_{m,n}(x)\) is asserted.

\section{Pullback, unit, root, publications and filler}

The pullback of complexes is
\begin{equation}\label{eq:detint-pgen}
 P^{\rm gen}_{y;m,n}
 :=P_n^0(c_y)\times_{A_n[0]}G_{m,n}(x)
 =\{(p,\gamma):\epsilon_P(p)=\epsilon_G(\gamma)\}.
\end{equation}
Its differential is
\begin{equation}\label{eq:detint-pgen-d}
 \partial(p,\gamma)=(\partial p,\partial\gamma).
\end{equation}
The common unit associated with the vertex \(y\) is
\begin{equation}\label{eq:detint-unidad}
 \mathbf1_y:=(r,[y]).
\end{equation}
Its two projections are group-like by
\eqref{eq:detint-vertice-grouplike} and
\eqref{eq:detint-coalgebra-p}; that is the precise meaning of a compatible
group-like unit of the pullback.

The root projection is the map of complexes
\begin{equation}\label{eq:detint-proyeccion-root}
 \varpi_{\rm root}:P^{\rm gen}_{y;m,n}\longrightarrow
 \operatorname{Root}_{y;m,n},
 \qquad
 \varpi_{\rm root}(p,\gamma):=(\epsilon_P(p)r,\gamma).
\end{equation}
The two internal publications are defined.
\begin{equation}\label{eq:detint-zpm}
 z_-(y):=(r,[y]),
 \qquad
 z_+(y):=(r+3c_yv_ye+v_yb,[y]).
\end{equation}
Both are cycles:
\begin{equation}\label{eq:detint-zpm-ciclos}
 \partial z_-(y)=0,
 \qquad
 \partial z_+(y)=3c_yv_yg-3c_yv_yg=0.
\end{equation}
Their roots coincide:
\begin{equation}\label{eq:detint-raiz-comun}
 \varpi_{\rm root}z_-(y)
 =(r,[y])
 =\varpi_{\rm root}z_+(y).
\end{equation}

The filler is constructed directly:
\begin{equation}\label{eq:detint-hstar}
 h_*(y):=(-v_yf,0)\in P^{\rm gen}_{y;m,n,1}.
\end{equation}
Then
\begin{equation}\label{eq:detint-filler-identidad}
 \partial h_*(y)
 =(-3c_yv_ye-v_yb,0)
 =z_-(y)-z_+(y).
\end{equation}
In particular,
\begin{equation}\label{eq:detint-clases-iguales}
 [z_-(y)]=[z_+(y)]
 \quad\text{in}\quad H_0(P^{\rm gen}_{y;m,n}),
\end{equation}
and the equality preserves its witness chain.

Under an extension \(\alpha:y\to y'\), we have
\begin{equation}\label{eq:detint-z-natural}
 z_-(y')=\Psi_\alpha z_-(y),
 \qquad
 z_+(y')-\Psi_\alpha z_+(y)=\partial K_\alpha,
\end{equation}
and
\begin{equation}\label{eq:detint-h-natural}
 h_*(y')-\Psi_\alpha h_*(y)=-K_\alpha.
\end{equation}
Therefore, naturality does not identify distinct values of \(v\): it records
their difference through the coherent homotopy \(K_\alpha\).

\section{Internal Orientation}

On \(P_n^0(c)\) is defined
\begin{equation}\label{eq:detint-omega-p}
 \Omega(r)=r,
 \qquad
 \Omega(q)=-q\quad(q\in\{e,b,f,g\}).
\end{equation}
It is an automorphism of complexes and of the relative coalgebra. In the nerve,
the oriented inversion is
\begin{equation}\label{eq:detint-omega-g}
 \mathfrak o_G[y_0\to\cdots\to y_q]
 :=(-1)^{q(q+1)/2}
 [\bar y_q\to\cdots\to\bar y_0].
\end{equation}
Together with \eqref{eq:detint-or-cv}, these actions give
\begin{equation}\label{eq:detint-or-z-h}
 \Omega z_-(y)=z_-(\bar y),
 \qquad
 \Omega z_+(c_y,v_y)=z_+(c_{\bar y},v_{\bar y}),
 \qquad
 \Omega h_*(y)=h_*(\bar y).
\end{equation}
Orientation acts on path, complex, publications, filler and determinant
line; it is not a label without an action.

\section{Terminal reduction and survival condition}
\label{sec:detint-terminal}

Only the common root is eliminated:
\begin{equation}\label{eq:detint-p-reducido}
 \overline P_{y;m,n}
 :=P^{\rm gen}_{y;m,n}/A_n\mathbf1_y,
 \qquad
 \widehat P_{y;m}^{\rm red}:=\varprojlim_n\overline P_{y;m,n}.
\end{equation}
The base is first ordered by the TPK order of routes and then by
\begin{equation}\label{eq:detint-orden-base}
 f\prec e\prec b\prec g.
\end{equation}
The following internal algorithm is executed.
\begin{enumerate}
 \item Find the coefficients that are units of \(\mathbf Z_3\).
 \item Choose the first according to \eqref{eq:detint-orden-base} and the order of
 paths.
 \item Cancel the corresponding pair and record the elementary
 operations, their sign, and their contribution to the orientation.
 \item Repeat until no cancellable unit pivot remains.
\end{enumerate}
The survival condition of \eqref{eq:detint-dominio} requires that,
cofinally in \(m\), the result be a two-term block
\begin{equation}\label{eq:detint-sx}
 \mathbf Z_3^{r_x}\xrightarrow{S_x}\mathbf Z_3^{r_x},
 \qquad
 \det S_x\ne0\text{ in }K_3,
\end{equation}
and that a refinement only add unit contractible blocks and recorded
changes of basis. If these conditions fail, the history does not belong to
\(\mathcal C_{\rm cont}^{\rm det}\); a nonexistent terminal is not forced.

The identity route provides an elementary witness. After removing the root,
the local complex is
\begin{equation}\label{eq:detint-testigo-contractil}
 A_nf\xrightarrow{\binom{3c}{1}}
 A_ne\oplus A_nb\xrightarrow{(1,-3c)}A_ng.
\end{equation}
Is contractible via
\begin{equation}\label{eq:detint-contraccion-local}
 s(g)=e,
 \qquad
 s(ae+\beta b)=\beta f,
 \qquad
 s(f)=0.
\end{equation}
Thus, the terminal conditions admit at least the internal identity witness
and do not describe a formally empty subdomain.

\section{Smith, \texorpdfstring{\(3\)}{3}-adic order and initial unit}

In oriented bases, a valuation diagonalization has the form
\begin{equation}\label{eq:detint-smith}
 U_xS_xV_x
 =\operatorname{diag}(3^{a_1}u_1,\ldots,3^{a_r}u_r),
 \qquad
 a_i\in\mathbf N,quad u_i\in\mathbf Z_3^\times.
\end{equation}
We define
\begin{equation}\label{eq:detint-orden-d}
 d_x:=v_3(\det S_x)=\sum_{i=1}^r a_i
\end{equation}
and the initial unit
\begin{equation}\label{eq:detint-leading-unit}
 \lambda_x:=3^{-d_x}\det S_x\in\mathbf Z_3^\times.
\end{equation}
The orientation coordinate trivializes the determinant line. If a
transition produces
\begin{equation}\label{eq:detint-estabilidad-bloques}
 S_x\oplus I_q=U S_{x'}V,
\end{equation}
the orientation is transported by the inverse factor of \(\det U\det V\).
The oriented unit \(\lambda_x^{\rm or}\) is then natural and does not change
when contractible blocks are added.

At finite precision,
\begin{equation}\label{eq:detint-smith-finito}
 S_{x,n}=S_x\bmod3^n,
 \qquad
 a_i^{(n)}=\min(a_i,n).
\end{equation}
The terminal is, therefore, the compatible family of all precisions
and not an isolated integer.

\section{Bockstein Tower and Higher Order Carry}

Let
\[
 C_x^1\xrightarrow{S_x}C_x^0
\]
the terminal block. If a class \([\bar y]\) admits a lift
\(y\in C_x^1\) with \(S_xy\in3^qC_x^0\), we define
\begin{equation}\label{eq:detint-beta-q}
 \beta_q[\bar y]
 :=\left[3^{-q}S_xy\bmod3\right].
\end{equation}
Changing the lift alters the representative by a boundary of the corresponding
page, so \(\beta_q\) is well defined. The higher
carry is
\begin{equation}\label{eq:detint-carry-superior}
 \operatorname{Carr}_q(y):=3^{-q}S_xy.
\end{equation}
It does not replace the TPK carry of \eqref{eq:detint-c-v-lifts}; it records its
prolongation in the coefficient tower.

For a block \(3^{a_i}u_i\), the class persists until page \(a_i\),
and on that page the nonzero differential is multiplication by
\(\bar u_i\in\mathbf F_3^\times\). Consequently,
\begin{equation}\label{eq:detint-bock-longitudes}
 \{\text{survival lengths}\}=\{a_1,\ldots,a_r\},
 \qquad
 d_x=\sum_i a_i.
\end{equation}
The oriented product of the page trivializations satisfies
\begin{equation}\label{eq:detint-bock-leading}
 \Theta_{\rm Bock}(x)=\overline{\lambda_x^{\rm or}}
 \in\mathbf F_3^\times.
\end{equation}
The proof is carried out in each diagonal block and transported by the recorded
unimodular changes. Smith, Bockstein and the initial unit are thus three
readings of the same terminal.

\section{The thirteen coordinates of the determinantal constraint}

For each \((m,n)\), it is defined
\begin{equation}\label{eq:detint-d-trece}
 D_{{\rm det},m,n}(x):=
 (\mathcal F,\operatorname{Ext}_\Gamma,\operatorname{Rel},
 \mathcal N,\operatorname{or},\operatorname{Carr},\operatorname{Mem},
 \partial,\gamma,\operatorname{Msect},\operatorname{Surv},
 \operatorname{Term},\mathcal L)_{{\rm det},m,n}(x),
\end{equation}
with the following content.
\begin{enumerate}
 \item The fiber is
 \[
  \mathcal F=(A_n,\mathsf T_m(x),
  \{P_n^0(c_y)\}_y,\{P^{\rm gen}_{y;m,n}\}_y).
 \]
 \item \(\operatorname{Ext}_\Gamma\) contains all morphisms
 \(\alpha\), the maps \(\Psi_\alpha\), the homotopies \(K_\alpha\) and
 the maps \(R_{m',m}^{n',n}\).
 \item \(\operatorname{Rel}\) contains \(\partial^2=0\), the
 pullback equation, AW, counit, \(\Psi_\beta\Psi_\alpha=\Psi_{\beta\alpha}\) and
 \(K_{\beta\alpha}=K_\beta+\Psi_\beta K_\alpha\).
 \item
 \(\mathcal N=(m,n,\widetilde c_y,\widetilde v_y,c_{y,n},v_{y,n},
 r_x,a_1,\ldots,a_{r_x},d_x)\).
 \item \(\operatorname{or}\) contains \(y\mapsto\bar y\),
 \(\mathfrak o_G\), \(\Omega\) and the orientation of the determinant
 line.
 \item \(\operatorname{Carr}\) separately preserves the integer carry,
 its reduction \(c_{y,n}\) and all higher carries
 \(\operatorname{Carr}_q\).
 \item \(\operatorname{Mem}\) contiene
 \[
  y_0\to\cdots\to y_m,
  \quad\{(\widetilde c_{y_j},\widetilde v_{y_j},\mathbf1_{y_j})\}_{j=0}^m,
  \quad\{K_{\alpha_j}\}.
 \]
 \item
 \(\partial=(\partial_G,\partial_P,\partial_{P^{\rm gen}},
 \partial_{\overline P},h_*)\).
 \item \(\gamma=[y_0\to\cdots\to y_q]\) preserves the complete oriented
 simplex.
 \item \(\operatorname{Msect}\) contains all compatible
 extensions, all \(3\)-adic lifts, and all unimodular
 presentations of the same oriented terminal; it does not choose one a posteriori.
 \item \(\operatorname{Surv}\) records arbitrary depth,
 coefficient compatibility, balance, acyclicity over \(K_3\),
 stability under unit blocks and persistence of the root.
 \item
 \[
  \operatorname{Term}=(\widehat P_x^{\rm red},S_x,
  \{a_i,u_i\},d_x,\lambda_x^{\rm or},
  \{\beta_q\},\Theta_{\rm Bock}).
 \]
 \item The internal reader is
 \[
  \mathcal L_{\rm det}^{\rm int}
  =(\mathbf1_y,
  \varpi_{\rm root}z_-=\varpi_{\rm root}z_+,
  [z_-]=[z_+],d_x,\lambda_x^{\rm or},\Theta_{\rm Bock}).
 \]
\end{enumerate}

The transitions satisfy, coordinate by coordinate,
\begin{equation}\label{eq:detint-naturalidad-trece}
 u^{\rm det}_{(m',n'),(m,n)}D_{{\rm det},m',n'}(x')
 =D_{{\rm det},m,n}(\tau_{m',m}x').
\end{equation}

\section{Assembler, publisher, and non-ablative chain}

We define
\begin{equation}\label{eq:detint-graph-d}
 \mathcal G_{\rm det}:=\operatorname{Graph}(D_{\rm det}),
 \qquad
 \operatorname{Res}_{\rm det}(x):=(x,D_{\rm det}x).
\end{equation}
Let \(\chi_{\rm det}(x)\) be the complete multisection of germs
\[
 (y,\widetilde c_y,\widetilde v_y,\operatorname{or}_y,\gamma_y)
\]
compatible at all depths. The dependent object is
\begin{equation}\label{eq:detint-x-dependiente}
 X_{\chi,{\rm det}}
 :=\{(\chi_{\rm det}(x),x,D_{\rm det}x):
 x\in\mathcal C_{\rm cont}^{\rm det}\},
\end{equation}
and
\begin{equation}\label{eq:detint-prx}
 \operatorname{pr}_{X,{\rm det}}(x,D_{\rm det}x)
 :=(\chi_{\rm det}(x),x,D_{\rm det}x).
\end{equation}

The raw assembler
\begin{equation}\label{eq:detint-ensamblador-crudo}
 a_{\rm det}:X_{\chi,{\rm det}}\longrightarrow
 \mathscr M_{\rm det}^{\rm amb}
\end{equation}
yields
\begin{equation}\label{eq:detint-ensamblador-salida}
\begin{split}
 a_{\rm det}(z)=\bigl(&
 \{G_{m,n},\Delta_{\rm AW},\epsilon_G\}_{m,n};\\
 &\{P_n^0(c),\Psi_\alpha,K_\alpha,P^{\rm gen},
 \mathbf1,z_-,z_+,h_*\}_{m,n};\\
 &\{\widehat P^{\rm red},S,d,\lambda^{\rm or},
 \beta_q,\Theta_{\rm Bock}\}\bigr).
\end{split}
\end{equation}
The macroobject preserves its entry:
\begin{equation}\label{eq:detint-graph-a}
 \mathfrak M_{\rm det}:=\operatorname{Graph}(a_{\rm det}),
 \qquad
 \mathcal A_{\rm det}(z):=(z,a_{\rm det}z).
\end{equation}

The raw publisher
\begin{equation}\label{eq:detint-publicador-crudo}
 p_{\rm det}:\mathfrak M_{\rm det}\longrightarrow\mathscr Y_{\rm det}
\end{equation}
publishes the certificate jointly
\begin{equation}\label{eq:detint-certificado-publicado}
\begin{gathered}
 \partial^2=0,quad \text{AW and counit},quad
 \mathbf1_y\text{ compatible group-like},\\
 \partial z_\pm=0,quad
 \varpi_{\rm root}z_-=\varpi_{\rm root}z_+,quad
 \partial h_*=z_--z_+,\\
 \Psi_\beta\Psi_\alpha=\Psi_{\beta\alpha},quad
 K_{\beta\alpha}=K_\beta+\Psi_\beta K_\alpha,\\
 d=v_3(\det S)=\sum_i a_i,quad
 \Theta_{\rm Bock}=\overline{\lambda^{\rm or}},quad
 \text{full naturality under refinement}.
\end{gathered}
\end{equation}
Finally,
\begin{equation}\label{eq:detint-graph-p}
 \mathcal C_{\rm det}^{\rm HMT}:=\operatorname{Graph}(p_{\rm det}),
 \qquad
 \mathcal P_{\rm det}(M):=(M,p_{\rm det}M).
\end{equation}
The complete chain is
\begin{equation}\label{eq:detint-cadena-completa}
 \mathcal C_{\rm cont}^{\rm disc}\supset
 \mathcal C_{\rm cont}^{\rm det}
 \xrightarrow{\operatorname{Res}_{\rm det}}\mathcal G_{\rm det}
 \xrightarrow{\operatorname{pr}_{X,{\rm det}}}X_{\chi,{\rm det}}
 \xrightarrow{\mathcal A_{\rm det}}\mathfrak M_{\rm det}
 \xrightarrow{\mathcal P_{\rm det}}\mathcal C_{\rm det}^{\rm HMT}.
\end{equation}
Each arrow has an inverse on its image given by a coordinate
projection. History is not replaced by value, multisection by selector,
complex by determinant or proof by name.

\section{Theorem of Internal Determinantal Realization}

\begin{theorem}[Complete internal determinant realization]
\label{thm:detint-realizacion}
For every \(x\in\mathcal C_{\rm cont}^{\rm det}\), the chain
\eqref{eq:detint-cadena-completa} constructs a macro-object natural in
depth and precision. It contains two cycles with a common root, an explicit filler
that identifies them, a compatible group-like unit, and a
\(3\)-adic terminal whose order, Smith form, and Bockstein tower produce the same
oriented reader. The thirteen coordinates of \eqref{eq:detint-d-trece}
remain recoverable in the output.
\end{theorem}

\begin{proof}
Equations \eqref{eq:detint-p0-diferencial}--
\eqref{eq:detint-d2} prove that \(P_n^0(c)\) is a complex. The two
augmentations are maps of complexes, so the pullback is typed.
Alexander--Whitney is natural and makes each vertex group-like; together with
\(r\), this gives the common unit. The direct calculation
\eqref{eq:detint-zpm-ciclos}--\eqref{eq:detint-filler-identidad} proves that
the publications are cycles, have the same root and are joined by
\(h_*\). The formulas for \(\Psi_\alpha\) prove functoriality, and
\(K_\alpha\) proves homotopic naturality while preserving the memory
increment. The definition of \(\mathcal C_{\rm cont}^{\rm det}\) ensures that
terminal reduction produces a square block without postulating it outside the
domain. Smith over \(\mathbf Z_3\) gives \(d=\sum_i a_i\); the calculation of each
block \(3^{a_i}u_i\) gives
\(\Theta_{\rm Bock}=\overline{\lambda^{\rm or}}\). Truncation, reduction and
addition of unit blocks preserve these identities. Finally, the
three graph constructions retain the input as the first coordinate,
so the complete composition is non-ablative.
\end{proof}

\section{Provenance and Forgers}

\paragraph{Origin.}
The TPK nerve, Alexander--Whitney, the counit, the common unit, the normal
form \(\partial f=3ce+b\), \(\partial e=g\),
\(\partial b=-3cg\), the common root, the filler, and the tower architecture
have the status of \texttt{PREEXISTING\_AUTHORIAL\_ARCHITECTURE}. The literal computation
of the two cycles, their root, and the identity
\(\partial h_*=z_--z_+\) has the status of a \texttt{RECOVERED\_RESULT}. The
local system \(\Psi_\alpha/K_\alpha\), the balanced surviving subdomain,
the exclusively \(3\)-adic terminal, and the packaging by
graphs have the status of a \texttt{NEW\_FORMALIZATION}; they formalize without altering
the conceptual provenance of the apparatus.

\begin{ablation}[Internal falsifiers of the determinantal branch]
\label{abl:detint-falsadores}
The construction is falsified at the corresponding level if any of
the following facts occurs.
\begin{enumerate}
 \item \(\partial^2=0\) fails; then the local complex does not exist.
 \item Alexander--Whitney does not commute with refinement, the counit fails, or
 the vertex ceases to be group-like.
 \item \(z_\pm\) are not cycles, their roots differ, or
 \(\partial h_*=z_--z_+\) fails.
 \item \(\Psi_\alpha\) is not a map of complexes or
 \(K_{\beta\alpha}=K_\beta+\Psi_\beta K_\alpha\) fails; in that case
 carry or memory has been lost.
 \item The reduced complex is not balanced or is not acyclic over
 \(K_3\); the history lies outside the typed terminal domain.
 \item The Bockstein lengths do not coincide with the \(a_i\), or
 \(\Theta_{\rm Bock}\ne\overline{\lambda^{\rm or}}\).
 \item Refinement modifies \(d\) or \(\lambda^{\rm or}\) outside
 unit blocks and recorded orientation changes.
 \item One of the thirteen coordinates is omitted: \emph{object replaced;
 conclusion not transferable}.
\end{enumerate}
These forgers check the internally constructed object; they do not replace it with a subsequent control.
\end{ablation}
\fi

\endgroup

\section{Cellular incidence and correlative filling of the same trajectory}

The sequence \eqref{eq:pdfv2-cor-sucesion-celda} assigns to each prefix the
two-term complex
\begin{equation}
\label{eq:pdfv2-cor-cono-celda}
 C_{k,N}(\widehat x):=
 \left[A_N^{I_k}\oplus A_N^{J_k}
 \xrightarrow{\kappa_{k,N}}A_N^{I_k\times J_k}\right].
\end{equation}
A refinement acts on ports by truncating the extension and on coefficients
by \(A_{N'}\twoheadrightarrow A_N\). When \(\alpha_u:I_{k'}\to I_k\) and \(\beta_u:J_{k'}\to J_k\) are those maps, their
action on generators is
\begin{equation}
\label{eq:pdfv2-cor-refinamiento-celda}
 (p,q)\longmapsto(p\circ\alpha_u,q\circ\beta_u),
 \qquad
 w\longmapsto w\circ(\alpha_u\times\beta_u),
\end{equation}
and satisfies
\begin{equation}
\label{eq:pdfv2-cor-refinamiento-cadena}
 \kappa_{k',N}\Gamma(u)_1=\Gamma(u)_0\kappa_{k,N}.
\end{equation}

The incidence is not fabricated by labeling four copies of a neutral emission.
It is generated by the free cellular completion of the module of the two sheets that
the same state already bears. For
\[
 V_{\widehat x}:=\mathbf F_3e_{\Sigma,\widehat x}
 \oplus\mathbf F_3e_{\Pi,\widehat x},
\]
the four internal lines are
\begin{equation}
\label{eq:pdfv2-cor-cuatro-rectas-tpk}
 \mathscr I_{\widehat x}:=
 \{L_0,L_\infty,L_+,L_-\}
 =\mathbf P^1(V_{\widehat x}),\qquad
 \begin{aligned}
 L_0&=\langle e_\Sigma\rangle,&
 L_\infty&=\langle e_\Pi\rangle,\\
 L_+&=\langle e_\Sigma+e_\Pi\rangle,&
 L_-&=\langle e_\Sigma-e_\Pi\rangle.
 \end{aligned}
\end{equation}
If \(v_L\) is, respectively,
\((1,0)^t,(0,1)^t,(1,1)^t,(1,-1)^t\), the action of the direction \(L\)
on the generators of sheet is the proyector
\begin{equation}
\label{eq:pdfv2-cor-proyectores-cuatro-direcciones}
 P_L:=v_L(v_L^tv_L)^{-1}v_L^t,
\end{equation}
that is, over \(\mathbf F_3\),
\[
 P_0=\begin{pmatrix}1&0\\0&0\end{pmatrix},\quad
 P_\infty=\begin{pmatrix}0&0\\0&1\end{pmatrix},\quad
 P_+=\begin{pmatrix}2&2\\2&2\end{pmatrix},\quad
 P_-=\begin{pmatrix}2&1\\1&2\end{pmatrix}.
\]
Direct calculation gives
\begin{equation}
\label{eq:pdfv2-cor-proyectores-no-colapso}
 P_L^2=P_L,\qquad \operatorname{rk}P_L=1,\qquad
 \operatorname{Im}P_L=L,\qquad
 P_L=P_{L'}\Longrightarrow L=L'.
\end{equation}
Therefore, the four directions are distinct outputs of the module
\(V_{\widehat x}\); they are not four names placed upon the same action.

\begin{definition}[Cell-free leaf TPK completion]
\label{def:pdfv2-cor-completacion-celular-tpk}
Let \(\operatorname{Cell}_{\rm TPK}(\widehat x)\) be the free complex
\[
 C_2(\widehat x)\xrightarrow{\partial_2}
 C_1(\widehat x)\xrightarrow{\partial_1}C_0(\widehat x)
\]
with
\begin{align*}
 C_0(\widehat x)
 &:=
 \mathbf F_3[v_\star]\oplus
 \bigoplus_{L\in\mathscr I}\mathbf F_3[v_L]\oplus
 \bigoplus_{(L,\epsilon)\in\mathscr B}
       \mathbf F_3[v_{L,\epsilon}]\oplus
 \bigoplus_{E\in\mathscr E}\mathbf F_3[v_E],\\
 C_1(\widehat x)
 &:=
 \bigoplus_{L\in\mathscr I}\mathbf F_3[u_L]\oplus
 \bigoplus_{(L,\epsilon)\in\mathscr B}
       \mathbf F_3[u_{L,\epsilon}]\oplus
 \bigoplus_{E=\{L,L'\}\in\mathscr E}
 \bigl(\mathbf F_3[u_{L'\mid L}]
       \oplus\mathbf F_3[u_{L\mid L'}]\bigr),\\
 C_2(\widehat x)
 &:=
 \bigoplus_{E\in\mathscr E}\mathbf F_3[\omega_E].
\end{align*}
Its action on generators is
\begin{align}
 \partial u_L&=v_L-v_\star,
 &
 \partial u_{L'\mid L}&=v_E-v_L,
 \label{eq:pdfv2-cor-cell-frontera-aristas}\\
 \partial u_{L,\epsilon}&=v_{L,\epsilon}-v_L,
 &
 \partial\omega_{\{L,L'\}}
 &=u_L+u_{L'\mid L}-u_{L'}-u_{L\mid L'}.
 \label{eq:pdfv2-cor-cell-frontera-celda}
\end{align}
\end{definition}

The eight generators \(v_{L,\epsilon}\) materially realize the screen
\(\mathscr B\). The preceding involution extends to the complex by
\begin{equation}
\label{eq:pdfv2-cor-pantalla-ocho-generadores}
 \vartheta_Lv_{K,\epsilon}:=v_{\vartheta_L(K,\epsilon)},\qquad
 \vartheta_Lu_{K,\epsilon}:=u_{\vartheta_L(K,\epsilon)},
\end{equation}
and fixes the other generators. Since the basis is free,
\begin{equation}
\label{eq:pdfv2-cor-pantalla-ocho-no-colapso}
 |\{v_{L,\epsilon}:(L,\epsilon)\in\mathscr B\}|=8,\qquad
 \vartheta_L^2=I,\qquad
 \vartheta_L\vartheta_{L'}=\vartheta_{L'}\vartheta_L.
\end{equation}
Moreover, \(\partial\vartheta_L=\vartheta_L\partial\), because both paths
send \(u_{L,\epsilon}\) to \(v_{L,-\epsilon}-v_L\) and fix the remaining
boundaries.

The two words
\[
 p_{L,L'}:=u_{L'\mid L}u_L,\qquad
 p_{L',L}:=u_{L\mid L'}u_{L'}
\]
are different paths in the free complex. Their difference is the boundary of
\(\omega_{\{L,L'\}}\), and
\begin{equation}
\label{eq:pdfv2-cor-cell-dos-cero}
 \partial^2\omega_E
 =(v_L-v_\star)+(v_E-v_L)
 -(v_{L'}-v_\star)-(v_E-v_{L'})=0.
\end{equation}
The representation of TPK on the leaves is fixed by
\[
 \varrho(u_L)=P_L,\qquad
 \varrho(u_{L'\mid L})=P_{L'},\qquad
 \varrho(p_{L,L'})=P_{L'}P_L.
\]
Even if two published matrices coincided in a specialization, the two
paths do not collapse: they remain distinct free generators of
\(C_1(\widehat x)\), joined by an explicit 2--cell.

The complete incidence coordinate is the group of 2--cochains.
\begin{equation}
\label{eq:pdfv2-cor-q-incidencia-total}
 Q_{\rm inc}(\widehat x):=
 C^2(\operatorname{Cell}_{\rm TPK}(\widehat x);\mathbf F_3)
 =\bigoplus_{E\in\mathscr E}\mathbf F_3\omega_E^*
 \simeq\mathbf F_3^6.
\end{equation}
The internal operation returns all of \(Q_{\rm inc}(\widehat x)\), not a
selected element. Each \(q\) has a unique expansion
\(q=\sum_Eq_E\omega_E^*\). For
\(a\in\mathbf F_3^{\mathscr I}\), its change of potential is
\begin{equation}
\label{eq:pdfv2-cor-collar-variacion}
 q\longmapsto q+Ba,\qquad (Ba)_{\{L,L'\}}=a_L+a_{L'}.
\end{equation}
As each line belongs to three pairs,
\[
 s(Ba)=3\sum_{L\in\mathscr I}a_L=0,
\]
and therefore
\begin{equation}
\label{eq:pdfv2-cor-wc-invariante-celular}
 W_c(q+Ba)=W_c(q)\qquad(q\in Q_{\rm inc}(\widehat x)).
\end{equation}
The unit of the completion only anchors the complex in the state:
\begin{equation}
\label{eq:pdfv2-cor-cell-unidad}
 \eta^{\rm cell}_{\widehat x}:V_{\widehat x}\longrightarrow
 V_{\widehat x}\otimes C_0(\operatorname{Cell}_{\rm TPK}(\widehat x)),\qquad
 e_{\diamond,\widehat x}\longmapsto
 e_{\diamond,\widehat x}\otimes v_\star,
\end{equation}
and marks the origin \(q=0\); it does not assert that the orbit of that origin is all of
\(Q_{\rm inc}\).

The construction is functorial under common truncation. If
\(\tau_{\ell,k}\widehat x_\ell=\widehat x_k\), we define
\begin{align*}
 \tau_V&:V_{\widehat x_\ell}\longrightarrow V_{\widehat x_k},&
 \tau_Ve_{\diamond,\widehat x_\ell}&=e_{\diamond,\widehat x_k},\\
 \tau_a&:\mathbf F_3^{\mathscr I_{\widehat x_\ell}}
          \longrightarrow\mathbf F_3^{\mathscr I_{\widehat x_k}},&
 (\tau_aa)_{\tau_{\mathscr I}L}&=a_L,\\
 \tau_Q&:Q_{\rm inc}(\widehat x_\ell)
          \longrightarrow Q_{\rm inc}(\widehat x_k),&
 (\tau_Qq)_{\tau_{\mathscr E}E}&=q_E,
\end{align*}
where \(\tau_{\mathscr I}[v]=[\tau_Vv]\) and
\(\tau_{\mathscr E}\{L,L'\}=\{\tau_{\mathscr I}L,
\tau_{\mathscr I}L'\}\). The map of cadenas
\(\operatorname{Cell}(\tau):C_\bullet(\widehat x_\ell)
\to C_\bullet(\widehat x_k)\) acts in generators by
\begin{equation}
\label{eq:pdfv2-cor-cell-truncamiento-generadores}
 v_\bullet(\widehat x_\ell)\mapsto v_\bullet(\widehat x_k),\quad
 u_\bullet(\widehat x_\ell)\mapsto u_\bullet(\widehat x_k),\quad
 \omega_E(\widehat x_\ell)\mapsto\omega_E(\widehat x_k).
\end{equation}
From \eqref{eq:pdfv2-cor-cell-frontera-aristas}--
\eqref{eq:pdfv2-cor-cell-frontera-celda} one obtains
\begin{equation}
\label{eq:pdfv2-cor-collar-naturalidad}
\begin{aligned}
 \operatorname{Cell}(\tau)\partial
 &=\partial\operatorname{Cell}(\tau),\\
 \tau_VP_L&=P_{\tau_{\mathscr I}L}\tau_V,\qquad
 \tau_QB=B\tau_a,\qquad W_c\circ\tau_Q=W_c,\\
 \operatorname{Cell}(\tau)\vartheta_L
 &=\vartheta_{\tau_{\mathscr I}L}\operatorname{Cell}(\tau),\\
 (\tau_V\otimes\operatorname{Cell}_0(\tau))
 \eta^{\rm cell}_{\widehat x_\ell}
 &=\eta^{\rm cell}_{\widehat x_k}\tau_V.
\end{aligned}
\end{equation}
Identity and composition are verified on each generator. Thus
\(\operatorname{Cell}_{\rm TPK}\) is the functorial cellular completion of
the same APP--TRIT--TPK fiber, not a carrier added from a subsequent
publication.

\subsection{Incidence measure, carry curvature, and adjoint naturality}
In the following development, they are identified exclusively as abbreviations
of the common completion,
\[
 I_{\rm gau}:=\mathscr I_{\widehat x},\qquad
 \mathscr B_{\rm gau}:=\mathscr B_{\widehat x},\qquad
 Q_{\rm gau}:=Q_{\rm inc}(\widehat x).
\]
With the uniform counting average on \(Q_{\rm gau}\), the function
\(W_c\) of \eqref{eq:pdfv2-cor-wc} satisfies
\begin{equation}
\label{eq:pdfv2-cor-wc-momentos}
 \mathbb E_QW_c=0,\qquad
 \mathbb E_QW_c^2=\frac29,\qquad
 \mathbb E_QW_c^3=\frac{2}{27}.
\end{equation}
These identities are properties of the same cell and the same carry, not
data from an external theory.
\begingroup
  \let\section\subsection
  % Extracción quirúrgica del propietario V2 05c: líneas 99--335.
% El resto permanece en este archivo como procedencia, pero no se publica.
\iffalse
\chapter[Internal Gauge Macrostructure]{Internal Gauge Macrostructure Generated by the Discrete Continuum}
\label{chap:gauint-macroestructura}

\begin{thesisbox}
The macrostructure of this chapter is constructed solely from a
surviving restriction of \(\mathcal C_{\rm cont}^{\rm disc}\). The numbers
\(4\), \(6\), and \(8\), the screen, the four involutions, the
projective measure, the transfer operator, the incidence mode, and the cellular
terminal are derived within that restriction. No subsequent evaluation
forms part of the domain or of the arrows that follow.
\end{thesisbox}

\section{The projective seed and the internal derivation of \texorpdfstring{\(4/6/8\)}{4/6/8}}

Let \(V_{\rm TPK}=\mathbf F_3^2\), with the ordered basis inherited from the two
sheets of the TPK state. We define the set of directions
\begin{equation}\label{eq:gauint-cuatro-direcciones}
 I_{\rm gau}:=\mathbf P^1(\mathbf F_3)
 =\{L_\infty,L_0,L_1,L_{-1}\},
\end{equation}
where
\[
 L_\infty=[1:0],\qquad L_0=[0:1],\qquad
 L_1=[1:1],\qquad L_{-1}=[1:-1].
\]
Therefore,
\begin{equation}\label{eq:gauint-cardinal-cuatro}
 |I_{\rm gau}|=\frac{3^2-1}{3-1}=4.
\end{equation}
A set of four indices is not introduced: it appears as the quotient of
the eight nonzero vectors by the action of \(\mathbf F_3^\times\).

The unoriented incidences are
\begin{equation}\label{eq:gauint-seis-incidencias}
 E_{\rm gau}:=\binom{I_{\rm gau}}2,
 \qquad |E_{\rm gau}|=\binom42=6,
 \qquad Q_{\rm gau}:=\mathbf F_3^{E_{\rm gau}}.
\end{equation}
The oriented screen duplicates each direction, not each incidence:
\begin{equation}\label{eq:gauint-ocho-caras}
 \mathscr B_{\rm gau}:=I_{\rm gau}\times\{+,-\},
 \qquad |\mathscr B_{\rm gau}|=2|I_{\rm gau}|=8.
\end{equation}
Thus, \(4\), \(6\), and \(8\) arise, respectively, from directions, pairs of
directions, and oriented faces of the same seed.

For each \(L\in I_{\rm gau}\), let \(\vartheta_L\) be the screen
involution that exchanges \((L,+)\) and \((L,-)\) and fixes the other six
faces. Then
\begin{equation}\label{eq:gauint-cuatro-involuciones}
 \vartheta_L^2=I,
 \qquad
 \vartheta_L\vartheta_{L'}=\vartheta_{L'}\vartheta_L
 \quad(L\ne L').
\end{equation}
The four involutions have been obtained prior to defining measure or operator.

\section{Incidence, internal mode and non-constant witness}

The incidence application is
\begin{equation}\label{eq:gauint-incidencia-b}
 B:\mathbf F_3^{I_{\rm gau}}\longrightarrow Q_{\rm gau},
 \qquad
 (Ba)_{\{L,L'\}}:=a_L+a_{L'}.
\end{equation}
If
\[
 s(q):=\sum_{e\in E_{\rm gau}}q_e\in\mathbf F_3,
\]
each direction appears exactly three times in the sum of the six
incidences, and therefore
\begin{equation}\label{eq:gauint-incidencia-invariante}
 s(Ba)=3\sum_{L\in I_{\rm gau}}a_L=0,
 \qquad
 s(q+Ba)=s(q).
\end{equation}

On \(Q_{\rm gau}\), only the normalized counting average is used.
The incidence mode is the rational function
\begin{equation}\label{eq:gauint-wc-definicion}
 W_c(q):=\mathbf 1_{\{s(q)=0\}}-\frac13.
\end{equation}
Since \(s:Q_{\rm gau}\to\mathbf F_3\) is surjective, each fiber has \(3^5\) elements.
Exact counting then yields,
\begin{equation}\label{eq:gauint-wc-momentos}
 \mathbb E_{Q}W_c=0,
 \qquad
 \mathbb E_{Q}W_c^2=\frac29,
 \qquad
 \mathbb E_{Q}W_c^3=\frac{2}{27}.
\end{equation}
Furthermore, \eqref{eq:gauint-incidencia-invariante} gives
\begin{equation}\label{eq:gauint-wc-invariancia}
 W_c(q+Ba)=W_c(q).
\end{equation}
Therefore, \(W_c\) is nonzero, centered, and invariant under the internally
generated incidence.

\fi
\section{Incidence, carry curvature, and common domain of naturalness}

For \(x\in\mathcal C_{\rm cont}^{\rm disc}\), let us denote by
\(Y_m(x)\) the finite set of gauge prefixes of depth \(m\) that
simultaneously conserve
\begin{equation*}
\begin{gathered}
 \text{fiber, extension relations, numbers, orientation, carry, memory,}\\
 \text{boundary, route, multisection, survival, terminal and reader}.
\end{gathered}
\end{equation*}
There exists a collar neighborhood application.
\begin{equation}\label{eq:gauint-collar}
 q_m:Y_m(x)\longrightarrow Q_{\rm gau}
\end{equation}
which records the six incidences without erasing the history that produces them.
For \(\ell\ge m\), the truncation is written
\[
 \tau_{\ell,m}:Y_\ell(x)\longrightarrow Y_m(x).
\]

For \(y,z\in Y_m(x)\), let
\begin{equation}\label{eq:gauint-todas-extensiones}
 \operatorname{Ext}^{\rm adm}_{m}(y,z)
\end{equation}
the set of \emph{all} admissible TPK extensions from \(y\) to \(z\).
An extension belongs to this set only if it transports the entire fiber,
the relations \(\operatorname{Ext}_\Gamma\), the native and operatorial carry,
the orientation, the accumulated memory, the boundary, the route, and the
survival condition. We define
\begin{equation}\label{eq:gauint-conteo-kernel}
 a_m(y,z):=\#\operatorname{Ext}^{\rm adm}_m(y,z),
 \qquad
 d_m(y):=\sum_{z\in Y_m(x)}a_m(y,z),
\end{equation}
and, when \(d_m(y)>0\),
\begin{equation}\label{eq:gauint-kernel-normalizado}
 k_m(y,z):=\frac{a_m(y,z)}{d_m(y)}.
\end{equation}
Thus, \(k_m\) is a normalized count of complete extensions, not a
function introduced after projecting the state.

The restriction \(\mathcal C_{\rm cont}^{\rm gau}\) is formed by the
surviving histories \(x\) for which the following conditions
hold at every depth.
\begin{enumerate}
 \item The sets \(Y_m(x)\) are nonempty, \(d_m(y)>0\), and the
 truncations \(\tau_{\ell,m}\) are surjective.
 \item There exists a nonempty multisection of strictly positive rational
 weights \(\mu_m:Y_m(x)\to\mathbf Q_{>0}\) such that
 \begin{equation}\label{eq:gauint-medida-proyectiva}
  \sum_{y\in Y_m}\mu_m(y)=1,
  \qquad
  \mu_m(y)=\sum_{\tau_{\ell,m}\widetilde y=y}\mu_\ell(\widetilde y).
 \end{equation}
 \item Each \(\mu_m\) is stationary:
 \begin{equation}\label{eq:gauint-estacionariedad}
  \sum_{y\in Y_m}\mu_m(y)k_m(y,z)=\mu_m(z).
 \end{equation}
 \item The marginal of \((q_m)_*\mu_m\) is the uniform counting measure on
 \(Q_{\rm gau}\). This condition makes the calculations
 of \eqref{eq:pdfv2-cor-wc-momentos} literally applicable to the lift \(W_c\circ q_m\).
 \item Direct lumpability holds
 \begin{equation}\label{eq:gauint-lump-directa}
  \sum_{\tau_{\ell,m}\widetilde z=z}
   k_\ell(\widetilde y,\widetilde z)
  =k_m(\tau_{\ell,m}\widetilde y,z)
 \end{equation}
 for every \(\widetilde y\in Y_\ell\) and \(z\in Y_m\).
 \item For the adjoint kernel
 \begin{equation}\label{eq:gauint-kernel-adjunto}
  k_m^*(z,y):=\frac{\mu_m(y)k_m(y,z)}{\mu_m(z)},
 \end{equation}
 the adjoint lumpability condition is satisfied
 \begin{equation}\label{eq:gauint-lump-adjunta}
  \sum_{\tau_{\ell,m}\widetilde y=y}
   k_\ell^*(\widetilde z,\widetilde y)
  =k_m^*(\tau_{\ell,m}\widetilde z,y).
 \end{equation}
 \item The categories of prefixes possess the TPK root as an initial object,
 so that their nerve has the explicit terminal contraction of
 \eqref{eq:gauint-sdr} below.
\end{enumerate}
These conditions define the typed domain; outside it, a measure or natural
operator is not published by mere declaration.

\section{Projective measure, kernel, and direct and adjoint naturality}

Having fixed a branch of the multisection \(\{\mu_m\}_m\), let
\begin{equation}\label{eq:gauint-hm}
 H_m:=\operatorname{Fun}(Y_m,\mathbf Q),
 \qquad
 \langle f,g\rangle_m:=\sum_{y\in Y_m}\mu_m(y)f(y)g(y).
\end{equation}
The operator generated by all extensions is
\begin{equation}\label{eq:gauint-km}
 (K_mf)(y):=\sum_{z\in Y_m}k_m(y,z)f(z).
\end{equation}
The Markov equations and stationarity give
\begin{equation}\label{eq:gauint-markov-estacionario}
 K_m\mathbf1=\mathbf1,
 \qquad
 K_m^*\mathbf1=\mathbf1,
 \qquad
 \|K_mf\|_m\le\|f\|_m.
\end{equation}

The lift and the truncated expectation are
\begin{equation}\label{eq:gauint-lift-esperanza}
 J_{\ell,m}f:=f\circ\tau_{\ell,m},
 \qquad
 E_{m,\ell}:=J_{\ell,m}^*.
\end{equation}
The projectivity of \(\mu\) implies
\begin{equation}\label{eq:gauint-ej-identidad}
 E_{m,\ell}J_{\ell,m}=I_{H_m}.
\end{equation}
The two lumpability conditions discharge, respectively,
\begin{equation}\label{eq:gauint-naturalidad-directa-adjunta}
 K_\ell J_{\ell,m}=J_{\ell,m}K_m,
 \qquad
 K_\ell^*J_{\ell,m}=J_{\ell,m}K_m^*.
\end{equation}
The second equality is not inferred from the first without the attached condition:
both remain registered in the domain.

The positive internal operator is defined
\begin{equation}\label{eq:gauint-transferencia-t}
 T_m:=K_m^*K_m.
\end{equation}
Then
\begin{equation}\label{eq:gauint-t-positivo-natural}
 \langle f,T_mf\rangle_m=\|K_mf\|_m^2\ge0,
 \qquad
 T_\ell J_{\ell,m}=J_{\ell,m}T_m.
\end{equation}

\section{Common law of \(8\) faces and \(4\) Fubini squares}

On \(Y_m^{\mathscr B_{\rm gau}}\), we define the rational mass
\begin{equation}\label{eq:gauint-ley-ocho-caras}
 \Omega_m^{(8)}(\mathbf y)
 :=\sum_{z\in Y_m}\mu_m(z)
 \prod_{(L,\varepsilon)\in\mathscr B_{\rm gau}}
 k_m(z,y_{L,\varepsilon}).
\end{equation}
It is a probability because each factor is Markov. All faces arise
from the same latent state \(z\) and the same kernel; they are not eight independently
added measures.

For functions \(f_{L,+},f_{L,-}\in H_m\), finite Fubini gives
\begin{equation}\label{eq:gauint-fubini-ocho}
\begin{aligned}
 &\sum_{\mathbf y}\Omega_m^{(8)}(\mathbf y)
  \prod_{L\in I_{\rm gau}}
  f_{L,+}(y_{L,+})f_{L,-}(y_{L,-})\\
 &\qquad=
 \sum_{z\in Y_m}\mu_m(z)
 \prod_{L\in I_{\rm gau}}
 (K_mf_{L,+})(z)(K_mf_{L,-})(z).
\end{aligned}
\end{equation}
Upon taking \(f_{L,+}=f_{L,-}=f_L\), the four squares appear
simultaneously
\begin{equation}\label{eq:gauint-cuatro-cuadrados}
 \sum_{z\in Y_m}\mu_m(z)
 \prod_{L\in I_{\rm gau}}(K_mf_L(z))^2.
\end{equation}
The marginal of any opposite pair satisfies
\begin{equation}\label{eq:gauint-marginal-t}
 \sum_{y_+,y_-}\Omega_{m,L}^{(2)}(y_+,y_-)
 f(y_+)g(y_-)
 =\langle K_mf,K_mg\rangle_m
 =\langle f,T_mg\rangle_m.
\end{equation}
The commutativity of the product in \eqref{eq:gauint-ley-ocho-caras} proves
\begin{equation}\label{eq:gauint-involuciones-ley}
 (\vartheta_L)_*\Omega_m^{(8)}=\Omega_m^{(8)}
 \qquad(L\in I_{\rm gau}).
\end{equation}

\section{Exact internal operator and self-test witness}

The projection onto the constant sector and its complement are
\begin{equation}\label{eq:gauint-pq}
 P_{\Omega,m}f:=\langle\mathbf1,f\rangle_m\mathbf1,
 \qquad
 Q_m:=I-P_{\Omega,m}.
\end{equation}
The internal cell operator is defined by
\begin{equation}\label{eq:gauint-d-hmt}
 D_m^{\rm HMT}:=3Q_m.
\end{equation}
The coefficient \(3\) is the cardinality of the residue field used in
\eqref{eq:pdfv2-cor-proyectores-cuatro-direcciones}; it does not represent an
added scale.
For the lift
\[
 W_{c,m}:=W_c\circ q_m
\]
one has, by \eqref{eq:pdfv2-cor-wc-momentos},
\begin{equation}\label{eq:gauint-w-eigen}
 P_{\Omega,m}W_{c,m}=0,
 \qquad
 D_m^{\rm HMT}W_{c,m}=3W_{c,m},
 \qquad
 \|W_{c,m}\|_m^2=\frac29.
\end{equation}
Thus, the complementary sector possesses a nonzero witness constructed by the
incidence of the six edges.

\section{Cellular terminal and strong retractor}

Let \(C_*(N\mathsf Y_m;A_n)\) be the normalized chain complex of the nerve of the
prefix category, and let \(y_*\) be its initial root. The inclusion and the
augmentation are
\[
 \iota_m:A_n[0]\longrightarrow C_*(N\mathsf Y_m;A_n),
 \quad \iota_m(1)=[y_*],
 \qquad
 \epsilon_m:C_*(N\mathsf Y_m;A_n)\longrightarrow A_n[0].
\]
The initial object provides the extra degeneracy
\begin{equation}\label{eq:gauint-homotopia-celular}
 H_m[y_0\to\cdots\to y_q]
 :=[y_*\to y_0\to\cdots\to y_q],
\end{equation}
with the normalized convention when a degeneration appears. It is verified
\begin{equation}\label{eq:gauint-sdr}
 \epsilon_m\iota_m=I,
 \qquad
 \partial H_m+H_m\partial=I-\iota_m\epsilon_m,
 \qquad
 H_m^2=H_m\iota_m=\epsilon_mH_m=0.
\end{equation}
This cellular terminal is distinct from the constant projection
\(P_{\Omega,m}\): the former contracts the nerve, while the latter separates the
sectors of the function space. Both are preserved.

\iffalse
\section{The thirteen coordinates of the gauge constraint}

For each \((m,n)\), it is defined
\begin{equation}\label{eq:gauint-d-trece}
 D_{{\rm gau},m,n}(x):=
 (\mathcal F,\operatorname{Ext}_\Gamma,\operatorname{Rel},
 \mathcal N,\operatorname{or},\operatorname{Carr},\operatorname{Mem},
 \partial,\gamma,\operatorname{Msect},\operatorname{Surv},
 \operatorname{Term},\mathcal L)_{{\rm gau},m,n}(x),
\end{equation}
with the following content.
\begin{enumerate}
 \item \(\mathcal F=(V_{\rm TPK},I_{\rm gau},E_{\rm gau},
 Q_{\rm gau},\mathscr B_{\rm gau},Y_m,H_m)\).
 \item \(\operatorname{Ext}_\Gamma\) contains all the sets
 \(\operatorname{Ext}^{\rm adm}_m(y,z)\), truncations, lifts,
 expectations, and their compositions.
 \item \(\operatorname{Rel}\) contains \(B\), \(sB=0\), the four
 involutions, Markov, stationarity, and both lumpabilities.
 \item \(\mathcal N=(3,4,6,8,m,n,|Y_m|,a_m,d_m)\); those numbers do not
 replace any of the remaining coordinates.
 \item \(\operatorname{or}\) preserves the faces \(+/-\), the four
 inversions \(\vartheta_L\), and the reversal of each path.
 \item \(\operatorname{Carr}\) contains the complete carry of every extension counted in \(a_m(y,z)\), including the carry already propagated.
 \item \(\operatorname{Mem}\) conserves the full prefix, the common
 state \(z\) of \(\Omega_m^{(8)}\), and the memory of its eight extensions.
 \item \(\partial=(q_m,\partial_{N\mathsf Y_m},Q_m,H_m)\) preserves
 collar boundary, cellular boundary, and complement.
 \item \(\gamma\) is the oriented simplex of the nerve that records the route
 and not only its endpoints.
 \item \(\operatorname{Msect}\) is the family of all compatible
 projective stationary measures and of all the extensions that
 enter into the counts; no retrospective branch is selected.
 \item \(\operatorname{Surv}\) records nonemptiness, arbitrary
 depth, projectivity, both lumpabilities, and the persistence of the
 uniform collar marginal.
 \item \(\operatorname{Term}=(\iota_m,\epsilon_m,H_m;
 P_{\Omega,m},Q_m)\) keeps the two terminals separate.
 \item \(\mathcal L=(\Omega_m^{(8)},K_m,T_m,D_m^{\rm HMT},W_{c,m})\)
 is the internal reader posterior to the construction of the history.
\end{enumerate}

For \(\ell\ge m\) and \(n'\ge n\), all the coordinates satisfy
\begin{equation}\label{eq:gauint-naturalidad-trece}
 u^{\rm gau}_{(\ell,n'),(m,n)}D_{{\rm gau},\ell,n'}(x')
 =D_{{\rm gau},m,n}(\tau_{\ell,m}x').
\end{equation}

\section{Assembler, publisher, and non-ablative chain}

Let
\begin{equation}\label{eq:gauint-graph-d}
 \mathcal G_{\rm gau}:=\operatorname{Graph}(D_{\rm gau}),
 \qquad
 \operatorname{Res}_{\rm gau}(x):=(x,D_{\rm gau}x).
\end{equation}
The dependent multisection \(\chi_{\rm gau}(x)\) contains all compatible branches
of collar, measure, and extension. We define
\begin{equation}\label{eq:gauint-x-dependiente}
 X_{\chi,{\rm gau}}
 :=\{(\chi_{\rm gau}(x),x,D_{\rm gau}x):
 x\in\mathcal C_{\rm cont}^{\rm gau}\},
\end{equation}
\begin{equation}\label{eq:gauint-prx}
 \operatorname{pr}_{X,{\rm gau}}(x,D_{\rm gau}x)
 :=(\chi_{\rm gau}(x),x,D_{\rm gau}x).
\end{equation}

The raw assembler
\begin{equation}\label{eq:gauint-ensamblador-crudo}
 a_{\rm gau}:X_{\chi,{\rm gau}}\longrightarrow
 \mathscr M_{\rm gau}^{\rm amb}
\end{equation}
produce the complete family
\[
 \bigl(I,E,\mathscr B,B,W_c;
 Y_m,\mu_m,k_m,J_{\ell,m},E_{m,\ell};
 K_m,T_m,\Omega_m^{(8)},P_{\Omega,m},Q_m,D_m^{\rm HMT};
 \iota_m,\epsilon_m,H_m\bigr)_{m,n}.
\]
The macroobject preserves its entry through
\begin{equation}\label{eq:gauint-graph-a}
 \mathfrak M_{\rm gau}:=\operatorname{Graph}(a_{\rm gau}),
 \qquad
 \mathcal A_{\rm gau}(z):=(z,a_{\rm gau}z).
\end{equation}

The raw publisher
\begin{equation}\label{eq:gauint-publicador-crudo}
 p_{\rm gau}:\mathfrak M_{\rm gau}\longrightarrow
 \mathscr Y_{\rm gau}
\end{equation}
publishes jointly the identities
\[
 \begin{gathered}
 |I|=4,quad |E|=6,quad |\mathscr B|=8,quad sB=0,
 \quad W_c(q+Ba)=W_c(q),\\
 K_\ell J=JK_m,quad K_\ell^*J=JK_m^*,
 \quad T_m=K_m^*K_m,\\
 (\vartheta_L)_*\Omega_m^{(8)}=\Omega_m^{(8)},
 \quad \text{the four quadratic Fubini factorizations},\\
 D_m^{\rm HMT}=3Q_m,quad D_m^{\rm HMT}W_{c,m}=3W_{c,m},
 \quad \partial H_m+H_m\partial=I-\iota_m\epsilon_m.
 \end{gathered}
\]
Finally,
\begin{equation}\label{eq:gauint-graph-p}
 \mathcal C_{\rm gau}^{\rm HMT}:=\operatorname{Graph}(p_{\rm gau}),
 \qquad
 \mathcal P_{\rm gau}(M):=(M,p_{\rm gau}M).
\end{equation}
The complete chain is
\begin{equation}\label{eq:gauint-cadena-completa}
 \mathcal C_{\rm cont}^{\rm disc}\supset
 \mathcal C_{\rm cont}^{\rm gau}
 \xrightarrow{\operatorname{Res}_{\rm gau}}\mathcal G_{\rm gau}
 \xrightarrow{\operatorname{pr}_{X,{\rm gau}}}X_{\chi,{\rm gau}}
 \xrightarrow{\mathcal A_{\rm gau}}\mathfrak M_{\rm gau}
 \xrightarrow{\mathcal P_{\rm gau}}\mathcal C_{\rm gau}^{\rm HMT}.
\end{equation}
Each arrow has an inverse on its image given by a projection of
coordinates. No output replaces its generative history.

\section{Theorem of internal gauge realization}

\begin{theorem}[Complete and natural internal realization]
\label{thm:gauint-realizacion}
For every \(x\in\mathcal C_{\rm cont}^{\rm gau}\), the chain \eqref{eq:gauint-cadena-completa} constructs, from the discrete structure of the continuum, a natural macrostructure with four directions, six incidences, eight faces, four involutions, a projective measure, a complete counting kernel, its positive transfer, a common eight-face law, a nonconstant incidence witness, and a cellular terminal. The thirteen coordinates of \eqref{eq:gauint-d-trece} remain recoverable in the output.
\end{theorem}

\begin{proof}
Equations \eqref{eq:gauint-cuatro-direcciones}--\eqref{eq:gauint-ocho-caras} derive \(4/6/8\). The incidence sum \eqref{eq:gauint-incidencia-invariante} proves the invariance of \(W_c\), and the uniform count proves \eqref{eq:gauint-wc-momentos}. The complete extensions produce \(k_m\) by \eqref{eq:gauint-kernel-normalizado}; the explicit domain conditions prove the Markov property, stationarity, and the two naturality equalities. Hence \(T_m=K_m^*K_m\) is positive and natural. Applying Fubini once to \eqref{eq:gauint-ley-ocho-caras} yields the four quadratic forms without changing the measure. The uniform marginal centers \(W_{c,m}\), whence \eqref{eq:gauint-w-eigen}. The initial object of the nerve gives the additional degeneracy and the SDR identity \eqref{eq:gauint-sdr}. Finally, the three graph constructions retain their inputs as first coordinates, so the composition forgets none of the thirteen inputs.
\end{proof}

\section{Provenance and Forgers}

\paragraph{Origin.}
The TPK seed, the oriented screen, the surviving extensions, the common measure, and the terminal contraction have the status \texttt{PREEXISTING\_AUTHORIAL\_ARCHITECTURE}. The counting transfer, the incidence mode, and the exact sector \(3Q\) are presented as a \texttt{RECOVERED\_RESULT}. The explicit derivation \(\mathbf P^1(\mathbf F_3)\to4/6/8\), the two lumpability conditions, and the non-ablative packaging by graphs have the status of a \texttt{NEW\_FORMALIZATION}; they formalize the preceding architecture without claiming a new conceptual origin.

\begin{ablation}[Internal falsifiers of the gauge branch]
\label{abl:gauint-falsadores}
The construction is falsified at the indicated level if any of the
following facts occurs.
\begin{enumerate}
 \item \(I_{\rm gau}\), \(E_{\rm gau}\), or
 \(\mathscr B_{\rm gau}\) is replaced by a list without the maps that derive them.
 \item \(sB=0\) fails or \(W_c\) ceases to be invariant and nonconstant.
 \item The count \(a_m(y,z)\) omits carry, orientation, memory, boundary, or
 route: in that case a different object has been counted.
 \item Projectivity, stationarity, or one of the two
 lumpabilities fails; then the corresponding naturality cannot be published.
 \item The eight faces do not follow from a unique \(\mu_m\) and a unique
 \(k_m\), or some \(\vartheta_L\) does not preserve \(\Omega_m^{(8)}\).
 \item \(P_{\Omega,m}W_{c,m}\ne0\) o
 \(D_m^{\rm HMT}W_{c,m}\ne3W_{c,m}\).
 \item The SDR identity \(\partial H+H\partial=I-\iota\epsilon\) fails.
 \item One of the thirteen coordinates is omitted: \emph{object replaced;
 conclusion not transferable}.
\end{enumerate}
These falsifiers control membership and naturality; they do not degrade a
history that satisfies all the identities.
\end{ablation}
\fi

\endgroup

The same rule defines its action on an elementary emission
\(\alpha:x\to y\). Since \(U_\alpha\) transports the ordered pair \emph{}
of sheets, its label map is
\(\overline U_\alpha e_{\Sigma,x}=e_{\Sigma,y}\) and
\(\overline U_\alpha e_{\Pi,x}=e_{\Pi,y}\). Its projectivization
\(\alpha_{\mathscr I}[v]:=[\overline U_\alpha v]\) acts in
\(\mathbf P^1(V_x)\) and, therefore,
\begin{equation}
\label{eq:pdfv2-cor-cell-accion-emision}
 \begin{gathered}
 \operatorname{Cell}_{\rm TPK}(\alpha)v_\star(x)=v_\star(y),\qquad
 \operatorname{Cell}_{\rm TPK}(\alpha)v_L(x)
   =v_{\alpha_{\mathscr I}L}(y),\\
 \operatorname{Cell}_{\rm TPK}(\alpha)v_{\{L,L'\}}(x)
   =v_{\{\alpha_{\mathscr I}L,\alpha_{\mathscr I}L'\}}(y),\qquad
 \operatorname{Cell}_{\rm TPK}(\alpha)v_{L,\epsilon}(x)
   =v_{\alpha_{\mathscr I}L,\epsilon}(y),\\
 \operatorname{Cell}_{\rm TPK}(\alpha)u_L(x)
   =u_{\alpha_{\mathscr I}L}(y),\qquad
 \operatorname{Cell}_{\rm TPK}(\alpha)u_{L,\epsilon}(x)
   =u_{\alpha_{\mathscr I}L,\epsilon}(y),\\
 \operatorname{Cell}_{\rm TPK}(\alpha)u_{L'\mid L}(x)
   =u_{\alpha_{\mathscr I}L'\mid\alpha_{\mathscr I}L}(y),\qquad
 \operatorname{Cell}_{\rm TPK}(\alpha)\omega_{\{L,L'\}}(x)
   =\omega_{\{\alpha_{\mathscr I}L,
                 \alpha_{\mathscr I}L'\}}(y).
 \end{gathered}
\end{equation}
In the leaf fiber and on its screen, one also has
\begin{equation}
\label{eq:pdfv2-cor-cell-entrelaza-proyectores-pantalla}
 \overline U_\alpha P_L=P_{\alpha_{\mathscr I}L}\overline U_\alpha,
 \qquad
 \operatorname{Cell}_{\rm TPK}(\alpha)\vartheta_L
 =\vartheta_{\alpha_{\mathscr I}L}
  \operatorname{Cell}_{\rm TPK}(\alpha).
\end{equation}
Boundary formulas give
\begin{equation}
\label{eq:pdfv2-cor-cell-accion-emision-natural}
 \partial\operatorname{Cell}_{\rm TPK}(\alpha)
 =\operatorname{Cell}_{\rm TPK}(\alpha)\partial,\qquad
 \operatorname{Cell}_{\rm TPK}(\beta\alpha)
 =\operatorname{Cell}_{\rm TPK}(\beta)
  \operatorname{Cell}_{\rm TPK}(\alpha).
\end{equation}
Thus the map of complexes used in the terminal register is the free action
induced by the same emission, not an added symbol.

\subsection{Module of paths, totalization and nonadic return cone}
Let \(J_d(x)\) be the full category of depth-\(d\) prefixes of the
common tree, \(\operatorname{Term}_d(x)\) its terminal set, and
\(\Gamma_{d,N}(\nu):=C_{d,N}(\nu)\) the cellular complex of
\eqref{eq:pdfv2-cor-cono-celda}. The module, totalization, and cone that
follow simultaneously preserve all continuations; they do not select
a history or separate a cohomological construction from the common continuum.
\begingroup
  \let\section\subsection
  % Extracción quirúrgica del propietario V2 05d: líneas 286--512.
% El resto permanece en este archivo como procedencia, pero no se publica.
\iffalse
% Macroestructura de coherencia estrictamente interna del continuo discreto.
% No usa geometrías, clases ni conclusiones clásicas.

\chapter{Internal macrostructure of cells, coherence and completeness}
\label{chap:pdfv2-coh-interna}

\begin{statusbox}
The sole input of this chapter is the discrete structure of the continuum. The
branch \(\mathrm{coh}\) is constructed on the rings
\(A_N=\mathbb Z/3^N\mathbb Z\), the APP cells, their histories, and their
terminals. Its nonablative chain is
\[
 \mathcal C_{\rm cont}^{\rm disc}
 \supset\mathcal C_{\rm cont}^{\rm coh}
 \xrightarrow{\operatorname{Res}_{\rm coh}}
 \mathcal G_{\rm coh}
 \xrightarrow{\operatorname{pr}_{X,{\rm coh}}}
 X_{\chi,{\rm coh}}
 \xrightarrow{\mathcal A_{\rm coh}}
 \mathfrak M_{\rm coh}
 \xrightarrow{\mathcal P_{\rm coh}}
 \mathcal C_{\rm coh}^{\rm HMT}.
\]
No classical object appears as an input or as an output. Each image literally preserves its generating history.
\end{statusbox}

\section{Typed domain and index separation}
\label{sec:pdfv2-coh-domain}

To avoid identifying depth with precision, the directed set is used.
\begin{equation}\label{eq:pdfv2-coh-lambda}
 \Lambda:=\mathbb N_{\rm prof}\times\mathbb N_{\rm prec},
 \qquad
 \lambda=(d,N),
 \qquad
 A_N:=\mathbb Z/3^N\mathbb Z.
\end{equation}
If \(\lambda\leq\lambda'=(d',N')\), the transition
\begin{equation}\label{eq:pdfv2-coh-rho-lambda}
 \rho_{\lambda',\lambda}:
 \mathcal C_{{\rm cont},\lambda'}^{\rm disc}
 \longrightarrow
 \mathcal C_{{\rm cont},\lambda}^{\rm disc}
\end{equation}
truncates the history after depth \(d\) and reduces coefficients by
\(A_{N'}\twoheadrightarrow A_N\), preserving sheet, orientation, residue,
quotient, carry, memory, boundary, path, survival, and terminal.

The substructure \(\mathcal C_{\rm cont}^{\rm coh}\) consists of the histories
\(x\in\mathcal C_{\rm cont}^{\rm disc}\) that satisfy, at every level:
\begin{enumerate}
 \item the prefix of \(x\) carries the oriented and based cells defined in
       the following section;
 \item each refinement induce a map of the complex of the celda;
 \item truncation, coefficient reduction, orientation, and nonadic
       extension commute;
 \item the surviving-extension multisections are nonempty at every
       depth;
 \item the nonadic return preserves the terminal difference between the history and its continuation.
\end{enumerate}
These conditions define the domain of the branch. They are not attributed to an
arbitrary history without verifying them.

\section{APP cell over \texorpdfstring{\(A_N\)}{AN}}
\label{sec:pdfv2-coh-cell}

An oriented cell is
\[
 \nu=(h_\nu,I_\nu,J_\nu,\epsilon_\nu),
\]
where \(h_\nu\) is a finite history, and \(I_\nu,J_\nu\) are nonempty
ordered sets of ports,
\[
 a_\nu:=|I_\nu|\geq1,
 \qquad
 b_\nu:=|J_\nu|\geq1,
\]
and \(\epsilon_\nu\in\{+1,-1\}\) is its orientation.

We define
\begin{equation}\label{eq:pdfv2-coh-kappa}
 \begin{aligned}
 \kappa_{\nu,N}:
 A_N^{I_\nu}\oplus A_N^{J_\nu}
 &\longrightarrow A_N^{I_\nu\times J_\nu},\\
 \kappa_{\nu,N}(u,v)_{ij}&:=u_i-v_j.
 \end{aligned}
\end{equation}
Once the base ports \(i_0\in I_\nu\), \(j_0\in J_\nu\) have been chosen, let
\begin{equation}\label{eq:pdfv2-coh-delta}
 \begin{aligned}
 \delta_{\nu,N}:A_N^{I_\nu\times J_\nu}
 &\longrightarrow
 A_N^{(I_\nu\setminus i_0)\times(J_\nu\setminus j_0)},\\
 \delta_{\nu,N}(w)_{ij}
 &:={}
 w_{ij}-w_{i j_0}-w_{i_0j}+w_{i_0j_0}.
 \end{aligned}
\end{equation}

\begin{proposition}[Exact sequence of the cell]
\label{prop:pdfv2-coh-cell-exact}
For every celda \(\nu\) and every precision \(N\), the sucesion
\begin{equation}\label{eq:pdfv2-coh-exact-sequence}
 0\longrightarrow A_N
 \xrightarrow{\ \Delta\ }
 A_N^{I_\nu}\oplus A_N^{J_\nu}
 \xrightarrow{\ \kappa_{\nu,N}\ }
 A_N^{I_\nu\times J_\nu}
 \xrightarrow{\ \delta_{\nu,N}\ }
 A_N^{(a_\nu-1)(b_\nu-1)}
 \longrightarrow0,
\end{equation}
where
\[
 \Delta(t)=\bigl((t)_{i\in I_\nu},(t)_{j\in J_\nu}\bigr),
\]
is exact and split over \(A_N\).
\end{proposition}

\begin{proof}
If \(\kappa_{\nu,N}(u,v)=0\), then \(u_i=v_j\) for every \(i,j\), so \((u,v)\) is diagonal. Thus,
\(\ker\kappa_{\nu,N}=\operatorname{im}\Delta\). Direct substitution gives
\(\delta_{\nu,N}\kappa_{\nu,N}=0\).

If \(\delta_{\nu,N}(w)=0\), set
\[
 u_i=w_{i j_0},
 \qquad
 v_j=w_{i_0j_0}-w_{i_0j}.
\]
Then
\[
 u_i-v_j=w_{i j_0}+w_{i_0j}-w_{i_0j_0}=w_{ij},
\]
so that \(\ker\delta_{\nu,N}=\operatorname{im}\kappa_{\nu,N}\). To prove the surjectivity of \(\delta_{\nu,N}\), the base
row and column are set to zero and the desired coordinates are placed
arbitrarily in the remaining positions. Those same formulas provide
splittings.
\end{proof}

The complex cell, in degrees \(1\to0\), is the homological cone of
\(\kappa_{\nu,N}\) with the convention that places its source in degree \(1\):
\begin{equation}\label{eq:pdfv2-coh-cell-complex}
 C_{\nu,N}:=\operatorname{Cone}(\kappa_{\nu,N})
 :=\left[
 A_N^{I_\nu}\oplus A_N^{J_\nu}
 \xrightarrow{\ \kappa_{\nu,N}\ }
 A_N^{I_\nu\times J_\nu}
 \right].
\end{equation}
The preceding proposition gives
\begin{equation}\label{eq:pdfv2-coh-cell-homology}
 H_1(C_{\nu,N})\simeq A_N,
 \qquad
 H_0(C_{\nu,N})\simeq A_N^{(a_\nu-1)(b_\nu-1)}.
\end{equation}
The cell defect is thus constructed, not added as a cardinal without
map.

\section{Orientation and carry of the cell}
\label{sec:pdfv2-coh-orientation-carry}

Let
\[
 \tau_1(u,v)=(v,u),
 \qquad
 P(w)_{ji}=w_{ij}.
\]
Oriented exchange is the complex map
\begin{equation}\label{eq:pdfv2-coh-theta}
 \Theta_{\nu,N}:=(\tau_1,-P):
 C_{a_\nu,b_\nu,N}\longrightarrow C_{b_\nu,a_\nu,N},
\end{equation}
because
\begin{equation}\label{eq:pdfv2-coh-orientation}
 \kappa_{b_\nu,a_\nu,N}\tau_1
 =-P\kappa_{a_\nu,b_\nu,N}.
\end{equation}
Moreover,
\(\Theta_{\nu^{\rm op},N}\Theta_{\nu,N}=I\). The opposite route is not
identified with the original: it changes the degree-zero sign and preserves the act
of inversion.

For \(r\in A_N\), let
\(\langle r\rangle_N\in\{0,\ldots,3^N-1\}\) be its canonical representative. The
local borrow/carry is
\begin{equation}\label{eq:pdfv2-coh-beta}
 \beta_N(r,s):=
 \frac{\langle r-s\rangle_N-\langle r\rangle_N+\langle s\rangle_N}
 {3^N}\in\{0,1\}.
\end{equation}
Thus,
\begin{equation}\label{eq:pdfv2-coh-carry-cell}
 \langle u_i-v_j\rangle_N
 =\langle u_i\rangle_N-\langle v_j\rangle_N
  +3^N\beta_N(u_i,v_j).
\end{equation}
The matrix
\[
 \operatorname{Carr}_{\nu,N}(u,v)
 :=\bigl(\beta_N(u_i,v_j)\bigr)_{ij}
\]
is preserved alongside \(\kappa_{\nu,N}(u,v)\): residue and carry are two
distinct coordinates.

When port sizes follow from the two APP leaves, the admissibility
of the cell includes the integer identity
\begin{equation}\label{eq:pdfv2-coh-app-defect}
 (a_\nu-1)(b_\nu-1)
 =\rho_{\Pi,\nu}-\rho_{\Sigma,\nu}+1
  +9(c_{\Pi,\nu}-c_{\Sigma,\nu}).
\end{equation}
The equality is required only for a cell with that genealogy; it is not asserted
for integers devoid of APP leaves.

If \(L_N(u)\) is the lift whole canonica of the matrix of a flecha
\(u\), it define the carry of composition by
\begin{equation}\label{eq:pdfv2-coh-carry-composition}
 L_N(v)L_N(u)=L_N(vu)+3^N c_N(v,u).
\end{equation}
The associativity of three arrows gives the exact cocycle.
\begin{equation}\label{eq:pdfv2-coh-carry-cocycle}
 c_N(w,vu)+L_N(w)c_N(v,u)
 =c_N(wv,u)+c_N(w,v)L_N(u).
\end{equation}

\section{Category of histories and cellular functor}
\label{sec:pdfv2-coh-category}

For a prefix \(x_{\leq d}\), let \(J_d(x)\) be the finite category whose
objects are all its cells \(\nu\). Its arrows are generated by:
\begin{enumerate}
 \item refinements \(u:\nu\to\nu'\) with surjective maps of
       ports
       \[
        \alpha_u:I_{\nu'}\twoheadrightarrow I_\nu,
        \qquad
        \beta_u:J_{\nu'}\twoheadrightarrow J_\nu;
       \]
 \item the exchanges of orientation \(\Theta_\nu\);
 \item the nonadic extensions \(\Gamma_9\nu\);
 \item identity and composition, subject to the boundary,
       path, orientation, and carry relations of the history.
\end{enumerate}
Finiteness refers to the prefix; the system of all prefixes is not declared finite.

The cellular functor is
\begin{equation}\label{eq:pdfv2-coh-gamma-functor}
 \Gamma_{d,N}:J_d(x)\longrightarrow
 \operatorname{Ch}^{[0,1]}(A_N),
 \qquad
 \Gamma_{d,N}(\nu)=C_{\nu,N}.
\end{equation}
For a refinement arrow \(u:\nu\to\nu'\), it acts by means of
\begin{align}
 \Gamma_{d,N}(u)_1(p,q)
 &:=(p\circ\alpha_u,q\circ\beta_u),
 \label{eq:pdfv2-coh-gamma-one}\\
 \Gamma_{d,N}(u)_0(w)
 &:=w\circ(\alpha_u\times\beta_u).
 \label{eq:pdfv2-coh-gamma-zero}
\end{align}
It is verified coordinate by coordinate.
\begin{equation}\label{eq:pdfv2-coh-gamma-chain-map}
 \kappa_{\nu',N}\Gamma_{d,N}(u)_1
 =\Gamma_{d,N}(u)_0\kappa_{\nu,N}.
\end{equation}
On an inversion of orientation it acts by
\(\Theta_{\nu,N}\), and on compositions it acts by composition, separately preserving
\(c_N(v,u)\).

The reduction \(r_{N+1,N}:A_{N+1}\twoheadrightarrow A_N\) satisfies
\begin{equation}\label{eq:pdfv2-coh-kappa-reduction}
 r_{N+1,N}^{I\times J}\kappa_{\nu,N+1}
 =\kappa_{\nu,N}
  \bigl(r_{N+1,N}^{I}\oplus r_{N+1,N}^{J}\bigr).
\end{equation}
Therefore, the cell diagram is naturally simultaneous in history and
precision.

\fi
\section{Module of paths, cellular totalization and nonadic terminal}
\label{sec:pdfv2-coh-route-module}

Let \(\operatorname{Term}_d(x)\) be the conjunto of terminal of memory of the
prefix. For \(\nu\in J_d(x)\), definase
\begin{equation}\label{eq:pdfv2-coh-w-module}
 W_{d,N}(\nu):=
 A_N\bigl[
 \operatorname{Path}_{J_d(x)}(\nu,\operatorname{Term}_d(x))
 \bigr].
\end{equation}
It is the free module on all terminal routes; it selects none. If
\(u:\nu\to\nu'\), the contravariant action is
\begin{equation}\label{eq:pdfv2-coh-w-action}
 W_{d,N}(u):W_{d,N}(\nu')\longrightarrow W_{d,N}(\nu),
 \qquad
 [\lambda]\longmapsto[\lambda\circ u].
\end{equation}
Each generator preserves orientation, carry, boundary, route, and memory.

\section{Bilateral bar complex and totalization of the bicolumn}
\label{sec:pdfv2-coh-bar}

For \(q\geq0\) and \(r\in\{0,1\}\), let
\begin{equation}\label{eq:pdfv2-coh-bar-bigraded}
 B_{q,r}^{d,N}(x):=
 \bigoplus_{\nu_0\xrightarrow{u_1}\cdots\xrightarrow{u_q}\nu_q}
 W_{d,N}(\nu_q)\otimes_{A_N}\Gamma_{d,N}(\nu_0)_r.
\end{equation}
The bar differential is defined without abbreviation by
\begin{align}
 d_{\rm bar}(w;u_1,\ldots,u_q;c)
 :={}&(w;u_2,\ldots,u_q;\Gamma(u_1)c)
 \nonumber\\
 &+\sum_{i=1}^{q-1}(-1)^i
 (w;u_1,\ldots,u_{i+1}u_i,\ldots,u_q;c)
 \nonumber\\
 &+(-1)^q
 (W(u_q)w;u_1,\ldots,u_{q-1};c).
 \label{eq:pdfv2-coh-bar-differential}
\end{align}
For \(q=1\), this is exactly the relation
\begin{equation}\label{eq:pdfv2-coh-coend-relation}
 d_{\rm bar}(w;u;c)
 =w\otimes\Gamma(u)c-W(u)w\otimes c.
\end{equation}
A transported incidence is thereby identified with the same incidence written
in the refined chart; the relation is not concealed beneath the name of the
totalizer.

The total complex is
\begin{equation}\label{eq:pdfv2-coh-total-complex}
 K_{d,N}(x):=\operatorname{Tot}B_{\bullet,\bullet}^{d,N}(x),
 \qquad
 D_{\rm tot}:=d_{\rm bar}+(-1)^q\partial_\Gamma
 \quad\text{on }B_{q,r}^{d,N}.
\end{equation}
The functoriality of \(W\) and \(\Gamma\), together with
\(\partial_\Gamma^2=0\), gives
\begin{equation}\label{eq:pdfv2-coh-total-square-zero}
 D_{\rm tot}^2=0.
\end{equation}
In compact notation, but now already discharged by
\eqref{eq:pdfv2-coh-bar-bigraded}--
\eqref{eq:pdfv2-coh-bar-differential},
\begin{equation}\label{eq:pdfv2-coh-derived-coend}
 K_{d,N}(x)
 \simeq
 W_{d,N}\otimes^{\mathbf L}_{A_N[J_d(x)]}\Gamma_{d,N}.
\end{equation}

\section{Depth and terminal nonadic colimit}
\label{sec:pdfv2-coh-terminal}

The prefixes form the colimit.
\begin{equation}\label{eq:pdfv2-coh-k-infinity}
 K_{\infty,N}(x):=\varinjlim_d K_{d,N}(x).
\end{equation}
The nine-phase extension carries depth \(d\) to depth \(d+9\).
Only after forming \eqref{eq:pdfv2-coh-k-infinity} does it induce a
well-typed endomorphism
\begin{equation}\label{eq:pdfv2-coh-u-gamma9}
 U_{\Gamma_9,N}:K_{\infty,N}(x)
 \longrightarrow K_{\infty,N}(x).
\end{equation}
The terminal is
\begin{equation}\label{eq:pdfv2-coh-terminal-cone}
 \mathcal T_{{\rm coh},N}(x)
 :=\operatorname{Cone}\bigl(I-U_{\Gamma_9,N}\bigr).
\end{equation}
With the homological convention
\[
 \operatorname{Cone}(f)_q
 =K_{\infty,N,q-1}\oplus K_{\infty,N,q},
\]
its differential is
\begin{equation}\label{eq:pdfv2-coh-cone-differential}
 \partial_{\rm Cone}(a,b)
 =\bigl(-\partial_Ka,\partial_Kb+f(a)\bigr).
\end{equation}

Although the visible phase of \(h\) and \(\Gamma_9h\) coincides, the generators
\([h]\) and \([\Gamma_9h]\) are different because their ledgers and depths
are different. The terminal object preserves
\begin{equation}\label{eq:pdfv2-coh-terminal-difference}
 (I-U_{\Gamma_9,N})[h]=[h]-[\Gamma_9h].
\end{equation}
Precision reductions commute with the terminal:
\begin{equation}\label{eq:pdfv2-coh-u-reduction}
 r_{N+1,N}U_{\Gamma_9,N+1}
 =U_{\Gamma_9,N}r_{N+1,N}.
\end{equation}

\section{Terminal cohomological reader: incidence cycle and class in \(H_0(\operatorname{Cone}(I-U_{\Gamma_9,N}))\)}
\label{sec:pdfv2-coh-reader}

Each cell of \(x_{\leq d}\) carries a vector of incidence
\[
 \omega_{\nu,N}(x)
 \in A_N^{I_\nu\times J_\nu}=C_{\nu,N,0}
\]
and the complete terminal route \(\lambda_{\nu,x}\). The raw cycle is
\begin{equation}\label{eq:pdfv2-coh-z-cycle}
 z_{d,N}(x):=
 \sum_{\nu\in\operatorname{Cell}_d(x)}
 \epsilon_\nu[\lambda_{\nu,x}]
 \otimes\omega_{\nu,N}(x)
 \in B_{0,0}^{d,N}(x).
\end{equation}
As the total complex is homologically non-negative,
\begin{equation}\label{eq:pdfv2-coh-z-closed}
 D_{\rm tot}z_{d,N}(x)=0.
\end{equation}
Both the cycle and its class are preserved.
\[
 [z_{d,N}(x)]\in H_0(K_{d,N}(x)).
\]
After including it in \(K_{\infty,N}\), its terminal class is
\[
 [(0,z_{d,N}(x))]
 \in H_0(\mathcal T_{{\rm coh},N}(x)).
\]
The complete reader is
\begin{equation}\label{eq:pdfv2-coh-reader-full}
 \mathcal L_{{\rm coh},d,N}(x)
 :=\bigl(
 z_{d,N}(x),[z_{d,N}(x)],
 U_{\Gamma_9,N}z_{d,N}(x),
 [(0,z_{d,N}(x))]
 \bigr).
\end{equation}
The publication is not merely a quotient class: the representative, the route that produced it, its continuation and the terminal remain.

\section{Surviving fibres and corrections without a selector}
\label{sec:pdfv2-coh-surviving-fibers}

Let \(\mathscr S_\lambda^{\rm coh}\) be the free \(A_N\)-module generated by the
admissible finite histories at level \(\lambda=(d,N)\). The reader extends
linearly to
\begin{equation}\label{eq:pdfv2-coh-e-lambda}
 e_\lambda:\mathscr S_\lambda^{\rm coh}
 \longrightarrow\mathscr O_\lambda^{\rm coh},
 \qquad
 \mathscr O_\lambda^{\rm coh}
 :=H_0(K_{d,N})\oplus H_0(\mathcal T_{{\rm coh},N}).
\end{equation}
The surviving codomain is
\begin{equation}\label{eq:pdfv2-coh-surviving-output}
 \mathscr O_\infty^{\rm surv}
 :=\operatorname{Im}\!\left(
 e_\infty:\mathcal C_{\rm cont}^{\rm coh}
 \longrightarrow\varprojlim_\lambda\mathscr O_\lambda^{\rm coh}
 \right).
\end{equation}
For \(a=(a_\lambda)_\lambda\in\mathscr O_\infty^{\rm surv}\), define
\begin{equation}\label{eq:pdfv2-coh-fiber}
 \mathscr F_\lambda(a):=
 \left\{s_\lambda\in\mathscr S_\lambda^{\rm coh}:
 \begin{array}{l}
 e_\lambda(s_\lambda)=a_\lambda,\\
 \exists y\in\mathcal C_{\rm cont}^{\rm coh}\text{ with }
 y|_\lambda=s_\lambda\text{ and }e_\infty(y)=a
 \end{array}
 \right\}.
\end{equation}
These fibers are nonempty by the definition of the surviving codomain. If
\(\lambda'\geq\lambda\), the restriction is surjective:
\begin{equation}\label{eq:pdfv2-coh-fiber-surjection}
 \rho_{\lambda',\lambda}:\mathscr F_{\lambda'}(a)
 \twoheadrightarrow\mathscr F_\lambda(a).
\end{equation}
Indeed, a complete history that witnesses
\(s_\lambda\in\mathscr F_\lambda(a)\) provides by truncation an
antecedent at any later level. That antecedent is not chosen as
part of the object.

The complete multisect is
\begin{equation}\label{eq:pdfv2-coh-extension-multisection}
 \Sigma_{\lambda',\lambda}^{\rm coh}(s_\lambda;a)
 :=\{s_{\lambda'}\in\mathscr F_{\lambda'}(a):
 \rho_{\lambda',\lambda}s_{\lambda'}=s_\lambda\}.
\end{equation}
For a provisional prolongation \(\widetilde s_{\lambda'}\), let
\begin{equation}\label{eq:pdfv2-coh-discrepancy}
 \delta_{\lambda'}:=
 a_{\lambda'}-e_{\lambda'}(\widetilde s_{\lambda'}).
\end{equation}
The relation of all corrections is
\begin{equation}\label{eq:pdfv2-coh-correction-relation}
 \operatorname{Corr}_{\lambda',\lambda}
 (\widetilde s_{\lambda'},\delta_{\lambda'})
 :=\left\{\eta\in\mathscr S_{\lambda'}^{\rm coh}:
 \begin{array}{l}
 \rho_{\lambda',\lambda}\eta=0,\\
 e_{\lambda'}(\eta)=\delta_{\lambda'},\\
 \eta\text{ is supported on new cells},\\
 \eta\text{ preserves orientation, carry, memory and terminal}
 \end{array}
 \right\}.
\end{equation}
If two lifts are comparable, their difference
\(\eta=s_{\lambda'}-\widetilde s_{\lambda'}\) belongs to this relation
exactly when it satisfies the four preceding conditions. It is not asserted
that \eqref{eq:pdfv2-coh-correction-relation} is nonempty for an
arbitrary discrepancy: the governing non-emptiness is that of
\eqref{eq:pdfv2-coh-extension-multisection} on the surviving domain.

\iffalse
\section{The thirteen coordinates of the constraint}
\label{sec:pdfv2-coh-thirteen}

For \(x_\lambda\in\mathcal C_{{\rm cont},\lambda}^{\rm coh}\), let
\begin{equation}\label{eq:pdfv2-coh-d-package}
 \begin{aligned}
 D_{{\rm coh},\lambda}(x_\lambda):=\bigl(&
 \mathcal F_{{\rm coh},\lambda},
 \operatorname{Ext}_{\Gamma,{\rm coh},\lambda},
 \operatorname{Rel}_{{\rm coh},\lambda},
 \mathcal N_{{\rm coh},\lambda},
 \operatorname{or}_{{\rm coh},\lambda},
 \operatorname{Carr}_{{\rm coh},\lambda},
 \operatorname{Mem}_{{\rm coh},\lambda},\\
 &\partial_{{\rm coh},\lambda},
 \gamma_{{\rm coh},\lambda},
 \operatorname{Msect}_{{\rm coh},\lambda},
 \operatorname{Surv}_{{\rm coh},\lambda},
 \operatorname{Term}_{{\rm coh},\lambda},
 \mathcal L_{{\rm coh},\lambda}\bigr).
 \end{aligned}
\end{equation}
Their coordinates are:
\begin{enumerate}
 \item \(\mathcal F_{{\rm coh},\lambda}\): all the modulos
       \(A_N^{I_\nu}\), \(A_N^{J_\nu}\),
       \(A_N^{I_\nu\times J_\nu}\), the complejos \(C_{\nu,N}\), the
       modulos \(W_{d,N}(\nu)\) and \(\mathscr S_\lambda^{\rm coh}\).
 \item \(\operatorname{Ext}_{\Gamma,{\rm coh},\lambda}\): all maps
       \(\Gamma(u)\), coefficient reductions, truncations, extensions
       \(\Gamma_9\), and the relations \(\Sigma_{\lambda',\lambda}\).
 \item \(\operatorname{Rel}_{{\rm coh},\lambda}\):
       \(\kappa\), \(\delta\), the exact sequence, the bar faces, the
       coend relation, the orientation, and the composition laws.
 \item \(\mathcal N_{{\rm coh},\lambda}\):
       \[
       (d,N,3^N,a_\nu,b_\nu,(a_\nu-1)(b_\nu-1),
       \rho_{\Sigma,\nu},\rho_{\Pi,\nu},c_{\Sigma,\nu},c_{\Pi,\nu}).
       \]
 \item \(\operatorname{or}_{{\rm coh},\lambda}\): the signs
       \(\epsilon_\nu\), the port orderings, and the maps
       \(\Theta_{\nu,N}\).
 \item \(\operatorname{Carr}_{{\rm coh},\lambda}\): the matrices
       \(\beta_N\), the cocycles \(c_N(v,u)\), residues, and quotients,
       conserved separately.
 \item \(\operatorname{Mem}_{{\rm coh},\lambda}\): the inherited ledger and the
       ordered word of refinements, carries, corrections, and advances
       \(\Gamma_9\), not merely its aggregate value.
 \item \(\partial_{{\rm coh},\lambda}\):
       \(\kappa_{\nu,N}\), \(D_{\rm tot}\),
       \(\partial_{\rm Cone}\) and the boundaries of the cells.
 \item \(\gamma_{{\rm coh},\lambda}\): all composable chains
       \(\nu_0\to\cdots\to\nu_q\) and their terminal continuations.
 \item \(\operatorname{Msect}_{{\rm coh},\lambda}\): all cells,
       terminal routes, extensions \(\Sigma\), and corrections
       \(\operatorname{Corr}\), without a selector.
 \item \(\operatorname{Surv}_{{\rm coh},\lambda}\): the system of fibres
       \(\mathscr F_\lambda(a)\) and surjective restrictions.
 \item \(\operatorname{Term}_{{\rm coh},\lambda}\):
       \(\operatorname{Cone}(I-U_{\Gamma_9,N})\), together with \(z\), \(Uz\)
       and their difference.
 \item \(\mathcal L_{{\rm coh},\lambda}\):
       \((z,[z],Uz,[(0,z)])\) and all its reduction compatibilities.
\end{enumerate}

For \(\lambda'\geq\lambda\), the coordinated transport satisfies
\begin{equation}\label{eq:pdfv2-coh-d-naturality}
 u_{\lambda',\lambda}^{\rm coh}
 D_{{\rm coh},\lambda'}(x_{\lambda'})
 =D_{{\rm coh},\lambda}
  (\rho_{\lambda',\lambda}x_{\lambda'}).
\end{equation}

\section{Constraint graph and dependent object}
\label{sec:pdfv2-coh-graphs}

The total restriction is the graph
\begin{equation}\label{eq:pdfv2-coh-g-graph}
 \mathcal G_{\rm coh}:=\operatorname{Graph}(D_{\rm coh}),
 \qquad
 \operatorname{Res}_{\rm coh}(x):=(x,D_{\rm coh}x).
\end{equation}
Thus,
\begin{equation}\label{eq:pdfv2-coh-res-inverse}
 \pi_1\operatorname{Res}_{\rm coh}=I,
 \qquad
 \operatorname{Res}_{\rm coh}\pi_1=I_{\mathcal G_{\rm coh}}.
\end{equation}

The full extractor is
\begin{equation}\label{eq:pdfv2-coh-chi}
 \chi_{\rm coh}(x):=
 \bigl(
 \{C_{\nu,N}\},
 \{z_\lambda,[z_\lambda]\},
 \{\mathscr F_\lambda(e_\infty x)\},
 \{\Sigma_{\lambda',\lambda}\},
 \{\operatorname{Corr}_{\lambda',\lambda}\}
 \bigr).
\end{equation}
We define
\begin{equation}\label{eq:pdfv2-coh-xchi}
 X_{\chi,{\rm coh}}
 :=\{(\chi_{\rm coh}(x),x,D_{\rm coh}x):
 x\in\mathcal C_{\rm cont}^{\rm coh}\}
\end{equation}
and
\begin{equation}\label{eq:pdfv2-coh-prx}
 \operatorname{pr}_{X,{\rm coh}}(x,D_{\rm coh}x)
 :=(\chi_{\rm coh}(x),x,D_{\rm coh}x).
\end{equation}
The projection onto \((x,D_{\rm coh}x)\) is inverse on this object
dependent.

\section{Internal assembler and publisher}
\label{sec:pdfv2-coh-assembler-publisher}

The ambient space \(\mathscr M_{\rm coh}^{\rm amb}\) consists of systems
compatible of exact sequences of cells, categories, functors, modules of
path, totalizers, terminals, readers, fibers and multisections. The
assembler raw is
\begin{equation}\label{eq:pdfv2-coh-assembler}
 \begin{aligned}
 a_{\rm coh}:X_{\chi,{\rm coh}}&\longrightarrow
 \mathscr M_{\rm coh}^{\rm amb},\\
 a_{\rm coh}(\chi(x),x,Dx)
 &:={}
 \bigl(
 \{C_{\nu,N},\kappa_{\nu,N},\delta_{\nu,N}\},
 \{J_d,\Gamma,W,K\},\\
 &\hspace{31mm}
 \{U_{\Gamma_9,N},\mathcal T_N\},
 \{e_\lambda,\mathscr F_\lambda,\Sigma,\operatorname{Corr}\}
 \bigr)_x.
 \end{aligned}
\end{equation}
The macrostructure is another graph:
\begin{equation}\label{eq:pdfv2-coh-macro-graph}
 \mathfrak M_{\rm coh}:=\operatorname{Graph}(a_{\rm coh}),
 \qquad
 \mathcal A_{\rm coh}(z):=(z,a_{\rm coh}z).
\end{equation}

Let \(\mathscr Y_{\rm coh}^{\rm int}\) be the type of certificates containing
\begin{equation}\label{eq:pdfv2-coh-certificate-list}
 \begin{gathered}
 \text{la sucesión exacta \eqref{eq:pdfv2-coh-exact-sequence}},\\
 \kappa_{b,a}\tau=-P\kappa_{a,b},
 \qquad D_{\rm tot}^2=0,
 \qquad\partial_{\rm Cone}^2=0,\\
 rU_{\Gamma_9}=U_{\Gamma_9}r,
 \qquad
 e_\lambda\rho_{\lambda',\lambda}
 =r_{\lambda',\lambda}e_{\lambda'},\\
 \rho_{\lambda',\lambda}:\mathscr F_{\lambda'}(a)
 \twoheadrightarrow\mathscr F_\lambda(a),
 \qquad
 \Sigma_{\lambda',\lambda}(s_\lambda;a)\neq\varnothing,
 \end{gathered}
\end{equation}
together with the raw cycles, their classes and their terminals. The publisher
\begin{equation}\label{eq:pdfv2-coh-publisher}
 p_{\rm coh}^{\rm raw}:\mathfrak M_{\rm coh}
 \longrightarrow\mathscr Y_{\rm coh}^{\rm int}
\end{equation}
produces that certificate without eliminating the macrostructure. Finally,
\begin{equation}\label{eq:pdfv2-coh-output-graph}
 \mathcal C_{\rm coh}^{\rm HMT}
 :=\operatorname{Graph}(p_{\rm coh}^{\rm raw}),
 \qquad
 \mathcal P_{\rm coh}(m):=(m,p_{\rm coh}^{\rm raw}m).
\end{equation}

\section{Internal Theorem of Coherent Generation}
\label{sec:pdfv2-coh-theorem}

\begin{theorem}[Internal Generation of the Coherence Macrostructure]
\label{thm:pdfv2-coh-internal-generation}
On the explicit domain \(\mathcal C_{\rm cont}^{\rm coh}\), the chain
\[
 \mathcal C_{\rm cont}^{\rm coh}
 \xrightarrow{\operatorname{Res}_{\rm coh}}
 \mathcal G_{\rm coh}
 \xrightarrow{\operatorname{pr}_{X,{\rm coh}}}
 X_{\chi,{\rm coh}}
 \xrightarrow{\mathcal A_{\rm coh}}
 \mathfrak M_{\rm coh}
 \xrightarrow{\mathcal P_{\rm coh}}
 \mathcal C_{\rm coh}^{\rm HMT}
\]
naturally constructs the cells, their cones, the diagram of histories, the
two-sided bar, the totalization, the nonadic terminal, and the
surviving fibers. It conserves the thirteen coordinates and all corrections as
multisections. Each stage admits, on its image, a projection that recovers
the preceding stage.
\end{theorem}

\begin{proof}
The \cref{prop:pdfv2-coh-cell-exact} constructs each cell over \(A_N\), with
diagonal kernel and defect \((a_\nu-1)(b_\nu-1)\). The identity
\eqref{eq:pdfv2-coh-orientation} makes the oriented exchange an involutive map of
complexes. The formulas
\eqref{eq:pdfv2-coh-gamma-one}--\eqref{eq:pdfv2-coh-gamma-zero} prove that
each refinement is a map of complexes, whereas
\eqref{eq:pdfv2-coh-kappa-reduction} proves compatibility in precision.

Precomposition of paths makes \(W\). The identities
functorial of \(W\) and \(\Gamma\) imply
\(d_{\rm bar}^2=0\); its commutation with the cellular differential gives
\(D_{\rm tot}^2=0\). Thus, \(K_{d,N}\) is constructed by the bar formulae,
not by to nominal invocation of the coend.

The extension \(\Gamma_9\) increases the depth. Therefore it induces the
endomorphism \(U_{\Gamma_9,N}\) only after passage to the colimit
\(K_{\infty,N}\). The cone
\(\operatorname{Cone}(I-U_{\Gamma_9,N})\) then conserves the difference
\([h]-[\Gamma_9h]\) and does not identify phase repetition with repetition
of state.

The cycle \(z_{d,N}(x)\) is formed with all the cells, signs, routes, and incidences of the history. Truncation and reduction commute with the reader. If \(s_\lambda\in\mathscr F_\lambda(a)\), the complete history appearing in the very definition of the fiber provides an antecedent at every level \(\lambda'\geq\lambda\); this proves the surjectivity \eqref{eq:pdfv2-coh-fiber-surjection}. The object preserves the entire set \(\Sigma_{\lambda',\lambda}\), not a choice. Differences between lifts are registered through \eqref{eq:pdfv2-coh-correction-relation}, with support, orientation, carry, memory, and terminal made explicit.

Finally, \(\mathcal G_{\rm coh}\), \(\mathfrak M_{\rm coh}\), and
\(\mathcal C_{\rm coh}^{\rm HMT}\) are graphs, whereas
\(X_{\chi,{\rm coh}}\) retains \(x\) and \(D_{\rm coh}x\) as coordinates.
The first projection of each object recovers its predecessor. Therefore, the
complete composition does not replace history with class, multisection with
selector, or terminal object with visible phase.
\end{proof}

\section{Probative origin and forgers}
\label{sec:pdfv2-coh-provenance-falsifiers}

\paragraph{Result recovered.}
Belonging to the already constructed architecture are the cell \(\kappa_{a,b}\), its
diagonal kernel and defect, the orientation
\(\kappa_{b,a}\tau=-P\kappa_{a,b}\), the APP carry, the cone of the cell, the
category of histories, the totalization by coend, the action \(\Gamma_9\), the
terminal \(\operatorname{Cone}(I-U_{\Gamma_9})\), and the architecture of
surviving extension.

\paragraph{New Formalization.}
Explicit formalizations include the exact sequence \eqref{eq:pdfv2-coh-exact-sequence} over \(A_N\), the separation between depth and precision, the complete faces of the bar, formation of the terminal object only after the colimit, packaging of the thirteen coordinates, corrections as a relation without a selector, and the four dependent objects or graphs that prevent loss of genealogy.

\begin{criterion}[Internal falsifiers of the branch \(\mathrm{coh}\)]
\label{crit:pdfv2-coh-falsifiers}
The construction fails if any of the following occurs:
\begin{enumerate}
 \item the sequence is not exact
       \eqref{eq:pdfv2-coh-exact-sequence};
 \item \(\kappa_{b,a}\tau=-P\kappa_{a,b}\) fails;
 \item some refinement does not satisfy
       \(\kappa_{\nu'}\Gamma(u)_1=\Gamma(u)_0\kappa_\nu\);
 \item the composition carry does not satisfy
       \eqref{eq:pdfv2-coh-carry-cocycle};
 \item \(D_{\rm tot}^2\neq0\);
 \item reduction, reader, or terminal do not commute with \(U_{\Gamma_9}\);
 \item \(\operatorname{Cone}(I-U_{\Gamma_9})\) is formed by previously identifying
       the depths \(d\) and \(d+9\);
 \item a fiber declared to be surviving is empty, or a restriction is not
       surjective;
 \item is replaced \(\Sigma_{\lambda',\lambda}\) by a single
       lift;
 \item only \([z]\) is published, and \(z\), route, carry, memory, or
       terminal is eliminated;
 \item any of the thirteen coordinates is omitted.
\end{enumerate}
In the last case, the binding diagnosis is: \emph{object replaced;
conclusion not transferable}.
\end{criterion}

\begin{warningbox}
This chapter concludes only the internal construction of
\(\mathcal C_{\rm coh}^{\rm HMT}\). It introduces neither a classical space,
nor a classical class, nor a classical publication. Any later intertwiner
belongs to another document and does not govern the existence of this
macrostructure.
\end{warningbox}
\fi

\endgroup

The local complex \eqref{eq:pdfv2-cor-complejo-local} and the vertex
\(y_\alpha\), endpoint of that same extension in the nerve, form the augmented
pullback. In it are defined
\begin{equation}
\label{eq:pdfv2-cor-zpm-h}
 \mathbf1_\alpha=(r,[y_\alpha]),\qquad
 z_-(\alpha)=(r,[y_\alpha]),
\end{equation}
\begin{equation}
\label{eq:pdfv2-cor-zplus-h}
 z_+(\alpha)=
 (r+3c_\alpha v_\alpha e+v_\alpha b,[y_\alpha]),
 \qquad
 h_*(\alpha)=(-v_\alpha f,0).
\end{equation}
The action on generators gives, without an additional abstract existence,
\begin{equation}
\label{eq:pdfv2-cor-filler}
 \partial z_-=\partial z_+=0,\qquad
 \partial h_*=z_--z_+.
\end{equation}
The correlative filling, the cell, and the collar preserve the same
\((\alpha,\operatorname{or}(\alpha),\kappa(\alpha),\partial\alpha,
\gamma\alpha,m(\gamma\alpha;m_0))\); they merely apply different internal readers
to the common record.

For the same finite based complex, the total determinant line is defined, without requiring balance or acyclicity.
\begin{equation}
\label{eq:pdfv2-cor-linea-determinante-total}
 \operatorname{Det}(C_{k,N}):=
 \bigotimes_q
 \left(\bigwedge^{\rm top}C_{k,N,q}\right)^{(-1)^q}.
\end{equation}
Refinement induces the determinantal morphism of the map of complexes and
simultaneously transports homology modules, Smith normal forms, and
Bockstein pages as fields of the register. Neither the existence of
\eqref{eq:pdfv2-cor-linea-determinante-total} nor its transport requires
homology to vanish. If the cone of refinement is contractible and oriented,
the morphism is canonically trivialized; in the general case the
cone is preserved and not replaced by a fictitious isomorphism.

\subsection{Smith Normal Form of the Terminal Operator and Primary Decomposition of Its Cokernel}
For each prefix \(y\), the unoriented carry and the oriented coefficient
of boundary of the common register admit the lifts
\begin{equation}
\label{eq:detint-c-v-lifts}
 \widetilde c_y,\widetilde v_y\in\mathbf Z,\qquad
 c_{y,n}:=\widetilde c_y\bmod3^n,\qquad
 v_{y,n}:=\widetilde v_y\bmod3^n.
\end{equation}
The determinantal output is neither an input nor a separate branch: it is the sector
produced by the same terminal reduction when the typed conditions
are verified
\begin{equation}
\label{eq:detint-dominio}
\begin{split}
 \mathcal C_{\rm cont}^{\rm det}:=\bigl\{x\in\mathcal C_{\rm cont}^{\rm disc}:{}&
 x\text{ survives at every depth};\\
 &\widehat P_x^{\rm red}\text{ is finite and based};\\
 &\chi(\widehat P_x^{\rm red})=0;\\
 &H_*(\widehat P_x^{\rm red}\otimes_{\mathbf Z_3}K_3)=0;\\
 &\text{the reduction produces a square block stable under refinement}
 \bigr\}.
\end{split}
\end{equation}
\begingroup
  \let\section\subsection
  % Extracción quirúrgica del propietario V2 05e: líneas 325--447.
% El resto permanece en este archivo como procedencia, pero no se publica.
\iffalse
\chapter[Internal Determinantal Macrostructure]{Internal Determinantal Macrostructure Generated by the Discrete Continuum}
\label{chap:detint-macroestructura}

\begin{thesisbox}
The \({\rm det}\) branch starts exclusively from surviving histories of
\(\mathcal C_{\rm cont}^{\rm disc}\). The nerve of extensions, the Alexander--Whitney diagonal,
the local complex, the unit, the two publications, the filler, the terminal
reduction, and the \(3\)-adic tower are constructed before any subsequent
evaluation. The balance and acyclicity conditions are incorporated into the
typed domain and are never tacitly assumed.
\end{thesisbox}

\section{Surviving determinant domain}

Let
\begin{equation}\label{eq:detint-an}
 A_n:=\mathbf Z/3^n\mathbf Z,
 \qquad
 \rho_{n+1,n}:A_{n+1}\longrightarrow A_n.
\end{equation}
For \(x\in\mathcal C_{\rm cont}^{\rm disc}\), let
\(\mathsf T_m(x)\) be the finite category of TPK prefixes of depth \(m\)
compatible with \(x\). Each object preserves sheet, orientation, residue,
quotient, integer carry, memory, boundary, route, multisection, survival,
and terminal; its morphisms are the admissible extensions that preserve those
coordinates.

From each object \(y\in\mathsf T_m(x)\), the internal integers are extracted
\begin{equation}\label{eq:detint-c-v-lifts}
 \widetilde c_y,\widetilde v_y\in\mathbf Z,
 \qquad
 c_{y,n}:=\widetilde c_y\bmod3^n,
 \qquad
 v_{y,n}:=\widetilde v_y\bmod3^n.
\end{equation}
The first coordinate is the unoriented carry and the second, the oriented
boundary coefficient. Reversal satisfies
\begin{equation}\label{eq:detint-or-cv}
 c_{\bar y,n}=c_{y,n},
 \qquad
 v_{\bar y,n}=-v_{y,n}.
\end{equation}

We write
\begin{equation}\label{eq:detint-z3-k3}
 \mathbf Z_3:=\varprojlim_n A_n,
 \qquad
 K_3:=\operatorname{Frac}(\mathbf Z_3).
\end{equation}
The determinantial restriction is defined by
\begin{equation}\label{eq:detint-dominio}
\begin{split}
 \mathcal C_{\rm cont}^{\rm det}:=
 \bigl\{x\in\mathcal C_{\rm cont}^{\rm disc}:{}&
 x\text{ survives at every depth};\\
 &\widehat P_x^{\rm red}\text{ is finite and based};\\
 &\chi(\widehat P_x^{\rm red})=0;\\
 &H_*(\widehat P_x^{\rm red}\otimes_{\mathbf Z_3}K_3)=0;\\
 &\text{the canonical reduction produces a square block}\\
 &\text{stable under refinement}\bigr\}.
\end{split}
\end{equation}
The objects \(\widehat P_x^{\rm red}\) and the square block will be constructed
in \cref{sec:detint-terminal}. Therefore
\eqref{eq:detint-dominio} is a verifiable membership condition and not
a conclusion inserted into the definition of the maps.

\section{Nerve, Alexander--Whitney, and counit}

The normalized complex is defined
\begin{equation}\label{eq:detint-gmn}
 G_{m,n}(x):=
 N_*^{\rm norm}\!\left(N\mathsf T_m(x);A_n\right).
\end{equation}
For a nondegenerate simplex
\(\sigma=[y_0\to y_1\to\cdots\to y_q]\), the diagonal is
\begin{equation}\label{eq:detint-aw}
 \Delta_{\rm AW}(\sigma)
 =\sum_{j=0}^q
 [y_0\to\cdots\to y_j]\otimes
 [y_j\to\cdots\to y_q].
\end{equation}
The counit is given by
\begin{equation}\label{eq:detint-counidad-g}
 \epsilon_G([y])=1,
 \qquad
 \epsilon_G(\sigma)=0\quad(q>0).
\end{equation}
In particular,
\begin{equation}\label{eq:detint-vertice-grouplike}
 \Delta_{\rm AW}[y]=[y]\otimes[y],
 \qquad
 \epsilon_G([y])=1.
\end{equation}

If \(m'\ge m\) and \(n'\ge n\), truncation and coefficient reduction
induce
\begin{equation}\label{eq:detint-rmn}
 R_{m',m}^{n',n}:G_{m',n'}(x)\longrightarrow G_{m,n}(x),
\end{equation}
and they are verified
\begin{equation}\label{eq:detint-aw-natural}
 R\partial=\partial R,
 \qquad
 (R\otimes R)\Delta_{\rm AW}=\Delta_{\rm AW}R,
 \qquad
 \epsilon_GR=\rho_{n',n}\epsilon_G.
\end{equation}
The path is not reduced to its endpoints: the entire simplex remains in
\(G_{m,n}(x)\).

\section{Internal local complex}

For \(c\in A_n\), define
\begin{equation}\label{eq:detint-p0-grados}
 P^0_{n,1}(c)=A_nf,
 \qquad
 P^0_{n,0}(c)=A_nr\oplus A_ne\oplus A_nb,
 \qquad
 P^0_{n,-1}(c)=A_ng,
\end{equation}
with differential
\begin{equation}\label{eq:detint-p0-diferencial}
 \partial f=3ce+b,
 \qquad
 \partial r=0,
 \qquad
 \partial e=g,
 \qquad
 \partial b=-3cg.
\end{equation}
The chain identity is literal:
\begin{equation}\label{eq:detint-d2}
 \partial^2f=3c\,\partial e+\partial b=3cg-3cg=0,
 \qquad
 \partial^2e=\partial^2b=\partial^2r=0.
\end{equation}

The increase
\begin{equation}\label{eq:detint-epsilon-p}
 \epsilon_P:P_n^0(c)\longrightarrow A_n[0]
\end{equation}
is fixed by
\[
 \epsilon_P(r)=1,
 \qquad
 \epsilon_P(e)=\epsilon_P(b)=\epsilon_P(f)=\epsilon_P(g)=0.
\]
The complex possesses the differential coalgebra.
\begin{equation}\label{eq:detint-coalgebra-p}
 \Delta_P(r)=r\otimes r,
 \qquad
 \Delta_P(q)=q\otimes r+r\otimes q
 \quad(q\in\{e,b,f,g\}).
\end{equation}
Thus, \(r\) is group-like and the other four generators are primitive
relative to \(r\). Equations
\eqref{eq:detint-p0-diferencial} prove that \(\Delta_P\) commutes with the
tensor differential.

\section{Local extension system and memory ledger}

For an extension \(\alpha:y\to y'\), we define
\begin{equation}\label{eq:detint-psi}
 \Psi_\alpha:P_n^0(c_y)\longrightarrow P_n^0(c_{y'})
\end{equation}
through means of
\begin{equation}\label{eq:detint-psi-generadores}
\begin{gathered}
 \Psi_\alpha(r_y)=r_{y'},\quad
 \Psi_\alpha(e_y)=e_{y'},\quad
 \Psi_\alpha(g_y)=g_{y'},\quad
 \Psi_\alpha(f_y)=f_{y'},\\
 \Psi_\alpha(b_y)=b_{y'}+3(c_{y'}-c_y)e_{y'}.
\end{gathered}
\end{equation}
It is a map of complexes because
\begin{equation}\label{eq:detint-psi-b-chain}
 \partial\Psi_\alpha(b_y)
 =-3c_{y'}g+3(c_{y'}-c_y)g
 =-3c_yg
 =\Psi_\alpha(\partial b_y),
\end{equation}
and
\begin{equation}\label{eq:detint-psi-f-chain}
 \Psi_\alpha(\partial f_y)
 =3c_ye+b+3(c_{y'}-c_y)e
 =3c_{y'}e+b
 =\partial\Psi_\alpha(f_y).
\end{equation}
Moreover,
\begin{equation}\label{eq:detint-psi-functorial}
 \Psi_\beta\Psi_\alpha=\Psi_{\beta\alpha}.
\end{equation}

The oriented coordinate change does not erase. It is recorded by
\begin{equation}\label{eq:detint-k-alpha}
 K_\alpha:=(v_{y'}-v_y)f_{y'}.
\end{equation}
For composable extensions,
\begin{equation}\label{eq:detint-k-coherencia}
 K_{\beta\alpha}=K_\beta+\Psi_\beta K_\alpha.
\end{equation}
This equality is the exact ledger of the boundary increment and makes visible
the memory that a purely nominal composition would lose.

The complete local diagram is the Grothendieck construction.
\begin{equation}\label{eq:detint-grothendieck-local}
 \int_{y\in\mathsf T_m(x)}P_n^0(c_y).
\end{equation}
In it, an arrow over \(\alpha:y\to y'\) carries \(\Psi_\alpha\)
as its linear component and \(K_\alpha\) as its coherence. When applied to an
element supported at a vertex, the notation is used
\begin{equation}\label{eq:detint-psi-soporte-vertice}
 \Psi_\alpha(p,[y]):=(\Psi_\alpha p,[y']).
\end{equation}
The edge \([y\to y']\) remains in the path coordinate of the nerve. This
convention will be used in the naturality formulas for
\(z_\pm\) and \(h_*\); no global map erasing the remaining
simplices of \(G_{m,n}(x)\) is asserted.

\section{Pullback, unit, root, publications and filler}

The pullback of complexes is
\begin{equation}\label{eq:detint-pgen}
 P^{\rm gen}_{y;m,n}
 :=P_n^0(c_y)\times_{A_n[0]}G_{m,n}(x)
 =\{(p,\gamma):\epsilon_P(p)=\epsilon_G(\gamma)\}.
\end{equation}
Its differential is
\begin{equation}\label{eq:detint-pgen-d}
 \partial(p,\gamma)=(\partial p,\partial\gamma).
\end{equation}
The common unit associated with the vertex \(y\) is
\begin{equation}\label{eq:detint-unidad}
 \mathbf1_y:=(r,[y]).
\end{equation}
Its two projections are group-like by
\eqref{eq:detint-vertice-grouplike} and
\eqref{eq:detint-coalgebra-p}; that is the precise meaning of a compatible
group-like unit of the pullback.

The root projection is the map of complexes
\begin{equation}\label{eq:detint-proyeccion-root}
 \varpi_{\rm root}:P^{\rm gen}_{y;m,n}\longrightarrow
 \operatorname{Root}_{y;m,n},
 \qquad
 \varpi_{\rm root}(p,\gamma):=(\epsilon_P(p)r,\gamma).
\end{equation}
The two internal publications are defined.
\begin{equation}\label{eq:detint-zpm}
 z_-(y):=(r,[y]),
 \qquad
 z_+(y):=(r+3c_yv_ye+v_yb,[y]).
\end{equation}
Both are cycles:
\begin{equation}\label{eq:detint-zpm-ciclos}
 \partial z_-(y)=0,
 \qquad
 \partial z_+(y)=3c_yv_yg-3c_yv_yg=0.
\end{equation}
Their roots coincide:
\begin{equation}\label{eq:detint-raiz-comun}
 \varpi_{\rm root}z_-(y)
 =(r,[y])
 =\varpi_{\rm root}z_+(y).
\end{equation}

The filler is constructed directly:
\begin{equation}\label{eq:detint-hstar}
 h_*(y):=(-v_yf,0)\in P^{\rm gen}_{y;m,n,1}.
\end{equation}
Then
\begin{equation}\label{eq:detint-filler-identidad}
 \partial h_*(y)
 =(-3c_yv_ye-v_yb,0)
 =z_-(y)-z_+(y).
\end{equation}
In particular,
\begin{equation}\label{eq:detint-clases-iguales}
 [z_-(y)]=[z_+(y)]
 \quad\text{in}\quad H_0(P^{\rm gen}_{y;m,n}),
\end{equation}
and the equality preserves its witness chain.

Under an extension \(\alpha:y\to y'\), we have
\begin{equation}\label{eq:detint-z-natural}
 z_-(y')=\Psi_\alpha z_-(y),
 \qquad
 z_+(y')-\Psi_\alpha z_+(y)=\partial K_\alpha,
\end{equation}
and
\begin{equation}\label{eq:detint-h-natural}
 h_*(y')-\Psi_\alpha h_*(y)=-K_\alpha.
\end{equation}
Therefore, naturality does not identify distinct values of \(v\): it records
their difference through the coherent homotopy \(K_\alpha\).

\section{Internal Orientation}

On \(P_n^0(c)\) is defined
\begin{equation}\label{eq:detint-omega-p}
 \Omega(r)=r,
 \qquad
 \Omega(q)=-q\quad(q\in\{e,b,f,g\}).
\end{equation}
It is an automorphism of complexes and of the relative coalgebra. In the nerve,
the oriented inversion is
\begin{equation}\label{eq:detint-omega-g}
 \mathfrak o_G[y_0\to\cdots\to y_q]
 :=(-1)^{q(q+1)/2}
 [\bar y_q\to\cdots\to\bar y_0].
\end{equation}
Together with \eqref{eq:detint-or-cv}, these actions give
\begin{equation}\label{eq:detint-or-z-h}
 \Omega z_-(y)=z_-(\bar y),
 \qquad
 \Omega z_+(c_y,v_y)=z_+(c_{\bar y},v_{\bar y}),
 \qquad
 \Omega h_*(y)=h_*(\bar y).
\end{equation}
Orientation acts on path, complex, publications, filler and determinant
line; it is not a label without an action.

\fi
\section{Lifting of the cokernel through the Bockstein tower and preservation of carry}
\label{sec:detint-terminal}

Only the common root is eliminated:
\begin{equation}\label{eq:detint-p-reducido}
 \overline P_{y;m,n}
 :=P^{\rm gen}_{y;m,n}/A_n\mathbf1_y,
 \qquad
 \widehat P_{y;m}^{\rm red}:=\varprojlim_n\overline P_{y;m,n}.
\end{equation}
The base is first ordered by the TPK order of routes and then by
\begin{equation}\label{eq:detint-orden-base}
 f\prec e\prec b\prec g.
\end{equation}
The following internal algorithm is executed.
\begin{enumerate}
 \item Find the coefficients that are units of \(\mathbf Z_3\).
 \item Choose the first according to \eqref{eq:detint-orden-base} and the order of
 paths.
 \item Cancel the corresponding pair and record the elementary
 operations, their sign, and their contribution to the orientation.
 \item Repeat until no cancellable unit pivot remains.
\end{enumerate}
The survival condition of \eqref{eq:detint-dominio} requires that,
cofinally in \(m\), the result be a two-term block
\begin{equation}\label{eq:detint-sx}
 \mathbf Z_3^{r_x}\xrightarrow{S_x}\mathbf Z_3^{r_x},
 \qquad
 \det S_x\ne0\text{ in }K_3,
\end{equation}
and that a refinement only add unit contractible blocks and recorded
changes of basis. If these conditions fail, the history does not belong to
\(\mathcal C_{\rm cont}^{\rm det}\); a nonexistent terminal is not forced.

The identity route provides an elementary witness. After removing the root,
the local complex is
\begin{equation}\label{eq:detint-testigo-contractil}
 A_nf\xrightarrow{\binom{3c}{1}}
 A_ne\oplus A_nb\xrightarrow{(1,-3c)}A_ng.
\end{equation}
Is contractible via
\begin{equation}\label{eq:detint-contraccion-local}
 s(g)=e,
 \qquad
 s(ae+\beta b)=\beta f,
 \qquad
 s(f)=0.
\end{equation}
Thus, the terminal conditions admit at least the internal identity witness
and do not describe a formally empty subdomain.

\section{Smith, \texorpdfstring{\(3\)}{3}-adic order and initial unit}

In oriented bases, a valuation diagonalization has the form
\begin{equation}\label{eq:detint-smith}
 U_xS_xV_x
 =\operatorname{diag}(3^{a_1}u_1,\ldots,3^{a_r}u_r),
 \qquad
 a_i\in\mathbf N,quad u_i\in\mathbf Z_3^\times.
\end{equation}
We define
\begin{equation}\label{eq:detint-orden-d}
 d_x:=v_3(\det S_x)=\sum_{i=1}^r a_i
\end{equation}
and the initial unit
\begin{equation}\label{eq:detint-leading-unit}
 \lambda_x:=3^{-d_x}\det S_x\in\mathbf Z_3^\times.
\end{equation}
The orientation coordinate trivializes the determinant line. If a
transition produces
\begin{equation}\label{eq:detint-estabilidad-bloques}
 S_x\oplus I_q=U S_{x'}V,
\end{equation}
the orientation is transported by the inverse factor of \(\det U\det V\).
The oriented unit \(\lambda_x^{\rm or}\) is then natural and does not change
when contractible blocks are added.

At finite precision,
\begin{equation}\label{eq:detint-smith-finito}
 S_{x,n}=S_x\bmod3^n,
 \qquad
 a_i^{(n)}=\min(a_i,n).
\end{equation}
The terminal is, therefore, the compatible family of all precisions
and not an isolated integer.

\section{Bockstein Tower and Higher Order Carry}

Let
\[
 C_x^1\xrightarrow{S_x}C_x^0
\]
the terminal block. If a class \([\bar y]\) admits a lift
\(y\in C_x^1\) with \(S_xy\in3^qC_x^0\), we define
\begin{equation}\label{eq:detint-beta-q}
 \beta_q[\bar y]
 :=\left[3^{-q}S_xy\bmod3\right].
\end{equation}
Changing the lift alters the representative by a boundary of the corresponding
page, so \(\beta_q\) is well defined. The higher
carry is
\begin{equation}\label{eq:detint-carry-superior}
 \operatorname{Carr}_q(y):=3^{-q}S_xy.
\end{equation}
It does not replace the TPK carry of \eqref{eq:detint-c-v-lifts}; it records its
prolongation in the coefficient tower.

For a block \(3^{a_i}u_i\), the class persists until page \(a_i\),
and on that page the nonzero differential is multiplication by
\(\bar u_i\in\mathbf F_3^\times\). Consequently,
\begin{equation}\label{eq:detint-bock-longitudes}
 \{\text{survival lengths}\}=\{a_1,\ldots,a_r\},
 \qquad
 d_x=\sum_i a_i.
\end{equation}
The oriented product of the page trivializations satisfies
\begin{equation}\label{eq:detint-bock-leading}
 \Theta_{\rm Bock}(x)=\overline{\lambda_x^{\rm or}}
 \in\mathbf F_3^\times.
\end{equation}
The proof is carried out in each diagonal block and transported by the recorded
unimodular changes. Smith, Bockstein and the initial unit are thus three
readings of the same terminal.

\iffalse
\section{The thirteen coordinates of the determinantal constraint}

For each \((m,n)\), it is defined
\begin{equation}\label{eq:detint-d-trece}
 D_{{\rm det},m,n}(x):=
 (\mathcal F,\operatorname{Ext}_\Gamma,\operatorname{Rel},
 \mathcal N,\operatorname{or},\operatorname{Carr},\operatorname{Mem},
 \partial,\gamma,\operatorname{Msect},\operatorname{Surv},
 \operatorname{Term},\mathcal L)_{{\rm det},m,n}(x),
\end{equation}
with the following content.
\begin{enumerate}
 \item The fiber is
 \[
  \mathcal F=(A_n,\mathsf T_m(x),
  \{P_n^0(c_y)\}_y,\{P^{\rm gen}_{y;m,n}\}_y).
 \]
 \item \(\operatorname{Ext}_\Gamma\) contains all morphisms
 \(\alpha\), the maps \(\Psi_\alpha\), the homotopies \(K_\alpha\) and
 the maps \(R_{m',m}^{n',n}\).
 \item \(\operatorname{Rel}\) contains \(\partial^2=0\), the
 pullback equation, AW, counit, \(\Psi_\beta\Psi_\alpha=\Psi_{\beta\alpha}\) and
 \(K_{\beta\alpha}=K_\beta+\Psi_\beta K_\alpha\).
 \item
 \(\mathcal N=(m,n,\widetilde c_y,\widetilde v_y,c_{y,n},v_{y,n},
 r_x,a_1,\ldots,a_{r_x},d_x)\).
 \item \(\operatorname{or}\) contains \(y\mapsto\bar y\),
 \(\mathfrak o_G\), \(\Omega\) and the orientation of the determinant
 line.
 \item \(\operatorname{Carr}\) separately preserves the integer carry,
 its reduction \(c_{y,n}\) and all higher carries
 \(\operatorname{Carr}_q\).
 \item \(\operatorname{Mem}\) contiene
 \[
  y_0\to\cdots\to y_m,
  \quad\{(\widetilde c_{y_j},\widetilde v_{y_j},\mathbf1_{y_j})\}_{j=0}^m,
  \quad\{K_{\alpha_j}\}.
 \]
 \item
 \(\partial=(\partial_G,\partial_P,\partial_{P^{\rm gen}},
 \partial_{\overline P},h_*)\).
 \item \(\gamma=[y_0\to\cdots\to y_q]\) preserves the complete oriented
 simplex.
 \item \(\operatorname{Msect}\) contains all compatible
 extensions, all \(3\)-adic lifts, and all unimodular
 presentations of the same oriented terminal; it does not choose one a posteriori.
 \item \(\operatorname{Surv}\) records arbitrary depth,
 coefficient compatibility, balance, acyclicity over \(K_3\),
 stability under unit blocks and persistence of the root.
 \item
 \[
  \operatorname{Term}=(\widehat P_x^{\rm red},S_x,
  \{a_i,u_i\},d_x,\lambda_x^{\rm or},
  \{\beta_q\},\Theta_{\rm Bock}).
 \]
 \item The internal reader is
 \[
  \mathcal L_{\rm det}^{\rm int}
  =(\mathbf1_y,
  \varpi_{\rm root}z_-=\varpi_{\rm root}z_+,
  [z_-]=[z_+],d_x,\lambda_x^{\rm or},\Theta_{\rm Bock}).
 \]
\end{enumerate}

The transitions satisfy, coordinate by coordinate,
\begin{equation}\label{eq:detint-naturalidad-trece}
 u^{\rm det}_{(m',n'),(m,n)}D_{{\rm det},m',n'}(x')
 =D_{{\rm det},m,n}(\tau_{m',m}x').
\end{equation}

\section{Assembler, publisher, and non-ablative chain}

We define
\begin{equation}\label{eq:detint-graph-d}
 \mathcal G_{\rm det}:=\operatorname{Graph}(D_{\rm det}),
 \qquad
 \operatorname{Res}_{\rm det}(x):=(x,D_{\rm det}x).
\end{equation}
Let \(\chi_{\rm det}(x)\) be the complete multisection of germs
\[
 (y,\widetilde c_y,\widetilde v_y,\operatorname{or}_y,\gamma_y)
\]
compatible at all depths. The dependent object is
\begin{equation}\label{eq:detint-x-dependiente}
 X_{\chi,{\rm det}}
 :=\{(\chi_{\rm det}(x),x,D_{\rm det}x):
 x\in\mathcal C_{\rm cont}^{\rm det}\},
\end{equation}
and
\begin{equation}\label{eq:detint-prx}
 \operatorname{pr}_{X,{\rm det}}(x,D_{\rm det}x)
 :=(\chi_{\rm det}(x),x,D_{\rm det}x).
\end{equation}

The raw assembler
\begin{equation}\label{eq:detint-ensamblador-crudo}
 a_{\rm det}:X_{\chi,{\rm det}}\longrightarrow
 \mathscr M_{\rm det}^{\rm amb}
\end{equation}
yields
\begin{equation}\label{eq:detint-ensamblador-salida}
\begin{split}
 a_{\rm det}(z)=\bigl(&
 \{G_{m,n},\Delta_{\rm AW},\epsilon_G\}_{m,n};\\
 &\{P_n^0(c),\Psi_\alpha,K_\alpha,P^{\rm gen},
 \mathbf1,z_-,z_+,h_*\}_{m,n};\\
 &\{\widehat P^{\rm red},S,d,\lambda^{\rm or},
 \beta_q,\Theta_{\rm Bock}\}\bigr).
\end{split}
\end{equation}
The macroobject preserves its entry:
\begin{equation}\label{eq:detint-graph-a}
 \mathfrak M_{\rm det}:=\operatorname{Graph}(a_{\rm det}),
 \qquad
 \mathcal A_{\rm det}(z):=(z,a_{\rm det}z).
\end{equation}

The raw publisher
\begin{equation}\label{eq:detint-publicador-crudo}
 p_{\rm det}:\mathfrak M_{\rm det}\longrightarrow\mathscr Y_{\rm det}
\end{equation}
publishes the certificate jointly
\begin{equation}\label{eq:detint-certificado-publicado}
\begin{gathered}
 \partial^2=0,quad \text{AW and counit},quad
 \mathbf1_y\text{ compatible group-like},\\
 \partial z_\pm=0,quad
 \varpi_{\rm root}z_-=\varpi_{\rm root}z_+,quad
 \partial h_*=z_--z_+,\\
 \Psi_\beta\Psi_\alpha=\Psi_{\beta\alpha},quad
 K_{\beta\alpha}=K_\beta+\Psi_\beta K_\alpha,\\
 d=v_3(\det S)=\sum_i a_i,quad
 \Theta_{\rm Bock}=\overline{\lambda^{\rm or}},quad
 \text{full naturality under refinement}.
\end{gathered}
\end{equation}
Finally,
\begin{equation}\label{eq:detint-graph-p}
 \mathcal C_{\rm det}^{\rm HMT}:=\operatorname{Graph}(p_{\rm det}),
 \qquad
 \mathcal P_{\rm det}(M):=(M,p_{\rm det}M).
\end{equation}
The complete chain is
\begin{equation}\label{eq:detint-cadena-completa}
 \mathcal C_{\rm cont}^{\rm disc}\supset
 \mathcal C_{\rm cont}^{\rm det}
 \xrightarrow{\operatorname{Res}_{\rm det}}\mathcal G_{\rm det}
 \xrightarrow{\operatorname{pr}_{X,{\rm det}}}X_{\chi,{\rm det}}
 \xrightarrow{\mathcal A_{\rm det}}\mathfrak M_{\rm det}
 \xrightarrow{\mathcal P_{\rm det}}\mathcal C_{\rm det}^{\rm HMT}.
\end{equation}
Each arrow has an inverse on its image given by a coordinate
projection. History is not replaced by value, multisection by selector,
complex by determinant or proof by name.

\section{Theorem of Internal Determinantal Realization}

\begin{theorem}[Complete internal determinant realization]
\label{thm:detint-realizacion}
For every \(x\in\mathcal C_{\rm cont}^{\rm det}\), the chain
\eqref{eq:detint-cadena-completa} constructs a macro-object natural in
depth and precision. It contains two cycles with a common root, an explicit filler
that identifies them, a compatible group-like unit, and a
\(3\)-adic terminal whose order, Smith form, and Bockstein tower produce the same
oriented reader. The thirteen coordinates of \eqref{eq:detint-d-trece}
remain recoverable in the output.
\end{theorem}

\begin{proof}
Equations \eqref{eq:detint-p0-diferencial}--
\eqref{eq:detint-d2} prove that \(P_n^0(c)\) is a complex. The two
augmentations are maps of complexes, so the pullback is typed.
Alexander--Whitney is natural and makes each vertex group-like; together with
\(r\), this gives the common unit. The direct calculation
\eqref{eq:detint-zpm-ciclos}--\eqref{eq:detint-filler-identidad} proves that
the publications are cycles, have the same root and are joined by
\(h_*\). The formulas for \(\Psi_\alpha\) prove functoriality, and
\(K_\alpha\) proves homotopic naturality while preserving the memory
increment. The definition of \(\mathcal C_{\rm cont}^{\rm det}\) ensures that
terminal reduction produces a square block without postulating it outside the
domain. Smith over \(\mathbf Z_3\) gives \(d=\sum_i a_i\); the calculation of each
block \(3^{a_i}u_i\) gives
\(\Theta_{\rm Bock}=\overline{\lambda^{\rm or}}\). Truncation, reduction and
addition of unit blocks preserve these identities. Finally, the
three graph constructions retain the input as the first coordinate,
so the complete composition is non-ablative.
\end{proof}

\section{Provenance and Forgers}

\paragraph{Origin.}
The TPK nerve, Alexander--Whitney, the counit, the common unit, the normal
form \(\partial f=3ce+b\), \(\partial e=g\),
\(\partial b=-3cg\), the common root, the filler, and the tower architecture
have the status of \texttt{PREEXISTING\_AUTHORIAL\_ARCHITECTURE}. The literal computation
of the two cycles, their root, and the identity
\(\partial h_*=z_--z_+\) has the status of a \texttt{RECOVERED\_RESULT}. The
local system \(\Psi_\alpha/K_\alpha\), the balanced surviving subdomain,
the exclusively \(3\)-adic terminal, and the packaging by
graphs have the status of a \texttt{NEW\_FORMALIZATION}; they formalize without altering
the conceptual provenance of the apparatus.

\begin{ablation}[Internal falsifiers of the determinantal branch]
\label{abl:detint-falsadores}
The construction is falsified at the corresponding level if any of
the following facts occurs.
\begin{enumerate}
 \item \(\partial^2=0\) fails; then the local complex does not exist.
 \item Alexander--Whitney does not commute with refinement, the counit fails, or
 the vertex ceases to be group-like.
 \item \(z_\pm\) are not cycles, their roots differ, or
 \(\partial h_*=z_--z_+\) fails.
 \item \(\Psi_\alpha\) is not a map of complexes or
 \(K_{\beta\alpha}=K_\beta+\Psi_\beta K_\alpha\) fails; in that case
 carry or memory has been lost.
 \item The reduced complex is not balanced or is not acyclic over
 \(K_3\); the history lies outside the typed terminal domain.
 \item The Bockstein lengths do not coincide with the \(a_i\), or
 \(\Theta_{\rm Bock}\ne\overline{\lambda^{\rm or}}\).
 \item Refinement modifies \(d\) or \(\lambda^{\rm or}\) outside
 unit blocks and recorded orientation changes.
 \item One of the thirteen coordinates is omitted: \emph{object replaced;
 conclusion not transferable}.
\end{enumerate}
These forgers check the internally constructed object; they do not replace it with a subsequent control.
\end{ablation}
\fi

\endgroup

\section{Naturalness of the ensemble and compatibility under refinement}

Before forming the horizon record, we make explicit the total action
of the five operations developed here on \emph{the same} edge
\(\mathfrak e_\alpha\). If \(P_x,Q_x\) denote the restrictions of
\eqref{eq:pdfv2-cor-sigma} to the fiber of \(x=s(\alpha)\), and analogously in
\(y=t(\alpha)\), we first set
\begin{equation}
\label{eq:pdfv2-cor-a-eta-emision}
 A_\alpha:=P_yU_\alpha P_x+Q_yU_\alpha Q_x,
 \qquad
 \eta_\alpha:=Q_yU_\alpha P_x+P_yU_\alpha Q_x.
\end{equation}
The five readings are dependent families indexed by the
\(\mathfrak e_\alpha\) itself. The index fixes source, target, route, and ledger, but
is not copied again within the value. Their types are the dependent
records
\(\operatorname{LedgerEntry}:=\operatorname{cod}(\lambda)\) and
\begin{align*}
 \mathsf Y_{\rm sp}[\alpha]&:=
 \operatorname{Rec}[U:\mathsf H_x\to\mathsf H_y;
 \sigma_x\in\operatorname{End}\mathsf H_x;
 \sigma_y\in\operatorname{End}\mathsf H_y;A,\eta:\mathsf H_x\to\mathsf H_y],\\
 \mathsf Y_{\rm sol}[\alpha]&:=
 \operatorname{Rec}[b,b^+,b^-\in\operatorname{LedgerEntry};
 \operatorname{upd}:\mathsf{Mem}_x\to\mathsf{Mem}_y],\\
 \mathsf Y_{\rm gau}[\alpha]&:=
 \mathcal P_{\rm fin}\operatorname{Rec}[
 \operatorname{Cell}(\alpha);q\in Q_{\rm inc}(y);B;W_c(q)],\\
 \mathsf Y_{\rm coh}[\alpha]&:=
 \operatorname{Rec}[I_y,J_y,\kappa_{y,N},\delta_y,C_{y,N};
 F_\alpha:C_{y,N}\to C_{x,N}],\\
 \mathsf Y_{\rm det}[\alpha]&:=
 \operatorname{Rec}[\Psi_\alpha,K_\alpha,\mathbf1_\alpha,z_\pm,h_*;
 \operatorname{Det},H_*,\operatorname{Smith},\operatorname{Boc}].
\end{align*}
Therefore, the common judgment is
\[
 \mathfrak e_\alpha:\operatorname{Ext}_k\ \vdash\
 \mathcal O_\bullet(\mathfrak e_\alpha)
 \in\mathsf Y_\bullet[\alpha],
 \qquad
 \bullet\in\{\mathrm{sp,sol,gau,coh,det}\}.
\]
\begin{align}
 \mathcal O_{\rm sp}(\mathfrak e_\alpha)
 &:=
 \bigl(U_\alpha,\sigma_x,\sigma_y,
 A_\alpha,\eta_\alpha\bigr),
 \label{eq:pdfv2-cor-accion-sp}\\
 \mathcal O_{\rm sol}(\mathfrak e_\alpha)
 &:=
 \bigl(b_\alpha,b_\alpha^+,b_\alpha^-;\,
 (M^+,M^-,D,H)\mapsto
 (M^++b_\alpha^+,M^-+b_\alpha^-,
  D+b_\alpha,H+|b_\alpha|)\bigr),
 \label{eq:pdfv2-cor-accion-sol}\\
 \mathcal O_{\rm gau}(\mathfrak e_\alpha)
 &:=
 \left\{
 \bigl(\operatorname{Cell}_{\rm TPK}(\alpha),
 q,B,W_c(q)\bigr):
 q\in Q_{\rm inc}(y)
 \right\},
 \label{eq:pdfv2-cor-accion-gau}\\
 \mathcal O_{\rm coh}(\mathfrak e_\alpha)
 &:=
 \bigl(I_y,J_y,\kappa_{y,N},\delta_y,C_{y,N},
 F_\alpha:C_{y,N}\longrightarrow C_{x,N}\bigr),
 \label{eq:pdfv2-cor-accion-coh}\\
 \mathcal O_{\rm det}(\mathfrak e_\alpha)
 &:=
 \bigl(\Psi_\alpha,K_\alpha,\mathbf1_\alpha,z_\pm(\alpha),
 h_*(\alpha),\operatorname{Det}(C_{y,N}),H_*(C_{y,N}),
 \operatorname{Smith}(C_{y,N}),\operatorname{Boc}(C_{y,N})\bigr).
 \label{eq:pdfv2-cor-accion-det}
\end{align}

The internal Archimedean reading is the first component of this simultaneous
action; it does not yet introduce any classical form:
\begin{equation}
\label{pII-full:eq:continuo-lector-arquimediano}
 \mathcal L_{\rm arq}^{\rm int}(\mathfrak e_\alpha)
 :=\mathcal O_{\rm sp}(\mathfrak e_\alpha).
\end{equation}
For every value in those families, including each element of the
multisection \(\mathcal O_{\rm gau}\), its support is the dependent index
\begin{equation}
\label{eq:pdfv2-cor-soporte-dependiente}
 \begin{aligned}
 \mathsf B_\alpha&:=
 \{(s(\alpha),t(\alpha),\gamma\alpha,
          \operatorname{Led}(\gamma\alpha))\},\\
 \operatorname{supp}_\alpha:\mathsf Y_\bullet[\alpha]&\longrightarrow
 \mathsf B_\alpha,\qquad
 o\longmapsto(s(\alpha),t(\alpha),\gamma\alpha,
          \operatorname{Led}(\gamma\alpha)).
 \end{aligned}
\end{equation}
Therefore, they remain typed, without a projection that recovers the edge,
\begin{equation}
\label{eq:pdfv2-cor-soporte-cinco-lecturas}
 \begin{gathered}
 s_\alpha(\mathcal O_\bullet):=s(\alpha),\qquad
 t_\alpha(\mathcal O_\bullet):=t(\alpha),\\
 \gamma_\alpha(\mathcal O_\bullet):=\gamma\alpha,
 \qquad
 \operatorname{Led}_\alpha(\mathcal O_\bullet)
 :=\operatorname{Led}(\gamma\alpha).
 \end{gathered}
\end{equation}
Here \(\operatorname{Cell}_{\rm TPK}(\alpha)\) is the map of complexes
\eqref{eq:pdfv2-cor-cell-accion-emision} induced by transport of the
same sheet fiber, and \(F_\alpha\) is
\eqref{eq:pdfv2-cor-refinamiento-celda} for that extension. In particular,
\(\mathcal O_{\rm gau}\) preserves the complete multisection of \(q\)'s; it does not
choose the origin or an orbit.

The five formulas above constitute a single dependent action.
\begin{equation}
\label{eq:pdfv2-cor-accion-unica-dependiente}
 \mathcal O(\mathfrak e_\alpha)
 :=\left\langle
 \mathcal O_{\rm sp},\mathcal O_{\rm sol},\mathcal O_{\rm gau},
 \mathcal O_{\rm coh},\mathcal O_{\rm det};
 \operatorname{Cpl}_\alpha
 \right\rangle,
\end{equation}
where \(\operatorname{Cpl}_\alpha\) is the certificate formed by the
identities
\begin{equation}
\label{eq:pdfv2-cor-acoplamientos-unicos}
 \begin{gathered}
 \operatorname{supp}_\alpha(\mathcal O_{\rm sp})
 =\operatorname{supp}_\alpha(\mathcal O_{\rm sol})
 =\operatorname{supp}_\alpha(\mathcal O_{\rm gau})
 =\operatorname{supp}_\alpha(\mathcal O_{\rm coh})
 =\operatorname{supp}_\alpha(\mathcal O_{\rm det}),\\
 \mathsf Q(t(\alpha))=\sigma_9\mathsf Q(s(\alpha)),\quad
 g(t(\alpha))=g(s(\alpha))+[1]_9,\quad
 q^{(9)}(t(\alpha))=q^{(9)}(s(\alpha))+\chi_9(g(s(\alpha))),\\
 A_\alpha^*A_\alpha+\eta_\alpha^*\eta_\alpha
 =P_xU_\alpha^*U_\alpha P_x+Q_xU_\alpha^*U_\alpha Q_x,\\
 (M^+,M^-,D,H)'=(M^++b_\alpha^+,M^-+b_\alpha^-,
 D+b_\alpha,H+|b_\alpha|),\\
 Q_{\rm inc}=C^2(\operatorname{Cell}_{\rm TPK};\mathbf F_3),\qquad
 sB=0,\qquad W_c(q+Ba)=W_c(q),\\
 C_{y,N}\text{ is the same complex in }
 \mathcal O_{\rm coh}\text{ and }\mathcal O_{\rm det},\qquad
 \partial h_*=z_--z_+,\qquad
 K_{\beta\alpha}=K_\beta+\Psi_\beta K_\alpha.
 \end{gathered}
\end{equation}
\Needspace{18\baselineskip}
The symbol \(\bullet\) ranges over the five readings of
\eqref{eq:pdfv2-cor-accion-unica-dependiente}; all preserve the same source, destination, route, and ledger.

We collect, without separating the state, the operations constructed up to here:
\begin{equation}
\label{eq:pdfv2-cor-constructor-unico}
 \begin{aligned}
 \mathfrak D_{k,N}(\widehat x)
 &:=\bigl(\{\mathfrak e_\alpha\}_\alpha;\\
 &\quad U_k,\sigma_k,P_k,Q_k,A_k,\eta_k;\\
 &\quad b_\alpha,b_\alpha^\pm,M_k^\pm,D_k^0,H_k^0,
 \kappa_k^{\rm can};\\
 &\quad V_{\rm TPK},\mathscr I,\mathscr E,\mathscr B,B;\\
 &\quad \operatorname{Cell}_{\rm TPK}(\widehat x),Q_{\rm inc}(\widehat x),W_c;\\
 &\quad I_k,J_k,\kappa_{k,N},\delta,C_{k,N};\\
 &\quad c_\alpha,v_\alpha,\Psi_\alpha,K_\alpha,
 \mathbf1_\alpha,z_\pm(\alpha),h_*(\alpha);\\
 &\quad
 \operatorname{Det}(C_{k,N}),H_*(C_{k,N}),
 \operatorname{Smith}(C_{k,N}),\operatorname{Boc}(C_{k,N});\\
 &\quad \{\mathcal O(\mathfrak e_\alpha)\}_\alpha
 \bigr).
 \end{aligned}
\end{equation}
The semicolon merely separates levels of interpretation of the same dependent tuple; it does not define independent structures. The input state is neither retained as the first coordinate nor recovered through a projection. Each operation is proved through its action on the generators \(\mathfrak e_\alpha\), which preserve the same support, the same route, and the same ledger.

For \(\ell\geq k\) and \(N'\geq N\), transport is defined component by
component through truncation of histories, reduction of coefficients, and
precomposition of ports. The naturality assertion that must be proved is
a single one:
\begin{equation}
\label{eq:pdfv2-cor-naturalidad-unica}
 u_{(\ell,N'),(k,N)}\mathfrak D_{\ell,N'}(\widehat x_\ell)
 =\mathfrak D_{k,N}(\tau_{\ell,k}\widehat x_\ell).
\end{equation}
Naturality is formulated over the generated history, not by pretending that an
isolated fine edge must remain an edge after truncating its two
endpoints. For
\(\gamma=(\varepsilon_1,\ldots,\varepsilon_\ell)\), we define
\[
 \mathsf E_{\leq\ell}(\gamma)
 :=(\mathfrak e_{\varepsilon_i})_{1\leq i\leq\ell},
 \qquad
 \mathcal O_{\leq\ell}(\gamma)
 :=(\mathcal O(\mathfrak e_{\varepsilon_i}))_{1\leq i\leq\ell}.
\]
The two transports are the explicit prefix maps.
\begin{align}
 u^{\mathfrak e}_{\ell,k}:
 \mathsf E_{\leq\ell}(\gamma)&\longrightarrow
 \mathsf E_{\leq k}(\gamma),&
 (\mathfrak e_{\varepsilon_i})_{i\leq\ell}
 &\longmapsto(\mathfrak e_{\varepsilon_i})_{i\leq k},
 \label{eq:pdfv2-cor-transporte-arista}\\
 u^{\mathcal O}_{\ell,k}:
 \mathcal O_{\leq\ell}(\gamma)&\longrightarrow
 \mathcal O_{\leq k}(\gamma),&
 (\mathcal O(\mathfrak e_{\varepsilon_i}))_{i\leq\ell}
 &\longmapsto
 (\mathcal O(\mathfrak e_{\varepsilon_i}))_{i\leq k}.
 \label{eq:pdfv2-cor-transporte-accion}
\end{align}
In each retained entry, the component \(\mathrm{sp}\) intertwines
\(U,\sigma,P,Q,A,\eta\); the component \(\mathrm{sol}\) applies the conditional
expectation to \(b,b^\pm,\kappa\); the component \(\mathrm{gau}\) applies
\(\operatorname{Cell}(\tau)\oplus\tau_Q\); the component \(\mathrm{coh}\)
precomposes the ports with \(F\); and the component \(\mathrm{det}\) applies
\(\Psi,K,H_*,\operatorname{Smith},\operatorname{Boc}\) to the same map of
complexes. The prefix acts simultaneously on those five components.
Its components already discharged are
\begin{equation}
\label{eq:pdfv2-cor-naturalidad-componentes}
 \begin{gathered}
 U_{m,k}=U_{m,\ell}U_{\ell,k}\quad(m\geq\ell\geq k),\qquad
 E_kb_f^\pm=b_c^\pm+\kappa_k^{\rm can},\\
 r_{N',N}\kappa_{k,N'}=
 \kappa_{k,N}(r_{N',N}\oplus r_{N',N}),\\
 \Psi_\beta\Psi_\alpha=\Psi_{\beta\alpha},
 \qquad K_{\beta\alpha}=K_\beta+\Psi_\beta K_\alpha,\\
 u^{\mathcal O}_{\ell,k}\mathcal O_{\leq\ell}(\gamma)
 =\mathcal O_{\leq k}(\gamma)
 =\operatorname{map}(\mathcal O)\bigl(u^{\mathfrak e}_{\ell,k}
                    \mathsf E_{\leq\ell}(\gamma)\bigr).
 \end{gathered}
\end{equation}
The first equality uses multiplication of the conditional weights over the same
system of common cylinders; the remaining ones are the preceding explicit
formulas. No visible projection
can be used to reconstruct or select \(\widehat x\).

To complete the last component, let us denote by
\(F_{(\ell,N'),(k,N)}:C_{\ell,N'}\to C_{k,N}\) the map of complexes of
\eqref{eq:pdfv2-cor-refinamiento-celda}. Its cone is preserved as part of the
transition. Functoriality of homology and of the Bockstein sequence gives
\begin{equation}
\label{eq:pdfv2-cor-naturalidad-homologia-bockstein}
 H_*(F_{m,k})=H_*(F_{\ell,k})H_*(F_{m,\ell}),
 \qquad
 \operatorname{Boc}(F_{m,k})
 =\operatorname{Boc}(F_{\ell,k})
  \operatorname{Boc}(F_{m,\ell}).
\end{equation}
The determinant line is transported by the identity of the exact triangle.
\begin{equation}
\label{eq:pdfv2-cor-naturalidad-determinante}
 \operatorname{Det}(C_{k,N})
 \simeq
 \operatorname{Det}(C_{\ell,N'})\otimes
 \operatorname{Det}\!\left(\operatorname{Cone}
 F_{(\ell,N'),(k,N)}\right).
\end{equation}
This formula also records the refinements that add homology. The
Smith reading is conserved without feigning invariance: each transition is
assigned the total certificate
\begin{equation}
\label{eq:pdfv2-cor-transicion-smith}
 \mathcal S(F):=
 \bigl(\operatorname{Smith}(C_{\ell,N'}),
       \operatorname{Smith}(C_{k,N}),
       \operatorname{Smith}(\operatorname{Cone}F)\bigr),
\end{equation}
which literally determines which summands are conserved, born, or die.

The action of a coherent symmetry \(g\) on the common tree induces
permutation isometries \(\widetilde g_k\) on
\(\mathscr H_k\). The bijection of children and preservation of the weights of
\eqref{eq:continuo-comun-naturalidad-medida} prove, on the basis vectors,
\begin{equation}
\label{eq:pdfv2-cor-naturalidad-transporte-hojas}
 \widetilde g_{k+1}U_k=U_k\widetilde g_k,
 \qquad
 \widetilde g_k\sigma_k=\sigma_k\widetilde g_k,
 \qquad
 \widetilde g_kP_k=P_k\widetilde g_k,
 \qquad
 \widetilde g_kQ_k=Q_k\widetilde g_k.
\end{equation}

\begin{theorem}[Totality and naturality of the correlative constructor]
\label{thm:pdfv2-cor-totalidad-naturalidad}
For every total state of \eqref{eq:continuo-comun-estado}, every finite horizon,
and every coefficient level, the expression
\eqref{eq:pdfv2-cor-constructor-unico} is defined. Its five
characteristics are simultaneous operations on the same extension ledger and
satisfy the single identity \eqref{eq:pdfv2-cor-naturalidad-unica}. The collar, the complex, and the determinant
line remain nonempty even when their subsequent readers are neither acyclic nor
invertible.
\end{theorem}

\begin{proof}
Nonemptiness of the tree and its ports is
\cref{prop:continuo-comun-arbol-finito-no-vacio}; the measure and the expectation
are total by \cref{thm:continuo-comun-medida-canonica}. The formulas
  \eqref{eq:pdfv2-cor-u-comun}--\eqref{eq:pdfv2-cor-memorias} define transport,
polarity, balance, and memory for all generators. The four
  directions of \eqref{eq:pdfv2-cor-cuatro-rectas-tpk} are distinct by
  \eqref{eq:pdfv2-cor-proyectores-no-colapso}; the cellular completion
  \cref{def:pdfv2-cor-completacion-celular-tpk} contains its two path orders and
  six cells, while
  \eqref{eq:pdfv2-cor-collar-naturalidad} proves its naturality. The
cellular sequence is exact by \eqref{eq:pdfv2-cor-sucesion-celda}, the
local complex satisfies \(\partial^2=0\), and the correlated filling is explicit by
\eqref{eq:pdfv2-cor-filler}. Every finite based construction possesses the line
\eqref{eq:pdfv2-cor-linea-determinante-total}; homology, Bockstein, and the
Smith certificate are transported by
\eqref{eq:pdfv2-cor-naturalidad-homologia-bockstein}--
\eqref{eq:pdfv2-cor-transicion-smith}. Finally, the identities are
verified on each generator \(\mathfrak e_\alpha\) through
\eqref{eq:pdfv2-cor-naturalidad-componentes},
\eqref{eq:pdfv2-cor-collar-naturalidad}, and
\eqref{eq:pdfv2-cor-naturalidad-transporte-hojas}; by linearity and by the
functorial composition of routes,
\eqref{eq:pdfv2-cor-naturalidad-unica} is obtained for the integral record.
\end{proof}

\section{Causal separation of the identities, deductions, and realizations of the correlator constructor}

\paragraph{No vacuity and exact identities of the correlator constructor}
The nonemptiness imported from \eqref{eq:continuo-comun-out-finito} and the
identities
\eqref{eq:pdfv2-cor-b-identidades},
\eqref{eq:pdfv2-cor-sucesion-celda},
\eqref{eq:pdfv2-cor-sb-cero},
\eqref{eq:pdfv2-cor-u-gram},
\eqref{eq:pdfv2-cor-balance-hojas},
\eqref{eq:pdfv2-cor-jensen-hojas},
\eqref{eq:pdfv2-cor-dos-cero},
\eqref{eq:pdfv2-cor-cociclo-matricial}, and
\eqref{eq:pdfv2-cor-filler} follow from the definitions and the exhibited
APP--TRIT--TPK relations. Their status is
\texttt{EXACT\_INTERNAL}.

\paragraph{Starting hypotheses for the isometry of \(U_k\), the cellular sequence, and the composition \(\Psi/K\)}
The construction of \(U_k\) takes as its starting structure the finiteness of the fibres, survival at every horizon, functional truncation, and the transports of both sheets already produced by TPK\@. Equality \(U_k^*U_k=I\) is equivalent to verifying the normalized average Gram matrix of \eqref{eq:pdfv2-cor-gram-promedio-normalizado}; in particular, it suffices to check \(U_\alpha^*U_\alpha=I\) for each emission. General identity \eqref{eq:pdfv2-cor-u-gram} is retained for all declared sheet transports. The cellular sequence uses both ordered APP ports. Composition \(\Psi/K\) uses integer carry, orientation, and boundary. Under these enumerated premises, the stated conclusions are affirmative.

\paragraph{Additional data for isometric realizations, balanced retracts, and trivialization of the determinant line}
The internal constructor is completed by the positive Gram form of
\eqref{eq:pdfv2-cor-u-gram}, the total space \(Q_{\rm inc}\), the cellular sequence,
and the determinant line with its transition cone. A specialized
publication may require additional data: the isometry of each
\(U_\alpha\), a concrete numerical identification of the APP defect, a
balanced acyclic retraction, or trivialization of a determinant
line. Those data are intertwiners of the published codomain; they do not
form part of the domain or of the existence of
\(\mathfrak D_{k,N}\). Their absence precludes exclusively the promotion that
mentions them.

\begin{criterion}[Falsifiers of the intrinsic correlation]
\label{crit:pdfv2-cor-falsadores}
The joint construction fails at the first level at which any of the
following facts occurs:
\begin{enumerate}
 \item an operation uses an extension, orientation, carry, memory, boundary or route different from that recorded in \(\mathfrak e_\alpha\);
 \item a child of
 \(\operatorname{Ch}_k(\widehat x_k)\) is chosen and the remaining ones are eliminated;
 \item the weights of \eqref{eq:pdfv2-cor-peso-conteo} do not sum to one, \eqref{eq:pdfv2-cor-u-gram} fails, or \(U_k^*U_k=I\) is declared without verifying the normalized average Gram matrix of \eqref{eq:pdfv2-cor-gram-promedio-normalizado}, for example through the isometry of each emission transport;
 \item \(\kappa_k^{\rm can}\) is identified with the integer carry
 \(\kappa_\alpha\);
 \item \(sB=0\) fails, \(Q_{\rm inc}(\widehat x)\) is contracted
 to its origin, the action \(q\mapsto q+Ba\) is altered, or the two path orders are identified before
 applying the reader;
 \item a refinement breaks
 \eqref{eq:pdfv2-cor-refinamiento-cadena} or the cocycle
 \eqref{eq:pdfv2-cor-cociclo-matricial};
 \item fails \(\partial^2=0\), \(z_\pm\) are not ciclos or
 \(\partial h_*\neq z_--z_+\);
 \item a nonzero determinantal scalar is asserted without the retract, the
 balance, and the acyclicity required by that publication;
 \item any of the five operations is presented as an object
 separate from the constructor \(\mathfrak D_{k,N}\).
\end{enumerate}
In all those cases the diagnosis is: \emph{common object substituted;
conclusion not transferable}.
\end{criterion}
